% concatenated from main.tex
\documentclass[accepted]{article}
% % \documentclass{article}
% \documentclass[preprint]{article}

\newif\ifshowappendix
\showappendixtrue % change to \showappendixtrue and comment externaldocuments to produce the full text including appendix


% if you need to pass options to natbib, use, e.g.:
%     \PassOptionsToPackage{numbers, compress}{natbib}
% before loading neurips_2025

\usepackage{hyperref}
\hypersetup{
    colorlinks=true,
    linkcolor=blue,
    filecolor=blue,
    citecolor=blue,
    urlcolor=blue,
}

\usepackage{xr-hyper}
% \externaldocument{neurips_2025_Hrad_suppl} % No .tex extension
% \externaldocument{labels_appendix_Hrad}
%%%%% Numeric
 \PassOptionsToPackage{numbers,sort&compress}{natbib}
\bibliographystyle{unsrtnat}

%ready for submission

%%%%%%%% Author,Year
% \usepackage{natbib}
% \bibliographystyle{apalike}

\usepackage{neurips_2025}


% to compile a preprint version, e.g., for submission to arXiv, add add the
% [preprint] option:
%     \usepackage[preprint]{neurips_2025}


% to compile a camera-ready version, add the [final] option, e.g.:
%     \usepackage[final]{neurips_2025}


% to avoid loading the natbib package, add option nonatbib:
%    \usepackage[nonatbib]{neurips_2025}

\usepackage[table]{xcolor}         % colors
\definecolor{Magenta}{rgb}{1,0,1}
\usepackage[utf8]{inputenc} % allow utf-8 input
\usepackage[T1]{fontenc}    % use 8-bit T1 fonts
\usepackage{hyperref}       % hyperlinks
\usepackage{url}            % simple URL typesetting
\usepackage{booktabs}       % professional-quality tables
\usepackage{amsfonts}       % blackboard math symbols
\usepackage{nicefrac}       % compact symbols for 1/2, etc.
\usepackage{microtype}      % microtypography
\usepackage{xcolor} 
\usepackage{booktabs}
\usepackage{graphicx} % colors
\usepackage{amsmath,amsthm,amssymb, mathtools}
% \usepackage{wrapfig}
\theoremstyle{plain} % default
\newtheorem{theorem}{Theorem}[section] 
\newtheorem{lemma}[theorem]{Lemma}
\newtheorem{proposition}[theorem]{Proposition}
\newtheorem{assumption}[theorem]{Assumption}
\newtheorem{corollary}[theorem]{Corollary}

\theoremstyle{definition}
\newtheorem{definition}{Definition}[section]
\newtheorem{example}[theorem]{Example}

\theoremstyle{remark}
\newtheorem{remark}{Remark}[section]

\newcommand{\eps}{\varepsilon}
\newcommand{\bQ}{\mathbf{Q}}
\newcommand{\E}{\mathbb{E}}
\newcommand{\R}{\mathbb{R}}
\newcommand{\N}{\mathbb{N}}
\newcommand{\X}{\mathcal{X}}
\newcommand{\ddP}{\frac{dP}{dP_0}}
\newcommand{\TV}{\text{TV}}
\newcommand{\KL}{\mathrm{KL}}
\newcommand{\chisq}{\chi^2}

\newcommand{\SA}[1]{\textbf{\textcolor{Magenta}{[#1]}}}
\newcommand{\HG}[1]{\textcolor{cyan}{[#1]}}
\newcommand{\BW}[1]{\textbf{\textcolor{blue}{[#1]}}}

\newcommand{\ngb}{N_\gamma(P_0)}
\newcommand{\Pt}{\tilde{\mathcal{P}}}
\newcommand{\Pth}
{\tilde{\mathcal{P}}_{c_1, c_2, h}}
\newcommand{\Ptm}
{\tilde{\mathcal{P}}_{c_1, c_2, \mu}}
\title{Locally Optimal Private Sampling: Beyond the Global Minimax}


% The \author macro works with any number of authors. There are two commands
% used to separate the names and addresses of multiple authors: \And and \AND.
%
% Using \And between authors leaves it to LaTeX to determine where to break the
% lines. Using \AND forces a line break at that point. So, if LaTeX puts 3 of 4
% authors names on the first line, and the last on the second line, try using
% \AND instead of \And before the third author name.


% \author{%
%   David S.~Hippocampus\thanks{Use footnote for providing further information
%     about author (webpage, alternative address)---\emph{not} for acknowledging
%     funding agencies.} \\
%   Department of Computer Science\\
%   Cranberry-Lemon University\\
%   Pittsburgh, PA 15213 \\
%   \texttt{hippo@cs.cranberry-lemon.edu} \\
%   % examples of more authors
%   % \And
%   % Coauthor \\
%   % Affiliation \\
%   % Address \\
%   % \texttt{email} \\
%   % \AND
%   % Coauthor \\
%   % Affiliation \\
%   % Address \\
%   % \texttt{email} \\
%   % \And
%   % Coauthor \\
%   % Affiliation \\
%   % Address \\
%   % \texttt{email} \\
%   % \And
%   % Coauthor \\
%   % Affiliation \\
%   % Address \\
%   % \texttt{email} \\
% }

\author{%
  Hrad Ghoukasian \\
  Department of Computing and Software \\
  McMaster University \\
  \texttt{ghoukash@mcmaster.ca}
  \And
  Bonwoo Lee \\
  Department of Mathematical Sciences \\
  Korea Advanced Institute of Science \& Technology \\
  \texttt{righthim@kaist.ac.kr}
  \And
  Shahab Asoodeh \\
  Department of Computing and Software \\
  McMaster University \\
  \texttt{asoodeh@mcmaster.ca}
}



\begin{document}


\maketitle


\begin{abstract}
  % We consider the problem of sampling from a distribution under local differential privacy (LDP). 
  % Given a private distribution $P\in \mathcal P$, a locally private sampler must output a single observation from a distribution that is close in an $f$-divergence to $P$ while satisfying LDP. Sampling abstracts the goal of generating small amounts of realistic-looking private data.  In contrast to the previous work of Park et al (NeurIPS'24) that focuses on the global minimax locally private samplers, we study the local behavior of the minimax error around a given distribution $P_0$ and provide pointwise bounds that depend on $P_0$ and privacy parameters. In particular, we show that the local minimax error is fully determined by the global minimax error with $\mathcal P$ restricted to a neighborhood around $P_0$.  We prove this reduction, by  (1) extending previous work from pure LDP to the functional LDP framework and then (2) showing the global optimal functional LDP sampler is closely related to the optimal sampler when the input universe is restricted to a certain neighborhood around $P_0$. We also derive a simple closed-form expression for the local minimax samplers. 
  % We demonstrate by several experiments that the optimal local minimax samplers strictly outperform global minimax samplers.    
We study the problem of sampling from a distribution under local differential privacy (LDP). Given a private distribution $P \in \mathcal{P}$, the goal is to generate a single sample from a distribution that remains close to $P$ in $f$-divergence while satisfying the constraints of LDP. This task captures the fundamental challenge of producing realistic-looking data under strong privacy guarantees.
While prior work by Park et al. (NeurIPS'24) focuses on global minimax-optimality across a class of distributions, we take a local perspective. Specifically, we examine the minimax risk in a neighborhood around a fixed distribution~$P_0$, and characterize its exact value, which depends on both~$P_0$ and the privacy level.
Our main result shows that the local minimax risk is determined by the global minimax risk when the distribution class $\mathcal{P}$ is restricted to a neighborhood around $P_0$.
To establish this, we (1) extend previous work from pure LDP to the more general functional LDP framework, and (2) prove that the globally optimal functional LDP sampler yields the optimal local sampler when constrained to distributions near $P_0$. Building on this, we also derive a simple closed-form expression for the locally minimax-optimal samplers which does not depend on the choice of $f$-divergence.
We further argue that this local framework naturally models private sampling with public data, where the public data distribution is represented by $P_0$. In this setting, we empirically compare our locally optimal sampler to existing global methods, and demonstrate that it consistently outperforms global minimax samplers.

\end{abstract}
\vspace{0.05cm}
\section{Introduction}\label{sec: introduction}
\vspace{0.1cm}
Differential privacy (DP) \citep{dwork2006calibrating} is the de facto standard for providing formal privacy guarantees in machine learning. A widely used variant, local differential privacy (LDP) \citep{kasiviswanathan2011can}, enables individuals to randomize their data on their own devices before sharing it with an untrusted curator. LDP ensures that the output of a randomized algorithm is statistically indistinguishable under arbitrary changes to the input, thereby limiting the privacy leakage of an individual’s data. By eliminating the need for a trusted curator, LDP has become particularly appealing to users and companies alike, and has seen wide deployment in practice—including systems developed by Google, Apple, and Microsoft\citep{erlingsson2014rappor, bittau2017prochlo, Apple_Privacy, ding2017collecting}.

Despite this progress, much of the LDP literature assumes that each user holds only a single data point. This assumption is often unrealistic, as modern devices routinely collect large volumes of data—such as images, messages, or time-series records~\citep{husain2020local}. To address this, a line of work known as \emph{user-level DP} assumes that each user holds a dataset of fixed size, with samples drawn i.i.d.\ from an underlying distribution~\citep{kent2024rate, mao2024privshape, bassily2023user, acharya2023discrete, huang2023federated, girgis2022distributed, levy2021learning, ghazi2021user, zhao2024huber}. However, the requirement of equal dataset size across users limits its practicality.

To overcome this limitation, \textit{locally private sampling} was introduced by \citet{husain2020local} and later extended by \citet{park2024exactly}. This framework considers a more general setting where clients possess large datasets of varying sizes and aim to privately release another dataset that closely resembles their original data. Each client's local dataset is modeled as an empirical distribution $P$, from which a single private sample is to be generated. Privacy is enforced by projecting $P$ onto a set of distributions— called a \textit{mollifier}—within a certain radius. While \citet{husain2020local} construct the mollifier in an ad-hoc manner using a reference distribution, \citet{park2024exactly} develop a principled minimax framework to identify the \textit{optimal} mollifier, which corresponds to the worst-case data distribution (namely, degenerate distributions in the discrete case).

In this work, we extend the framework of \citet{park2024exactly} to a public-private setting, where each individual also has access to a public dataset, represented by a distribution $P_0$. From a theoretical perspective, our goal is to move beyond the pessimistic worst-case formulation and instead characterize optimal samplers through a local minimax lens. From a practical standpoint, this framework is motivated by the increasing relevance of personalized collaborative learning~\citep{ozkara2023a,Beirami_personalization,Personalization3}, in which individual models are trained collaboratively, and the public data may reflect (privatized) information shared by other users. Since real-world data distributions are rarely worst-case, our local formulation better captures achievable performance in practice.

Our local minimax framework aims to identify locally private samplers that minimize the $f$-divergence between the true distribution $P$ and the sampling distribution, for all $P$ within a neighborhood $N_\gamma(P_0)$ around $P_0$, parameterized by a constant $\gamma \geq  1$. This neighborhood is defined via one of the strictest notions of distance between distributions: we say $Q \in N_\gamma(P_0)$ if $Q(x) \leq \gamma P_0(x)$ and $P_0(x) \leq \gamma Q(x)$ for all $x$. We will discuss this notion of neighborhood in details in Section \ref{sec: local functional}. 

Interestingly, we show that the local minimax-optimal private sampler under this notion of neighborhood is fully characterized by the global minimax-optimal sampler for a restricted class of data distributions. This relationship is most transparent in the context of functional LDP~\citep{dong2022gaussian}. We therefore first extend the global minimax framework of \citet{park2024exactly} from pure LDP to functional LDP, and then use it to derive the local minimax-optimal samplers. 



\textbf{Contributions.} Our main contributions are as follows: 
\begin{itemize}
\item We characterize the exact global minimax risk of private sampling under the functional LDP framework and derive optimal samplers for both continuous and discrete domains. Leveraging properties of functional LDP~\citep{dong2022gaussian}, our results generalize the pure-LDP framework of~\citet{park2024exactly} to approximate and Gaussian LDP (GLDP) notions  (Section~\ref{sec: global functional}). As in their case, our optimal samplers are independent of the choice of $f$-divergence.

\item  We introduce a local minimax formulation to identify samplers that minimize $f$-divergence over all $P \in N_\gamma(P_0)$ for a given $P_0$ and $\gamma \geq 1$. We then show that this minimax risk is fully determined by the global minimax risk when the distribution class is restricted to a neighborhood around \( P_0 \) (Section~\ref{sec: local functional}). Building on our global minimax results, we derive closed-form expressions for the local minimax-optimal functional samplers. We further specialize this to express a pointwise optimal private sampler under pure LDP that achieves the local minimax risk (Section~\ref{sec: local pure}), enabling direct comparison with the global samplers of~\citet{park2024exactly}.

\item We numerically validate our local minimax samplers, showing they consistently—and often substantially—outperform the global samplers of~\citet{park2024exactly} across privacy regimes (Section~\ref{sec: numeric}). Figure~\ref{fig: Laplace pure and functional local vs global} illustrates that our samplers yield distributions significantly closer to the original under both pure LDP and GLDP.\footnote{Experiment code is publicly available at 
\url{https://github.com/hradghoukasian/private_sampling}, 
with reproduction instructions provided in Appendix~\ref{appendix: reproduce}.}


\end{itemize}




\textbf{Related work.}
Locally private sampling was initially introduced by~\citet{husain2020local} and later formalized through a global minimax framework by~\citet{park2024exactly}. In their work, two fundamentally different families of minimax-optimal samplers were identified: one derived from the randomized response mechanism, referred to as the \emph{linear sampler}, and another based on clipping the data distribution, known as the \emph{non-linear sampler}. They rigorously showed that the non-linear sampler is pointwise better than the linear one. In a similar vein, we develop two families of locally minimax-optimal pure LDP samplers and demonstrate that one is pointwise better than the other.
% \HG{Is this sentence saying  that because in \emph{local} minimax we have projection into the neighborhood and then sampling, it is not trivial? Because our functional LDP local minimax sampler is indeed randomized response inside the neighborhood.} \SA{Yes some of our results are linear, yet optimal! At least for the pure LDP, we have the pointwise optimal one though. I feel like we need to explain a bit technically why instantiation of Thm 4.1 on $g_\eps$ is strictly weaker than Thm 5.1. We can say the former and latter differs in only when $P\notin N_\gamma$ and in that case the former is a linear sampler and the latter non-linear. Thus intuitively we have pointiwse optimality. This discription makes Fig 2 more understandable as well.}\HG{I see your point! Should sth be added here?} \SA{Done!}

Sampling with public data was recently explored by~\citet{zamanlooy2024locally}, who study minimax-optimal samplers in the presence of a public prior. While related, their work differs in key ways: they focus on a global minimax formulation restricted to discrete domains and linear samplers (i.e., perturbing the data distribution linearly), whereas we study a local minimax problem, support both discrete and continuous settings, and consider arbitrary samplers. Moreover, the role of public data differs between the two approaches. \citet{zamanlooy2024locally} impose it as a hard constraint: their samplers are required to preserve the public prior exactly, ensuring that the public data remains unchanged while privacy is enforced only on the private data. By contrast, we use the public distribution to define a local neighborhood, which in turn yields a local minimax formulation applicable to general sample spaces. A key implication of this distinction is that in their setting, access to public data does not necessarily reduce the minimax risk, whereas in ours it does.

Private sampling under central DP was introduced by~\citet{NEURIPS2021_f2b5e92f}, who proposed algorithms for generating private samples from $k$-ary and product Bernoulli distributions. They showed that sampling can incur significantly lower privacy costs than private learning for some range of parameters. This was extended to broader distribution families~\citep{NEURIPS2023_f4eaa4b8,kamath_sampling,tasnim2025reveal} and to multi-sample settings~\citep{MultiSampling}.

Private sampling has also gained traction in fine-tuning large language models~\citep{flemings2024differentially,flemings2024adaptively}, and is closely related to private generative modeling. However, most prior work in this area focuses on central DP~\citep{raskhodnikova2021differentially,ghazi2023differentially,ebadi2016sampling,abay2018privacy,xie2018differentially,xin2020private,torkzadehmahani2019dp,liu2019ppgan}, assumes single-record LDP~\citep{cunningham2022geopointgan,shibata2023local,zhang2023publishing,gwon2024ldp}, or targets specific estimation tasks~\citep{imola2024metric}.

Finally, we remark that local minimax formulations have also been explored in the context of estimation and distribution learning under DP \cite{chen2024lq, mcmillan2022instance, Duchi_Complexity,asi2020instance,feldman2024instance,rohde2020geometrizing}. In particular, \citet{chen2024lq} analyze discrete distribution estimation under LDP by considering neighborhoods around a reference distribution. Similarly, \citet{mcmillan2022instance} and \citet{Duchi_Complexity} develop locally minimax estimators for one-dimensional parameters in the central and local models, respectively, while \citet{asi2020instance} extend this framework to function estimation. More recently, \citet{feldman2024instance} utilize the notion of neighborhoods around a distribution in the context of nonparametric density estimation.
Although these works focus primarily on estimating statistical functionals or learning distributions, they are conceptually distinct from our setting, which centers on private sampling. Nonetheless, our notion of neighborhood in the local minimax formulation is closely aligned with that of~\citet{feldman2024instance}. 

 
%In this work, we focus on private sampling, where the goal is to generate a single data point from a user's distribution under LDP constraints.

% \citet{husain2020local} formalized the private sampling problem and proposed a locally differentially private algorithm that generates a sample approximating a user's input distribution. However, the performance of their method is highly sensitive to the choice of an arbitrary reference distribution, making it less reliable in practice.
% \citet{park2024exactly}, however, developed a minimax formulation for locally private sampling and   characterized the minimax risk of the private sampling problem for general $f$-divergences and identifying an optimal sampling mechanism that achieves this risk. Their formulation assumes that input distributions belong to a global class, offering strong worst-case guarantees. However, such global formulations are often overly conservative, as they focus on the hardest possible instances. A more adaptive alternative is the local minimax framework, which tailors the analysis to a specific instance and often yields more practical guarantees. Local minimax-optimality has been widely studied in both private and non-private settings, leading to optimal estimators across a wide range of estimation tasks~\citep{WeiNing_Pointwise,chen2024lq,cai2015framework,donoho1994minimax,feldman2024instance,Duchi_Complexity,asi2020instance,asi2020near,chen2025decision,huang2021instance,mcmillan2022instance}.









\setlength{\textfloatsep}{0.4cm}
\begin{figure}[t]
  \centering
  \includegraphics[width=0.72\linewidth]{figures/LaplaceMixture_2x3Comparison.pdf}
  \caption{Comparison of global and local minimax-optimal sampler under pure LDP ($\varepsilon = 1$) and $\nu$-GLDP ($\nu = 1.5$). Full details of this experiment are provided in Appendix~\ref{appendix: laplace mixture}. 
}
  \label{fig: Laplace pure and functional local vs global}
\end{figure}
























% % Commented Fig
% \begin{figure}[ht]
%   \centering
%   \includegraphics[width=0.8\linewidth]{figures/LaplaceMixture_2x3Comparison.pdf}
%   \caption{Comparison of global and local minimax-optimal mechanisms under pure LDP ($\varepsilon = 1$) and $\nu$-GDP ($\nu = 1.5$).}
%   \label{fig: Laplace pure and functional local vs global}
% \end{figure}

% \begin{figure}[ht]
%   \centering
%   \includegraphics[width=0.8\linewidth]{figures/Lap_GlobalvsLocal_Pure_eps=1_c=3.png}
%   \caption{Comparison of global and local minimax-optimal mechanisms under pure LDP for $\eps = 1$.}
%   \label{fig: Laplace pure local vs global eps = 1}
% \end{figure}

% \begin{figure}[ht]
%   \centering
%   \includegraphics[width=0.8\linewidth]{figures/Lap_GlobalvsLocal_GDP_mu=1.5_c=3.png}
%   \caption{Comparison of global and local minimax-optimal mechanisms under functional LDP with $\nu$-GDP ($\nu= 1.5$).}
%   \label{fig: Laplace GDP local vs global nu = 1.5}
% \end{figure}

% \begin{figure}[ht]
%   \centering
%   \includegraphics[width=0.8\linewidth]{figures/Lap_GlobalvsLocal_GDP_mu=1_c=3.png}
%   \caption{Comparison of global and local minimax-optimal mechanisms under functional LDP with $\nu$-GDP ($\nu= 1$).}
%   \label{fig: Laplace GDP local vs global nu = 1}
% \end{figure}


% \textcolor{blue}{[I think most part of the last two paragraphs of intro can be removed to related work. In the intro, I think it is enough saying 1. global is bad. 2. local is good, shown in prior works (cite only neighborhood-based). 3. So we develop local sampler and it was shown to be good in these aspects A,B,... You can elaborate on the existing approaches in private local estimator, and why you choose neighbor-based approach in detail in related work section.]}


 %    \textcolor{blue}{[Do we need to define decomposability if we will state results for continuous case and discrete case each?]\HG{The original plan was to present the continuous case first, then the general case, and finally the discrete case—to highlight the generality of our framework. We also planned to note that the discrete case follows from the general one, so we only focus on the continuous case in later results. I agree that introducing the decomposability assumption could add complexity, though it does lead to a stronger statement, which I like. Let’s check with Shahab.}
 % }\textcolor{blue}{[ So you are going from specific to general? Then I think it is good enough. I got confused by introducing continuous setting in the start of section3->general->discrete case.]}\HG{Exactly, I was going from specific (continuous) to general then specific (discrete). @Shahab, do you suggest keeping the order as it is or do you think adding decomposability assumption for general case can be confusing?} 
 

%\subsection{Related work}\label{subsec: related work}
\textbf{Notation.} Let $\mathcal{P}(\mathcal{X})$ denote the set of all probability distributions over the sample space $\mathcal{X}$, and let $\mathcal{C}(\mathbb{R}^n)$ denote the set of all continuous probability distributions on $\mathbb{R}^n$. For $P\in \mathcal{C}(\mathbb{R}^n)$, we denote its probability density function (PDF) by the lowercase letter $p$. For any integer $k \in \mathbb{N}$, we define $[k] \coloneqq \{1, 2, \dots, k\}$. The notation $\mathrm{clip}(x; s_1, s_2) \coloneqq \max\{s_1, \min\{s_2, x\}\}$ refers to the clipping function that bounds $x$ between $s_1$ and $s_2$.
% We use $\mathcal{Q}_{\mathcal{X}, \tilde{\mathcal{P}}, g}$, $\mathcal{Q}_{\mathcal{X}, \tilde{\mathcal{P}}, \varepsilon}$, and $\mathcal{Q}_{\mathcal{X}, \tilde{\mathcal{P}}, \varepsilon, \delta}$ to denote the sets of all mechanisms $\bQ : \tilde{\mathcal{P}} \to \mathcal{P}(\mathcal{X})$ that satisfy, respectively, $g$-FLDP, $\varepsilon$-LDP, and $( \varepsilon, \delta )$-LDP. 
 % The convex conjugate of a function \( g \) is denoted by \( g^* \). 
 The convex conjugate of a function $g:D\subseteq\mathbb{R}\to\mathbb{R}$ is the function $g^*:\mathbb{R}\to(-\infty,+\infty]$ defined by $g^*(y) \coloneqq \sup_{x\in D} \{\, xy - g(x) \,\}$ for $y\in\mathbb{R}$.
We denote the \( n \)-dimensional Laplace distribution with mean \( m \in \mathbb{R}^n \) and scale \( b > 0 \) by \( \mathcal{L}(m, b) \), and its density at \( x \in \mathbb{R}^n \) by \( \mathcal{L}(x \mid m, b) \). For the Gaussian case, \( \mathcal{N}(m, \Sigma) \) denotes the \( n \)-dimensional normal distribution with mean \( m \in \R^n\) and covariance matrix \( \Sigma \).

\vspace{0.1cm}
\section{Preliminaries}\label{sec: prilim}
\vspace{0.1cm}
In this section, we formally introduce definitions and concepts necessary for the subsequent sections.


Suppose a client holds a distribution \( P \in \tilde{\mathcal{P}} \subset \mathcal{P}(\mathcal{X}) \) over a sample space \( \mathcal{X} \), and wishes to release a sample that resembles being drawn from \( P \) while preserving privacy. To ensure privacy, the client applies a private sampler \( \mathbf{Q}: \tilde{\mathcal{P}} \to \mathcal{P}(\mathcal{X}) \), which maps the input distribution \( P \) to a privatized distribution \( \mathbf{Q}(P) \in \mathcal{P}(\mathcal{X}) \). The released sample is then drawn from \( \mathbf{Q}(P) \), thereby preserving privacy while approximating the original distribution. We equivalently view \( \mathbf{Q}(P) \) as the conditional distribution \( \mathbf{Q}(\cdot \mid P) \).


% To ensure privacy, the client applies a private sampler $\mathbf{Q}: \tilde{\mathcal{P}} \to \mathcal{P}(\mathcal{X})$, the output of the sampler is a distribution in $\mathcal{P}(\mathcal{X})$ and then to draw a sample from  and the released sample is drawn from the sampling distribution. 

 % defined as a conditional distribution from \( \tilde{\mathcal{P}} \) to \( \mathcal{X} \), and the released sample is drawn from the sampling distribution \( \mathbf{Q}(\cdot \mid P) \). 
 %The sampler $\bQ$ is required to satisfy different variants of LDP—such as pure, approximate, or functional—depending on the application \citep{husain2020local, duchi2013local, park2024exactly}.
 % \SA{above, you said sampler is from $\tilde{\mathcal{P}}$ to $\mathcal X$? but below is from $\tilde{\mathcal{P}}$ to $\mathcal{P}(\mathcal X)$ Pick the second one and keep it that way. }\HG{I see your point. what I had was confusing. I needed the conditional distributional interpretation of $\bQ(.|P)$ for the privacy definitions to be meaningful though. I deleted the domain and do-domain of the conditional distribution. Is it clear now?}  \SA{It is great now!}
\begin{definition}[Approximate LDP]
\label{def:approx-LDP}
Let $\varepsilon \geq 0$ and $\delta \in [0,1]$.  
A sampler $\mathbf{Q}: \tilde{\mathcal{P}} \to \mathcal{P}(\mathcal{X})$ is said to satisfy \emph{$(\varepsilon,\delta)$-local differential privacy ($( \varepsilon,\delta)$-LDP)} if, for every pair of input distributions $P, P' \in \tilde{\mathcal{P}}$ and every measurable set $A \subseteq \mathcal{X}$, we have
\[
\mathbf{Q}(A \mid P) \leq e^{\varepsilon} \, \mathbf{Q}(A \mid P') + \delta.
\]
When $\delta = 0$, $\mathbf{Q}$ is $\eps$-LDP sampler \cite{husain2020local}, also referred to as \emph{pure LDP}. We let $\mathcal{Q}_{\mathcal{X}, \tilde{\mathcal{P}}, \varepsilon}$ and $\mathcal{Q}_{\mathcal{X}, \tilde{\mathcal{P}}, \varepsilon, \delta}$ denote the sets of all $\varepsilon$-LDP and $(\varepsilon, \delta)$-LDP samplers, respectively.

\end{definition}


% Definition \ref{def:approx-LDP} ensures privacy by requiring that the output of the sampling mechanism remains indistinguishable under any change in the input. In contrast, functional LDP interprets privacy within the framework of hypothesis testing. Specifically, let \(P\) and \(Q\) denote the output distributions of a private mechanism corresponding to two different inputs. Functional LDP measures privacy guarantee of the mechanism by quantifying the difficulty an adversary faces in distinguishing between these two distributions. This can be equivalently framed as the following binary hypothesis testing problem \citep{wasserman2010statistical,kairouz2015composition}:


Definition~\ref{def:approx-LDP} ensures privacy by requiring that the output of the sampler remains indistinguishable under any arbitrary changes in the input. In contrast, functional LDP \cite{dong2022gaussian} interprets privacy through the lens of hypothesis testing: given two inputs, let \(P\) and \(Q\) be the corresponding output distributions. The privacy guarantee is quantified by how hard it is for an adversary to distinguish \(P\) from \(Q\), which can be framed as a binary hypothesis testing problem \citep{wasserman2010statistical,kairouz2015composition}:
\[
H_0:\ \text{output} \sim P 
\quad\text{vs}\quad 
H_1:\ \text{output} \sim Q.
\]
Given a rejection rule \(\phi:\mathcal{X}\rightarrow[0,1]\), the Type I and Type II error rates  are defined as
\[
    a_\phi := \mathbb{E}_{P}[\phi],
    \qquad
    b_\phi := 1 - \mathbb{E}_{Q}[\phi],
\]
respectively.
%Based on this formulation, we define the trade-off function and the notion of $g$-FLDP.

\begin{definition}[Trade--off function]\label{def:trade-off}
Let \(P\) and \(Q\) be probability measures in $\mathcal{P}(\X)$. 
The \emph{trade--off function}
\(T(P,Q):[0,1]\to[0,1]\) is defined as
\[
   T(P,Q)(u)
   \;=\;
   \inf_{\phi}\,\bigl\{\,b_\phi : a_\phi \le u \bigr\},
\]
where the infimum is taken over all measurable rejection rules~\(\phi\).
\end{definition}

Notice that \(T(P,Q)(u)\) represents the smallest achievable Type II error
subject to Type I error being at most \(u\). A function $g$ is called a trade-off function if $g = T(P,Q)$ for some distributions $P$ and $Q$.

\begin{definition}[Functional LDP \citep{dong2022gaussian,lee2023minimax}]\label{def:gFLDP}
% Let \(g:[0,1]\to[0,1]\) be a convex, continuous, non-increasing function such that 
% \(g(t)\le 1-t\) for all \(t\in[0,1]\).
Given a trade-off function $g:[0,1]\to[0,1]$,
a sampler 
\(\mathbf{Q}~:~\tilde{\mathcal{P}}~\to~\mathcal{P}(\mathcal{X})\)
is said to satisfy \emph{\(g\)-functional local differential privacy} (\(g\)-FLDP) if for every pair of input distributions \(P, P' \in \tilde{\mathcal{P}}\) and every
\(u \in [0,1]\),
\[
   T\bigl(\mathbf{Q}(\cdot\mid P),\,
            \mathbf{Q}(\cdot\mid P')\bigr)(u)
   \;\ge\; g(u).
\]
We let $\mathcal{Q}_{\mathcal{X}, \tilde{\mathcal{P}}, g}$ denote the sets of all samplers $\mathbf{Q} : \tilde{\mathcal{P}} \to \mathcal{P}(\mathcal{X})$ satisfying $g$-FLDP.
\end{definition}
Note that functional LDP generalizes both approximate and pure LDP. Specifically, for any privacy parameters \( (\varepsilon, \delta) \), there exists a trade-off function \( g_{\varepsilon,\delta} \) such that \( g_{\varepsilon,\delta} \)-FLDP is equivalent to \( (\varepsilon,\delta) \)-LDP \citep[Proposition 3]{dong2022gaussian}. This function is given by $g_{\varepsilon, \delta}(\theta) = \max \left\{ 0, 1 - \delta - e^{\varepsilon} \theta,\ e^{-\varepsilon}(1 - \delta - \theta) \right\}$. The case of pure LDP corresponds to the special case $g_{\eps} \coloneqq g_{\eps,0}$.
%The case of pure LDP corresponds to the special case \( g_\varepsilon := g_{\varepsilon, 0} \). 
Another notable instance of \( g \)-FLDP is \( \nu \)-GLDP, where the trade-off function is given by
$G_\nu(\theta) = \Phi\left(\Phi^{-1}(1 - \theta) - \nu\right)$,
with \( \Phi \) denoting the standard normal CDF \citep{dong2022gaussian,lee2023minimax}.


% A function
% $g$ that satisfies the conditions in Definition~\ref{def:gFLDP} is a trade-off function. 
% \textcolor{blue}{[If such characterization of trade-off function is not used in the main body, it would be better to say $g$ is a trade-off function, and add in Definition 2.2 sentence like \textit{we say $g$ is a trade-off function if $g=TV(P,Q)$ for some distributions $P$ and $Q$.}]\HG{Yeah, the characterization of of $g$ in terms of its properties is only used in the appendix in som part of proof. It's updated now. }}


In order to quantify utility loss incurred by sampler $\bQ$, we use the \( f \)-divergence between the original distribution \( P \) and the resulting sampling distribution \( \bQ(P) \), defined as follows.
\begin{definition}[$f$-divergence \citep{ali1966general,csiszar1967information}]
\label{def:f-divergence}
Let \( f: (0,\infty) \to \mathbb{R} \) be a convex function with \( f(1) = 0 \). 
The $f$-divergence between $P, Q \in \mathcal{P}(\mathcal{X})$ with $P \ll Q$ is defined as
% \[D_f(P \| Q) = \mathbb{E}_Q\big[f(dP/dQ)\big].\]
$
D_f(P \| Q) := \mathbb{E}_Q\Big[ f\left( \frac{dP}{dQ} \right) \Big]
$.


% If $P \not\ll Q$, we set $D_f(P \| Q) := +\infty$.
\end{definition}
The above definition can be extended to cases where 
$\displaystyle P\ll Q$ is no longer satisfied as in Appendix~\ref{subsec: f-div}.
 Common examples of \( f \)-divergences include KL divergence, total variation distance, and squared Hellinger distance. We denote total variation distance between $P,Q\in\mathcal{P}(\mathcal{X})$ by \( D_{\mathrm{TV}}(P, Q) \). A key instance in this work is the \( E_\gamma \)-divergence (or hockey-stick divergence), defined as the $f$-divergence with \( f(t) = \max\{t - \gamma, 0\} \) for \( \gamma \geq 1 \)~\citep{sharma2013fundamental}. This divergence is particularly relevant as it defines the local neighborhood in our minimax formulation. We denote $E_\gamma$-divergence of two distributions $P$ and $Q$ as $E_\gamma(P \,\|\, Q)$.



% \SA{This definition can be removed if we need space!}
% \begin{definition}[$E_\gamma$-divergence \citep{sharma2013fundamental}]
% \label{def:Egamma-divergence}
% Given $\gamma \geq 1$, for two probability measures $P,Q \in \mathcal{P}(\mathcal{X})$ with $P \ll Q$, the
% \emph{$E_\gamma$-divergence} is
% \[
%     E_\gamma(P \,\|\, Q)
%     \;:=\;
%     \mathbb{E}_Q\!\left[
%         \max\!\Bigl\{\tfrac{dP}{dQ}-\gamma,\,0\Bigr\}
%     \right]= \frac{1}{2} \int \left| \mathrm{d}P - \gamma\,\mathrm{d}Q \right| - \frac{1}{2} (\gamma - 1).
% \]
% % \[
% % E_\gamma(P \| Q) := \frac{1}{2} \int \left| \mathrm{d}P - \gamma\,\mathrm{d}Q \right| - \frac{1}{2} (\gamma - 1).
% % \]
% \end{definition}


% \textcolor{blue}{[I like the part where you explain the meaning of the condition of hockey stick gamma being zero. But I am not sure how the connection between the divergence and the privacy is relevant to the paper.]}\HG{You're right. The connection with privacy does not play a significant role right now, I asked Shahab about this. Also, we were originally planning to mention in the future work section that a possible generalization would be to relax the neighborhood to $E_\gamma(P \,\|\, P_0) \leq \delta$, and potentially include a complementary result in the appendix. While we may not have time to derive a new result now, I think it's still worth keeping as a future direction. I defined the neighborhood using $E_\gamma$ mainly to make this connection more natural. What do you think?
% } \textcolor{blue}{[I also agree. By introducing $E_\gamma$ we can smoothly extend the class of neighborhood, and it includes the total variation-based neighborhood as well.]}
Our most general result, covering both continuous and discrete spaces, requires a technical assumption on distributions, stated below and discussed further in Section~\ref{sec: global functional}.


\begin{definition}[Decomposability \citep{park2024exactly}]\label{def: decomposability}
    Let $\alpha \in (0,1)$ and $t,u \in \mathbb{N}$ with $t > u$. A probability measure $\mu \in \mathcal{P}(\mathcal{X})$ is \emph{$(\alpha, t, u)$-decomposable} if there exist sets $A_1, \dots, A_t \subseteq \mathcal{X}$ such that $\mu(A_i) = \alpha$ for all $i \in [t]$, and for every $x \in \mathcal{X}$, the number of sets $A_i$ containing $x$ is at most $u$.
%     \textcolor{blue}{[Do we need to define decomposability if we will state results for continuous case and discrete case each? If we are going to mention this in main body, we should mention that each cases satisfy the decomposability condition.]\HG{The original plan was to present the continuous case first, then the general case, and finally the discrete case—to highlight the generality of our framework. We also planned to note that the discrete case follows from the general one, so we only focus on the continuous case in later results. I agree that introducing the decomposability assumption could add complexity, though it does lead to a stronger statement, which I like. Let’s check with Shahab.}
% \HG{Regarding your second point,do you mean I should explain that the decomposability assumption is satisfied for continuous and discrete right after introducing it? After Theorem~\ref{thm: global functional}, I note that it holds in the continuous case and in the discrete case when 
% $P_0$ is uniform. Should I emphasize this elsewhere in the main body?} }
% \textcolor{blue}{[I see. I must have missed it. Sorry. So you are going from specific to general? Then I think it is good enough. I got confused by introducing continuous setting in the start of section3->general->discrete case.]}\HG{Exactly, I was going from specific (continuous) to generel then specific (discrete). }
\end{definition}
Establishing the exact optimal minimax risk in our proofs involves deriving matching upper bounds (the achievability part) and lower bounds (the converse part). For a unified proof over general sample space $\mathcal{X}$, the converse part relies on $(\alpha, t, 1)$-decomposability; namely,  \( \mathcal{X} \) contains \( t \) disjoint measurable subsets, each of measure \( \alpha \).  We note that this mild condition holds naturally for absolutely continuous distributions and for the uniform reference measure in the finite case.



% \textcolor{blue}{.But if we are not going to use the full definition of the decomposability, we could just add it in Assumption 3.2 with other assumptions on the measure space, then explain. Plus, I think the specific examples given below Theorem 3.5 ($\mathcal{X}=[k]$ etc.) is good. Also, the paper use different notions for the same value $t=\frac{1}{\alpha}=\frac{c_2-c_1}{1-c_1}$. We could also define $t:=\frac{c_2-c_1}{1-c_1}\in\mathbb{N}$ and unify the notion.]}
% \HG{Regarding adding decomposability assumption in Assumption 3.2, I am using Assumption 3.2 for all continuous results but I only need decomposability for Theorem \ref{thm: global functional}} \textcolor{blue}{[I see. Then it will be inappropriate.]}





\vspace{0.1cm}
\section{Global minimax-optimal samplers under functional LDP}\label{sec: global functional}
\vspace{0.1cm}

% As outlined in the introduction, our goal is to characterize the locally minimax-optimal sampler for the private sampling problem. It can be shown that the technical framework developed by \citet{park2024exactly} cannot be naturally extended to approximate LDP (due to the notion of decomposability). In order to make a unified proof technique for different notions of LDP, we first study the global minimax formulation under functional LDP—a generalization of the pure LDP setting in \citet{park2024exactly} that also encompasses approximate LDP and GDP. This section establishes the global minimax mechanism and its optimal value, which will serve as the foundation for our local analysis in the next section. More precisely, the global minimax private sampling problem is formulated as:

In this section, we study the global minimax samplers under the general framework of functional LDP and derive compact characterizations of the optimal samplers. These global samplers will serve as proxies for constructing local minimax-optimal samplers in subsequent sections.

Using $f$-divergence as the utility measure, we define the global minimax risk as
\begin{equation}\label{eqn: global minimax functional problem}
\mathcal{R}\big(\mathcal{Q}_{\mathcal{X},\tilde{\mathcal{P}},g},\tilde{\mathcal{P}},f\big) \coloneqq \inf_{\bQ \in \mathcal{Q}_{\mathcal{X},\tilde{\mathcal{P}},g}} \hspace{0.1cm} \sup_{P \in \tilde{\mathcal{P}}}\hspace{0.1cm} D_f\big(P \, \| \, \bQ(P)\big).
\end{equation}
As noted by~\citet{park2024exactly}, this risk becomes vacuous in the continuous setting if the distribution class $\tilde{\mathcal{P}}$ is unrestricted. To address this, they consider a restricted class of distributions defined as
\begin{equation}\label{eqn: set of distributions}
   \tilde{\mathcal{P}} = \tilde{\mathcal{P}}_{c_1, c_2, h} := \left\{ P \in \mathcal{C}(\mathbb{R}^n) : c_1 h(x) \leq p(x) \leq c_2 h(x), \quad \forall x \in \mathbb{R}^n \right\}, 
\end{equation}

for a reference function $h : \mathcal{X} \to [0, \infty)$ satisfying $\int_{\mathbb{R}^n} h(x)\,dx < \infty$, and constants $c_2 > c_1 \geq 0$. 
%This assumption is followed by prior works on sampling and generative models literature~\citep{husain2020local, goodfellow2020generative, park2024exactly}, that assume the set of all data distributions $\tilde{\mathcal{P}}$ satisfies $\tilde{\mathcal{P}} \subseteq \tilde{\mathcal{P}}_{c_1, c_2, h}$.
This class captures a broad range of practical distributions, including mixtures of Gaussians and Laplace distributions. See the example below and further discussion in Appendices~\ref{appendix: laplace mixture} and~\ref{appendix: gaussian ring}.
%\footnote{Proof of all subsequent results is deferred to the appendix.} 

\begin{example}\label{example: Gaussian Laplace universe}
    For $n,k \in \mathbb{N}$,
define $\tilde{\mathcal{P}}_{\mathcal{L}}  = \Big\{ \sum_{i=1}^k \lambda_i \mathcal{L}( m_i, b) :  \lambda_i \geq 0, \, \sum_{i=1}^k \lambda_i = 1, \, \|m_i\|_1 \leq 1 \Big\}$ and $\tilde{\mathcal{P}}_{\mathcal{N}}  = \Big\{ \sum_{i=1}^k \lambda_i \mathcal{N}(m_i, \sigma^2 I_n)  :  \lambda_i \geq 0, \, \sum_{i=1}^k \lambda_i = 1, \, \|m_i\|_2 \leq 1 \Big\}$. Here $m_i, x \in \mathbb{R}^n$ and $I_n$ denotes the $n \times n$ identity matrix.  It can be verified that $\tilde{\mathcal{P}}_{\mathcal{N}} \subseteq \tilde{\mathcal{P}}_{0,\, 1,\, h_{\mathcal{N}}}$ and  $\tilde{\mathcal{P}}_{\mathcal{L}} \subseteq \tilde{\mathcal{P}}_{e^{-1/b},\, e^{1/b},\, h_{\mathcal{L}}}$, where $h_{\mathcal{N}}(x) = \frac{1}{(2\pi \sigma^2)^\frac{n}{2}} \exp\big(-\frac{[\max(0, \|x\|_2 - 1)]^2}{2\sigma^2} \big)$ and  $h_{\mathcal{L}}(x) = \mathcal{L}(x \mid \mathbf{0}, b)$. 
\end{example}
Without loss of generality, we assume that $\int_{\mathbb{R}^n} h(x),dx = 1$ and $0 \leq c_1 < 1 < c_2$. Moreover, for technical reasons---explained in Appendix~\ref{proof: theorem global functional}--- we also assume that \( \frac{c_2 - c_1}{1 - c_1} \in \mathbb{N} \). This assumption is not restrictive and holds in many practical settings, e.g., Gaussian mixtures (see Example~\ref{example: Gaussian Laplace universe}). When it is not satisfied, one can slightly increase \( c_2 \) or decrease \( c_1 \) to enforce the condition, resulting in a superset of the original distribution class, i.e., \( \tilde{\mathcal{P}} \subseteq \tilde{\mathcal{P}}_{c_1, c_2, h} \). 

While many of our examples focus on continuous domains, our analysis applies more generally to \emph{any} sample space~\( \mathcal{X} \), provided the distributions are absolutely continuous with respect to a reference measure~\( \mu \). In this case, we define  $
\tilde{\mathcal{P}}_{c_1, c_2, \mu}
:= \{P\in \mathcal{P}(\mathcal{X}):P \ll \mu, ~c_1 \leq \frac{dP}{d\mu} \leq c_2, ~~ \mu\text{-a.e.}\}$, which we take as the generalized distribution class.
Thus, we make the following assumption. 
% \SA{The following assumption can be removed if we need to save up space!}
\begin{assumption}\label{assumption: norm}
We assume $\int_{\mathbb{R}^n} h(x)\,dx = 1$ and $\mu(\mathcal{X}) = 1$ for the classes $\tilde{\mathcal{P}}_{c_1, c_2, h}$ and $\tilde{\mathcal{P}}_{c_1, c_2, \mu}$, respectively. Additionally, we assume that $0 \leq c_1 < 1 < c_2$, and $\frac{c_2 - c_1}{1 - c_1} \in \mathbb{N}$. 
% \SA{Don't we want to assume $c_1<1<c_2$ for the latter case?}\HG{We do! I initially, framed the assumption in this way: For $\tilde{\mathcal{P}}_{c_1, c_2, h}$, we assume $\int_{\mathbb{R}^n} h(x)\,dx = 1$, $0 \leq c_1 < 1 < c_2$, and $\frac{c_2 - c_1}{1 - c_1} \in \mathbb{N}$. In the geneal sample space case, the condition   $\int_{\mathbb{R}^n} h(x)\,dx = 1$ is replaced by $\mu(\mathcal{X}) = 1$. Isn't this okay? } \SA{I changed it!}
\end{assumption}
The following proposition shows that an additional condition is required to exclude trivial samplers.
\begin{proposition}\label{prop: triviality functional}
Let $g$ be a trade-off function. Under Assumption~\ref{assumption: norm},  if \( 1 + g^*(-e^\beta) > \frac{(c_2 - c_1 e^\beta)(1 - c_1)}{c_2 - c_1} \) for all \( \beta \geq 0 \), then \( \mathbf{Q}(P) = P \) satisfies \( g \)-FLDP. 
% \SA{$\mathbf{Q}^I$ do we need this notation later? If not, remove it and write $\mathbf{Q}(P) = P $. Avoid using unnecessary notation!}\HG{right, no we don't need it in the main body. I changed it.}
\end{proposition}
In view of this result, we assume throughout that for any given trade-off function $g$ the following condition holds:  \begin{equation}\label{eqn: non-trivial condition condition functional}
    \exists\, \beta \geq 0 \quad \text{such that} \quad 1 + g^*(-e^\beta) \leq \tfrac{(c_2 - c_1 e^\beta)(1 - c_1)}{c_2 - c_1}.
\end{equation}
Now, we are in order to state our first technical result. 
%The following theorem provides the optimal value and an optimal sampler for problem~\eqref{eqn: global minimax functional problem} for the continuous sample space $\X = \R^n$.
\begin{theorem}\label{thm: continuous global functional} 
Let $\tilde{\mathcal{P}}=\tilde{\mathcal{P}}_{c_1,c_2,h}$. Under Assumption \ref{assumption: norm}, the sampler $\mathbf{Q}_{c_1,c_2,h,g}^\star$ defined as a continuous distribution whose density is given by
\begin{equation}\label{eqn: continuous global functional opt mech}
   q^\star_g(P) (x)\coloneqq \lambda^\star_{c_1,c_2,g} p(x) + \left(1 - \lambda^\star_{c_1,c_2,g}\right) h(x), \quad \lambda^\star_{c_1,c_2,g} = \inf_{\beta \geq 0} 
\tfrac{e^\beta + \frac{c_2 - c_1}{1 - c_1} ( 1 + g^*(-e^\beta)) - 1}
     {(1 - c_1)e^\beta + c_2 - 1},
\end{equation}
satisfies $g$-FLDP, and is minimax-optimal under any \( f \)-divergence, that is
\begin{equation}\label{eqn: continuous global functional optimal value}
    \sup_{P \in \tilde{\mathcal{P}}} D_f\big(P \,\|\, \bQ^\star_{c_1,c_2,h,g}(P)\big)
    = \mathcal{R}\big(\mathcal{Q}_{\mathbb{R}^n, \tilde{\mathcal{P}}, g}, \tilde{\mathcal{P}}, f\big)
    = \frac{1 - r_1}{r_2 - r_1} f(r_2) + \frac{r_2 - 1}{r_2 - r_1} f(r_1),
\end{equation}
% where \( \beta^\star_g \) is the optimizer of the expression for $\lambda^\star_g$ in~\eqref{eqn: continuous global functional opt mech}, and
% \[
% r_1 = \frac{c_1}{c_2 - c_1} \cdot \frac{(1 - c_1)e^{\beta^\star_g} + c_2 - 1}{1 - (1 + g^*(-e^{\beta^\star_g}))}, \quad
% r_2 = \frac{c_2}{c_2 - c_1} \cdot \frac{(1 - c_1)e^{\beta^\star_g} + c_2 - 1}{e^{\beta^\star_g} + \frac{c_2 - 1}{1 - c_1} (1 + g^*(-e^{\beta^\star_g}))}.
% \]
for $r_1=\frac{c_1}{1-(1-c_1)\lambda_{c_1,c_2,g}^\star}$ and $r_2=\frac{c_2}{(c_2-1)\lambda_{c_1,c_2,g}^\star+1}$.
\end{theorem}


We note that the non-triviality condition in Proposition~\ref{prop: triviality functional} is equivalent to \( \lambda_{c_1,c_2,g}^\star \leq 1 \), which allows a natural interpretation of the sampler $\mathbf{Q}_{c_1,c_2,h,g}^\star$: It samples from \( P \) with probability \( \lambda_{c_1,c_2,g}^\star \) and from the reference distribution with probability \(1- \lambda_{c_1,c_2,g}^\star \). This probabilistic mixture reflects the privacy–utility trade-off: a larger $ \lambda^\star_{c_1, c_2, g} $ yields outputs closer to $ P $, while sampling from $ h $ introduces the randomness required for privacy.
The same mixture structure yields a minimax-optimal sampler for the discrete case where \( \mathcal{X} = [k] \) and \( \tilde{\mathcal{P}} = \mathcal{P}([k]) \).
\begin{theorem}\label{thm: discrete global functional}
Let \( \tilde{\mathcal{P}} = \mathcal{P}([k]) \) for some \( k \in \mathbb{N} \), and let \( \mu_k \) denote the uniform distribution on \([k]\). Then the sampler defined as
\[
\bQ^\star_{k,g}(P) \coloneqq \lambda^\star_{k,g} P + (1 - \lambda^\star_{k,g}) \mu_k, \quad 
\lambda^\star_{k,g} = \inf_{\beta \geq 0} 
\tfrac{e^\beta + k ( 1 + g^*(-e^\beta)) - 1}
     {e^\beta + k - 1},
\]
satisfies $g$-FLDP, and is minimax-optimal under any \( f \)-divergence, that is
\[
\sup_{P \in \tilde{\mathcal{P}}} D_f\big(P \,\|\, \bQ^\star_{k,g}(P)\big)
= \mathcal{R}\big(\mathcal{Q}_{[k], \tilde{\mathcal{P}}, g}, \tilde{\mathcal{P}}, f\big).
\]
\end{theorem}


%Similar to Theorem~\ref{thm: continuous global functional}, the optimal value of the minimax problem in the finite setting can be expressed in closed form. Theorem~\ref{thm: discrete global functional} states only the optimal sampler. The optimal value can similarly be derived in all subsequent results where it is omitted.


Theorems~\ref{thm: continuous global functional} and~\ref{thm: discrete global functional} provide optimal samplers for continuous and finite sample spaces, respectively. The next theorem addresses the same task for general sample spaces under the assumption of decomposability.
 

\begin{theorem}\label{thm: global functional}
% Let \( \tilde{\mathcal{P}} = \tilde{\mathcal{P}}_{c_1, c_2, \mu} \) under Assumption \ref{assumption: norm}. Suppose \( \mu \) is \(( \alpha,\frac{1}{\alpha},1) \)-decomposable with \( \alpha = \frac{1 - c_1}{c_2 - c_1} \). Then the mechanism
% \begin{equation}\label{eqn: global functional opt mech}
% \bQ^\star_g(P) \coloneqq \lambda^\star_g P + (1 - \lambda^\star_g) \mu, \quad \lambda^\star_g = \inf_{\beta \geq 0} 
% \frac{e^\beta + \frac{c_2 - c_1}{1 - c_1} (1 + g^*(-e^\beta)) - 1}
%      {(1 - c_1)e^\beta + c_2 - 1},
% \end{equation}
% belongs to \( \mathcal{Q}_{\mathcal{X}, \tilde{\mathcal{P}}, g} \) and is minimax-optimal under any \( f \)-divergence:
% \begin{equation}\label{eqn: global functional optimal value}
%     \sup_{P \in \tilde{\mathcal{P}}} D_f\big(P \,\|\, \bQ^\star_g(P)\big)
%     = \mathcal{R}\big(\mathcal{Q}_{\mathcal{X}, \tilde{\mathcal{P}}, g}, \tilde{\mathcal{P}}, f\big).
% \end{equation}
Let \( \tilde{\mathcal{P}} = \tilde{\mathcal{P}}_{c_1, c_2, \mu} \) under Assumption~\ref{assumption: norm}, and suppose \( \mu \) is \( (\alpha, \frac{1}{\alpha}, 1) \)-decomposable with \( \alpha = \frac{1 - c_1}{c_2 - c_1} \). Then the sampler defined as \( \bQ^\star_{c_1,c_2,\mu,g}(P) = \lambda^\star_{c_1,c_2,g} P + (1 - \lambda^\star_{c_1,c_2,g}) \mu \) satisfies $g$-FLDP, and is minimax-optimal with respect to any $f$-divergence, where $\lambda_{c_1,c_2,g}^\star$ is defined in Theorem~\ref{thm: continuous global functional}.
% i.e., it is minimax-optimal under any \( f \)-divergence:

\end{theorem}


Similar to Theorem~\ref{thm: global functional}, all subsequent results can be extended to a general sample space. For clarity of presentation, however, we state them in the continuous setting, while all proofs are given in full generality. 
  % \SA{This paragraph needs to be edited.}\HG{I changed and made it shorter. Is it better now?}




% As we extended the continuous sample space $\mathbb{R}^n$ in Theorem~\ref{thm: continuous global functional} to an arbitrary measurable space $\mathcal{X}$ in Theorem~\ref{thm: global functional}, 


Having established the optimal sampler for general \( g \)-FLDP, we present specific instantiations of the result for two widely studied settings: pure LDP and $\nu$-GLDP. First, we derive the optimal $\eps$-LDP sampling algorithm by setting $g=g_\eps$ in Theorem~\ref{thm: continuous global functional}, as stated in the following corollary.

% \begin{corollary}[$g_\eps$-FLDP as a special case]\label{cor: special case of f-LDP}\HG{check the functional LDP non-triviality condition!}
% Let $\tilde{\mathcal{P}} = \tilde{\mathcal{P}}_{c_1, c_2, \mu}$ satisfy the normalization condition~\eqref{eqn: normalization condition} and non-triviality conditions (\ref{eqn: non-trivial condition condition functional}), and assume that $\frac{c_2 - c_1}{1 - c_1} \in \mathbb{N}$. Suppose the reference measure $\mu$ is $(\alpha,\frac{1}{\alpha},1)$-decomposable with $\alpha = \frac{1 - c_1}{c_2 - c_1}$. Then the mechanism
% \begin{equation}
%     \bQ^\star_{g_\varepsilon}(P) \coloneqq \lambda^\star_{g_\varepsilon} P + (1 - \lambda^\star_{g_\varepsilon})\mu
% \end{equation}
% belongs to the class $\mathcal{Q}_{\mathcal{X}, \tilde{\mathcal{P}}, g_\varepsilon}$ and is minimax-optimal with respect to any $f$-divergence $D_f$, that is,
% \begin{equation}
%     \sup_{P \in \tilde{\mathcal{P}}} D_f\big(P \, \| \, \bQ^\star_{g_\varepsilon}(P)\big)
%     = \mathcal{R}\big(\mathcal{Q}_{\mathcal{X}, \tilde{\mathcal{P}}, g_\varepsilon}, \tilde{\mathcal{P}}, f\big),
% \end{equation}
% where 
% \[
% \lambda^\star_{g_\varepsilon} = \frac{e^\varepsilon - 1}{(1 - c_1)e^\varepsilon + c_2 - 1}.
% \]
% \end{corollary}

\begin{corollary}\label{cor: special case of f-LDP}
Let $\tilde{\mathcal{P}}=\tilde{\mathcal{P}}_{c_1,c_2,h}$. Under Assumption \ref{assumption: norm}, the sampler $\mathbf{Q}_{c_1,c_2,h,g_\eps}^\star$ defined as a continuous distribution whose density is given by
   $q^\star_{g_\eps}(P) (x)\coloneqq \lambda^\star_{c_1,c_2,g_\eps} p(x) + \left(1 - \lambda^\star_{c_1,c_2,g_\eps}\right) h(x)$, with $ \lambda^\star_{c_1,c_2,g_\eps} = \frac{e^\varepsilon - 1}{(1 - c_1)e^\varepsilon + c_2 - 1},
$
satisfies $g_\eps$-FLDP, and is minimax-optimal with respect to any \( f \)-divergence, that is
\begin{equation}\label{eqn: continuous global functional optimal value g_eps}
    \sup_{P \in \tilde{\mathcal{P}}} D_f\big(P \,\|\, \bQ^\star_{c_1,c_2,h,g_\eps}(P)\big)
    = \mathcal{R}\big(\mathcal{Q}_{\mathbb{R}^n, \tilde{\mathcal{P}}, g_\eps}, \tilde{\mathcal{P}}, f\big)
    = \frac{1 - r_1}{r_2 - r_1} f(r_2) + \frac{r_2 - 1}{r_2 - r_1} f(r_1),
\end{equation}
for $
r_1 = c_1 \cdot \frac{(1 - c_1)e^{\eps} + c_2 - 1}{c_2 - c_1 }$, and $
r_2 = \frac{c_2}{c_2 - c_1} \cdot \frac{(1 - c_1)e^{\eps} + c_2 - 1}{e^{\eps}}$.
% Let \( h: \mathbb{R}^n \to [0, \infty) \) be a probability density function, and let \( H \) be the corresponding probability measure. Define \( \tilde{\mathcal{P}} = \tilde{\mathcal{P}}_{c_1, c_2, h} \) under Assumption \ref{assumption: norm}. Then the sampler $
%     \bQ^\star_{g_\varepsilon}(P) \coloneqq \lambda^\star_{g_\varepsilon} P + (1 - \lambda^\star_{g_\varepsilon}) H$, with 
% $\lambda^\star_{g_\varepsilon} = \frac{e^\varepsilon - 1}{(1 - c_1)e^\varepsilon + c_2 - 1}$
% belongs to \( \mathcal{Q}_{\mathbb{R}^n, \tilde{\mathcal{P}}, g_\varepsilon} \) and is minimax-optimal under any \( f \)-divergence:
% \[
% \sup_{P \in \tilde{\mathcal{P}}} D_f\big(P \,\|\, \bQ^\star_{g_\varepsilon}(P)\big)
% = \mathcal{R}\big(\mathcal{Q}_{\mathbb{R}^n, \tilde{\mathcal{P}}, g_\varepsilon}, \tilde{\mathcal{P}}, f\big) = \frac{1 - r_1}{r_2 - r_1} f(r_2) + \frac{r_2 - 1}{r_2 - r_1} f(r_1),
% \]
% where
% $
% r_1 = c_1 \cdot \frac{(1 - c_1)e^{\eps} + c_2 - 1}{c_2 - c_1 }$, and $
% r_2 = \frac{c_2}{c_2 - c_1} \cdot \frac{(1 - c_1)e^{\eps} + c_2 - 1}{e^{\eps}}$.
\end{corollary} 
 % \textcolor{blue}{[Maybe we can just remove the corollary and add 'Theorem 3.4 recovers the known result of park et al.' below Theorem 3.4, if we need space]} \SA{Since we need this result for experiments, I think we better keep pure LDP in the main body, but we don't have to present as a corollary. We can mention the result in the text and say it is exactly Theorem blah in Park et al. }\HG{I wrote it with the new notation. If we are short of space, or if stating as a corollary makes it complicated I can move it into text.}

Indeed, by framing $\eps$-LDP as a special case of functional LDP, Corollary \ref{cor: special case of f-LDP} reproduces the minimax risk of the optimal 
$\eps$-LDP sampler in \citep[Theorem 3.3]{park2024exactly}. %highlighting the generality of the functional LDP framework.
Next, we focus on the $\nu$-GLDP as a special case of functional LDP. 
\begin{corollary}[$\nu$-GLDP as a special case]\label{cor: special case GDP}
Let \( \tilde{\mathcal{P}} = \tilde{\mathcal{P}}_{c_1, c_2, h} \).
% Let \( h: \mathbb{R}^n \to [0, \infty) \) be a probability density function, and \( \tilde{\mathcal{P}} = \tilde{\mathcal{P}}_{c_1, c_2, h} \). 
% \SA{We already said $h$ is a pdf, so do we still need to say it again?} \HG{right, we don't need to say it again. I changed it.} 
Under Assumption \ref{assumption: norm}, the sampler  \eqref{eqn: continuous global functional opt mech} for $g=G_\nu$ belongs to $\mathcal{Q}_{\R^n, \tilde{\mathcal{P}}, G_\nu}$ and is minimax-optimal with respect to any $f$-divergence, that is,
\begin{equation}
    \sup_{P \in \tilde{\mathcal{P}}} D_f\big(P \, \| \, \bQ^\star_{G_\nu}(P)\big)
    = \mathcal{R}\big(\mathcal{Q}_{\R^n, \tilde{\mathcal{P}},G_\nu}, \tilde{\mathcal{P}}, f\big).
\end{equation}
\end{corollary}
% \HG{ I know  $\lambda_{c_1,c_2,G_\nu}^\star$ does not have a very clean expression and makes the result complicated. But is it okay to mention Corollary \ref{cor: special case GDP} in the above format? I feel it does not give an additional information than Theorem \ref{thm: continuous global functional}. I mean is it okay not to mention its specific $\lambda_{c_1,c_2,G_\nu}^\star$? Should I refer to appendix for how to derive the sampler in detail?} \SA{I think so! Maybe we need to add a line after corollary say that the explicit description of the sampler is given in Appendix..}
The explicit description of the optimal sampler under $\nu$-GLDP is given in Appendix~\ref{proof: GDP}.

As outlined in the introduction, the global minimax formulation under functional LDP serves as the basis for our local analysis. We now turn to the local minimax setting, first under general  $g$-FLDP (Section~\ref{sec: local functional}) and then under pure LDP (Section~\ref{sec: local pure}), leveraging Theorem \ref{thm: continuous global functional} and Corollary~\ref{cor: special case of f-LDP}, respectively. 


% In this section, we studied the global minimax sampler under functional LDP setting and derived optimal mechanisms and values for any $f$-divergence. In the next section, we turn to the local minimax setting, where input distributions lie within a neighborhood of a reference distribution $P_0$.

% where 

% \begin{align*}
% \lambda^\star_{G_\nu} & =\inf_{\beta \geq 0} 
%     \frac{e^\beta + \Big(\frac{c_2 - c_1}{1 - c_1}\Big) \bigg( 1 -e^\beta \Phi\left(-\frac{\nu}{2} - \frac{\beta}{\nu}\right) - \Phi\left(-\frac{\nu}{2} + \frac{\beta}{\nu}\right) \bigg) - 1}
%          {(1 - c_1)e^\beta + c_2 - 1}\\
%          & = \inf_{\beta \geq 0} 
%     \frac{e^\beta + \Big(\frac{c_2 - c_1}{1 - c_1}\Big) \bigg( \Phi\left(\frac{\nu}{2} - \frac{\beta}{\nu}\right)  -e^\beta \Phi\left(-\frac{\nu}{2} - \frac{\beta}{\nu}\right)  \bigg) - 1}
%          {(1 - c_1)e^\beta + c_2 - 1}.    
% \end{align*}




\vspace{0.1cm}
\section{Local minimax-optimal sampling under functional LDP}\label{sec: local functional}
\vspace{0.1cm}
In this section, we begin by formalizing the local minimax setting using a neighborhood structure induced by the $E_\gamma$-divergence. We then develop a general framework for identifying local minimax-optimal samplers over arbitrary sample spaces.




The local minimax-optimal sampling problem adopts the same structural form as the global minimax problem, with a key distinction: instead of selecting the worst-case distribution from the entire space of the general class \( \tilde{\mathcal{P}} \), the search is restricted to a neighborhood around a fixed distribution $P_0$. This reference distribution $P_0$  may represent, for example, a publicly available dataset held by an individual or shared collaboratively in a distributed setting. Formally, the local minimax problem under \( g \)-FLDP is defined as:
\begin{equation}\label{eqn: local minimax functional problem} \mathcal{R}\big(\mathcal{Q}_{\mathcal{X},\tilde{\mathcal{P}},g},N_\gamma(P_0),f\big) \coloneqq \inf_{\bQ \in \mathcal{Q}_{\mathcal{X},\tilde{\mathcal{P}},g}} \hspace{0.1cm} \sup_{P \in N_\gamma(P_0)}\hspace{0.1cm} D_f\big(P \, \| \, \bQ(P)\big),
\end{equation}
where $N_\gamma(P_0)$ denotes a neighborhood around $P_0$, recently adopted by  \citet{feldman2024instance}, and is defined for any $\gamma\geq 1$ as
% \begin{equation}\label{eqn: ngb}N_\gamma(P_0) \coloneqq  \left\{ P \in \mathcal{P}(\mathcal{X}) :\quad \frac{1}{\gamma} \leq \frac{P(x)}{P_0(x)} \leq \gamma \quad \forall x \in \mathcal{X} \right\} \end{equation}
\begin{equation}\label{eqn: ngb}N_\gamma(P_0) \coloneqq  \left\{ P \in \mathcal{P}(\mathcal{X}) :\quad E_\gamma(P \, \| \, P_0) =  E_\gamma(P_0 \, \| \, P) = 0  \right\}. \end{equation}
The zero-\( E_\gamma \) neighborhood generalizes the classical notion of total variation distance. Specifically, for \( \gamma \geq 1 \), \( E_\gamma(P \| Q) = 0 \) implies \( D_{\mathrm{TV}}(P, Q) \leq 1 - \tfrac{1}{\gamma} \) \citep[Proposition~4]{liu2016e_}. Moreover, we can straightforwardly extend this neighborhood to the more general case by replacing \( E_\gamma(P \| P_0) = 0 \) with \( E_\gamma(P \| P_0) \leq \zeta \), meaning the likelihood ratio lies in \( [1/\gamma, \gamma] \) with probability at least \( 1 - \zeta \) under \( P_0 \). This, in turn, enables further generalization to any \( f \)-divergence with a twice differentiable function \( f \), which is a common trick in information theory \citep{polyanskiy2015dissipation,raginsky2016strong,zamanlooy2023strong,asoodeh2024contraction}.
As in the global case (Section~\ref{sec: global functional}), we assume \( \gamma \in \mathbb{N} \) for the same technical reasons. Theorem~\ref{thm: local functional} presents the corresponding locally minimax-optimal sampler, connecting the global minimax analysis of Section~\ref{sec: global functional} with the local setting under functional LDP. 



% \begin{figure}[ht]
%   \centering
%   \includegraphics[width=0.6\linewidth]{figures/local_functional.png}
%   \caption{Illustration of the optimal mechanism $\bQ^\star_{g, N_\gamma(P_0)}$ described in Theorem~\ref{thm: local functional}.}
%   \label{fig:local-functional}
% \end{figure}

% \begin{wrapfigure}{r}{0.3\textwidth}
%   \centering
%   \includegraphics[width=0.45\textwidth]{figures/local_functional.png}
 % adjust as needed
%   \caption{Illustration of the optimal mechanism $\bQ^\star_{g, N_\gamma(P_0)}$.}
%   \label{fig:local-functional}
% \end{wrapfigure}
\begin{theorem}\label{thm: local functional}
Let \( P_0 \) be a continuous distribution on \( \mathbb{R}^n \). Define \( N_\gamma(P_0) \)  as in~\eqref{eqn: ngb}, and assume \( N_\gamma(P_0) \subseteq \tilde{\mathcal{P}} \).
Let \( \bQ^\star_g \) denote the global sampler from~\eqref{eqn: continuous global functional opt mech}, instantiated with parameters \( c_1 = \frac{1}{\gamma} \), \( c_2 = \gamma \), and \( h = p_0 \).
We define the sampler \( \bQ^\star_{g, N_\gamma(P_0)} \) as
\[
\bQ^\star_{g, N_\gamma(P_0)}(P) \coloneqq
\begin{cases}
\bQ^\star_g(P), & \text{if } P \in N_\gamma(P_0), \\[1pt]
\bQ^\star_g(\hat{P}), & \text{otherwise},
\end{cases}
\]
where \( \hat{P} \in N_\gamma(P_0) \) is a distribution that minimizes \(  D_f(P \,\|\, P') \) over all $P' \in N_\gamma(P_0)$.
Then \( \bQ^\star_{g, N_\gamma(P_0)} \in \mathcal{Q}_{\mathbb{R}^n, \tilde{\mathcal{P}}, g} \) and is locally minimax-optimal under any \( f \)-divergence, that is
\[
\sup_{P \in N_\gamma(P_0)} D_f\big(P \,\|\, \bQ^\star_{g, N_\gamma(P_0)}(P)\big)
= \mathcal{R}\big(\mathcal{Q}_{\mathbb{R}^n, \tilde{\mathcal{P}}, g}, N_\gamma(P_0), f\big).
\]
\end{theorem}
Theorem~\ref{thm: local functional} establishes that the local minimax-optimal sampler for the problem in~\eqref{eqn: local minimax functional problem} coincides with the global solution in~\eqref{eqn: global minimax functional problem} when the universe is restricted to the $N_\gamma$ neighborhood around \( P_0 \). In this sense, the global solution presented in Theorem~\ref{thm: continuous global functional} directly yields the local solution in Theorem~\ref{thm: local functional}. 
% By setting \( g = g_\varepsilon \) in Theorem~\ref{thm: local functional}, we recover the local minimax sampler for pure LDP, as shown in Corollary~\ref{cor: special case of f-LDP}. In the next section, however, we focus specifically on the pure LDP setting and introduce a mechanism with improved pointwise performance.

By instantiating Theorem~\ref{thm: local functional} with $g = g_\eps$, we obtain an $\eps$-LDP sampler that is locally minimax-optimal. However, as we demonstrate in the next section, this sampler can still be improved in a \textit{pointwise} sense.

%In the next section, we focus on the local minimax samplers under pure LDP and identify a fundamentally different construction that is \textit{pointwise} superior to the pure-LDP sampler obtained by instantiating 




%As a special case, setting \( g = g_\varepsilon \) in Theorem~\ref{thm: local functional} yields a local minimax sampler for pure LDP. Section \ref{sec: local pure} focuses on local minimax under $\eps$-LDP and presents a sampler with better pointwise performance than this instantiation.


\vspace{0.1cm}
\section{Local minimax-optimal sampling under pure LDP}\label{sec: local pure}
\vspace{0.1cm}

In this section, we turn our attention to pure-LDP sampling and develop a new framework that is fundamentally different from the approach presented in the previous section. Our goal is to identify a sampler that is \textit{pointwise} better than the one obtained from Theorem~\ref{thm: local functional} with $g = g_\eps$.  The construction of such sampler is presented in the next theorem.
%The proposed sampler is conceptually distinct from that of the previous section and is provably pointwise better than the construction given by Theorem~\ref{thm: local functional}.



%Theorem~\ref{thm: local functional}, instantiated with $g=g_\eps$, provides a local minimax-optimal sampler under $\eps$-LDP. However, we show that there exists another local minimax sampler that achieves strictly better utility. 
%Under $\varepsilon$-LDP, the local minimax problem is formulated analogously to~\eqref{eqn: local minimax functional problem} as follows:
%\begin{equation}\label{eqn: local minimax pure problem}
    %\mathcal{R}\big(\mathcal{Q}_{\mathcal{X},\tilde{\mathcal{P}},\eps},N_\gamma(P_0),f\big) \coloneqq \inf_{\bQ \in \mathcal{Q}_{\mathcal{X},\tilde{\mathcal{P}},\eps}} \hspace{0.1cm} \sup_{P \in N_\gamma(P_0)}\hspace{0.1cm} D_f\big(P \, \| \, \bQ(P)\big),
%\end{equation}
% \begin{figure}[ht]
%   \centering
%   \includegraphics[width=0.6\linewidth]{figures/local_pure.png}
%   \caption{Illustration of the optimal mechanism $\bQ^\star_{\eps, N_\gamma(P_0)}$ described in Theorem~\ref{thm: local pure}.}
%   \label{fig:local-pure}
% \end{figure}

% \setlength{\textfloatsep}{0.4cm}
% \begin{figure}[t]
%   \centering
%   \includegraphics[width=0.7\linewidth]{figures/local_functional_pure_visualization.png}
%   \caption{
%     Illustrations of the optimal mechanisms described in Theorems~\ref{thm: local functional} and~\ref{thm: local pure}. 
%     \textbf{Left:} $\bQ^\star_{g, N_\gamma(P_0)}$ for functional LDP. 
%     \textbf{Right:} $\bQ^\star_{\varepsilon, N_\gamma(P_0)}$ for pure LDP. In the right figure, $M_\varepsilon$ is the set onto which projecting any distribution in $\ngb$ ensures $\varepsilon$-LDP and minimizes the local minimax risk.
%   }
%   \label{fig:local-mechanisms}
% \end{figure}
%In the following theorem, we propose a sampler that achieves the optimal value of~\eqref{eqn: local minimax pure problem}.
\setlength{\textfloatsep}{0.6cm}
\begin{figure}[t]
  \centering
  \includegraphics[width=0.8\linewidth]{figures/Theorems_Visualization.pdf}
  \caption{
    Illustrations of the optimal samplers described in Theorems~\ref{thm: local functional} and~\ref{thm: local pure}. 
    \textbf{Left:} $\bQ^\star_{g, N_\gamma(P_0)}$ for functional LDP. 
    \textbf{Right:} $\bQ^\star_{\varepsilon, N_\gamma(P_0)}$ for pure LDP. Following \citep[Proposition 3.4]{park2024exactly}, $M_\eps$  is defined as $M_\eps \coloneqq \{ Q \in \mathcal{C}(\mathbb{R}^n): \hspace{0.1cm} \frac{\gamma + 1}{\gamma + e^\varepsilon} \, p_0(x)  \leq q(x) \leq \frac{(\gamma + 1)e^\varepsilon}{\gamma + e^\varepsilon} \, p_0(x), \hspace{0.1cm}\forall x \in \mathbb{R}^n \}$.} 
  \label{fig:local-mechanisms}
  % \SA{Do we have the definition of $M_\eps$ in appendix?The caption reads a bit cryptically! }\HG{I don't have the definition in the appendix right now, but it is basically the set that $f$-divergence projection of $P$ into $M_\eps$ gives us the optimal sampler. (Similar to Hyun-Young's result but for local neighborhood) Should I include it in the appendix? I can create a section in appendix and formally show that optimal clipping is equivalent to $f$-divergence projection similar to Proposition 3.4 in \citep{park2024exactly}.} \SA{Maybe just restate the Hyun-Young result and clearly defined $M_\eps$ and cite to it here.}\HG{done! Is it okay now?} \SA{Awesome!}
  % is the set onto which projecting any distribution in $\ngb$ ensures $\varepsilon$-LDP and minimizes the local minimax risk.
\end{figure}
\begin{theorem}\label{thm: local pure}
Let \( P_0 \) be a continuous distribution on \( \mathbb{R}^n \) with density \( p_0\) and $N_\gamma(P_0)$ be the neighborhood around $P_0$ defined in \eqref{eqn: ngb}. Assuming \( N_\gamma(P_0) \subseteq \tilde{\mathcal{P}} \), let 
\( \bQ^\star_\varepsilon \in \mathcal{Q}_{\mathbb{R}^n, N_\gamma(P_0), \varepsilon} \) be  a sampler with \( \bQ^\star_\varepsilon(P) \)  having density
\[
q(x) 
= \mathrm{clip}\left(
\frac{1}{r_P} \, p(x); \,
\frac{\gamma + 1}{\gamma + e^\varepsilon} \, p_0(x),\ 
\frac{(\gamma + 1)e^\varepsilon}{\gamma + e^\varepsilon} \, p_0(x)
\right),
\]
where \( r_P \) being the normalizing constant ensuring \( \int_{\mathbb{R}^n} q(x) \, dx = 1 \).
Furthermore, define the extended sampler 
\( \bQ^\star_{\varepsilon, N_\gamma(P_0)} \) by
\[
\bQ^\star_{\varepsilon, N_\gamma(P_0)}(P) \coloneqq
\begin{cases}
\bQ^\star_\varepsilon(P), & \text{if } P \in N_\gamma(P_0), \\[1pt]
\bQ^\star_\varepsilon(\hat{P}), & \text{otherwise},
\end{cases}
\]
where \( \hat{P} \in N_\gamma(P_0) \) is a distribution that minimizes \(  D_f(P \,\|\, P') \) over all $P' \in N_\gamma(P_0)$.
Then, we have \( \bQ^\star_{\varepsilon, N_\gamma(P_0)} \in \mathcal{Q}_{\mathbb{R}^n, \tilde{\mathcal{P}}, \varepsilon} \) and 
$
\sup_{P \in N_\gamma(P_0)} D_f\big(P \,\|\, \bQ^\star_{\varepsilon, N_\gamma(P_0)}(P)\big)
= \mathcal{R}\big(\mathcal{Q}_{\mathbb{R}^n, \tilde{\mathcal{P}}, \varepsilon}, N_\gamma(P_0), f\big)$.
%is locally minimax-optimal under any \( f \)-divergence:
\end{theorem}


%Figure~\ref{fig:local-mechanisms} provides a visual representation of Theorems~\ref{thm: local functional} and~\ref{thm: local pure}. In both cases, the local minimax-optimal samplers follow a two-step procedure: first, project the input distribution onto the neighborhood \( \ngb \) (if it lies outside), and then apply the corresponding optimal sampler within~\( \ngb \).

Similar to Theorem~\ref{thm: local functional}, the above theorem defines a sampler that distinguishes between the cases \( P \in N_\gamma(P_0) \) and \( P \notin N_\gamma(P_0) \). In the latter case, both theorems are aligned in that the sampling distribution depends on \( \hat{P} \), the closest distribution to \( P \) within \( N_\gamma(P_0) \). However, for the case \( P \in N_\gamma(P_0) \), the two constructions differs fundamentally. Theorem~\ref{thm: local functional} assigns a \textit{linear} sampler, whereas Theorem~\ref{thm: local pure} carefully constructs a non-linear sampler tailored to this regime. In fact, this non-linear sampler is obtained by projecting $P$ onto a mollifier $M_\eps$; see proof of the theorem in Appendix~\ref{proof: local pure} and Figure~\ref{fig:local-mechanisms} for a visual representation. This projection is carried out through a clipping operation, which guarantees that the resulting distribution $ \mathbf{Q}^\star_\varepsilon(P) $ satisfies the LDP constraint while maintaining a close match to the original distribution $ P $ in terms of $ f $-divergence. This approach yields a nonlinear transformation of $ P $ that carefully balances privacy constraints and utility preservation.
It can be verified---using a mutatis mutandis adaptation of \cite[Proposition C.5]{park2024exactly}--- that the non-linear sampler pointwise outperforms the linear one under the same privacy guarantees; that is, for any $\eps\geq 0$ and $P\in N_\gamma(P_0)$
\[
D\big(P \,\|\, \bQ^*_{\eps}(P)\big) \leq D\big(P \,\|\, \bQ^*_{g_\eps}(P)\big).
\]
In fact, the non-linear sampler outperforms the linear one because it is not only worst-case optimal but also instance-optimal: for each input distribution $P$ in our distribution class and any $f$-divergence $D_f$, it minimizes $D_f(P || \mathbf{Q})$ over all admissible $\mathbf{Q}$ satisfying $\varepsilon$-LDP (see Proposition \ref{prop: pointwise optimal Park et al.}). In contrast, the linear sampler uses a fixed transformation for all input distributions. This instance-specific optimization explains the empirical advantage of the non-linear sampler. As a result, the overall sampler in Theorem~\ref{thm: local pure} is pointwise better than the one in Theorem~\ref{thm: local functional}. This intuition is formalized in the result below. 



 %With the optimal sampler for the general $g$-FLDP setting in Section \ref{sec: local functional}, a local minimax-optimal sampler for pure LDP can be obtained by instantiating Theorem~\ref{thm: local functional} with $g = g_\varepsilon$, yielding the sampler $\bQ^\star_{g_\varepsilon,\ngb}$. However, the following proposition shows that the sampler $\bQ^\star_{\varepsilon,\ngb}$ proposed in Theorem~\ref{thm: local pure} achieves better pointwise performance; i.e., for any $P \in \ngb$, it achieves a smaller $f$-divergence between the input distribution and the sampler's output.


\begin{proposition}\label{prop: pointwise}
Let \( \tilde{\mathcal{P}} \) be the global universe and \( P_0 \) a probability measure on \( \mathbb{R}^n \), with \( N_\gamma(P_0) \subseteq \tilde{\mathcal{P}} \) for $\ngb$ defined as in~\eqref{eqn: ngb}. Let \( \bQ^\star_{\varepsilon,\ngb} \) be the optimal $\eps$-LDP sampler from Theorem~\ref{thm: local pure}, and  \( \bQ^\star_{g_\varepsilon,\ngb} \) the instantiation of Theorem~\ref{thm: local functional} with \( g = g_\varepsilon \). Then, for all \( P \in N_\gamma(P_0) \),
\[
D_f(P \,\|\, \bQ^\star_{\varepsilon,\ngb}) \leq D_f(P \,\|\, \bQ^\star_{g_\varepsilon,\ngb}).
\]
\end{proposition}
\vspace{0.1cm}
\section{Numerical results}\label{sec: numeric}
\vspace{0.1cm}
In this section, we numerically compare the worst-case \( f \)-divergence of our locally minimax sampler against the globally optimal sampler of~\citep{park2024exactly} under the $\varepsilon$-LDP setting. Experiments span both finite and continuous domains, evaluating KL divergence, total variation distance, and squared Hellinger distance across \( \varepsilon \in \{0.1, 0.5, 1, 2\} \). Additional results under \( \nu \)-GLDP appear in Appendix~\ref{appendix: experiment GDP local global discrete} and~\ref{appendix: experiment GDP local global continuous}. In addition, we report complementary results for the continuous case under pure LDP in Appendix~\ref{appendix: high dimension laplace}.





\subsection{Finite sample space}\label{experiment: finite}

We compare local and global minimax-optimal samplers in the finite setting \( \mathcal{X} = [k] \), where the global class is \( \tilde{\mathcal{P}}_{\textsf{global}} = \mathcal{P}([k]) = \tilde{\mathcal{P}}_{0, k, \mu_k} \) for the uniform measure \( \mu_k \). The local neighborhood is defined as \( \tilde{\mathcal{P}}_{\textsf{local}} = \mathcal{N}_\gamma(\mu_k) = \tilde{\mathcal{P}}_{\frac{1}{\gamma},\, \gamma,\, \mu_k} \) with \( \gamma = \frac{k}{2} - 1 \), ensuring that all local neighborhood assumptions are satisfied and \( \mathcal{N}_\gamma(\mu_k) \subseteq \mathcal{P}([k]) \).
Indeed, both global and local worst-case \( f \)-divergences can be computed in closed form: the global risk is given by~\citep[Theorem 3.1]{park2024exactly}, while the local risk is derived from a finite-space version of Theorem~\ref{thm: local pure} (see Appendix~\ref{appendix: finite space}). Figure~\ref{fig: finite worst-case k = 20} compares the two for \( k = 20 \); additional results for other \( k \) values are provided in Appendix~\ref{appendix: finite space}. The local minimax-optimal sampler consistently outperforms the global minimax-optimal sampler across all \( \varepsilon \) and \( f \)-divergences.

 
% In this subsection, we compare the local and global minimax-optimal mechanisms in the finite space $\mathcal{X} = [k]$ with $\tilde{\mathcal{P}} = \mathcal{P}([k])$. We note that $\mathcal{P}([k]) = \Pt_{0,[k],\mu_k}$ for uniform reference distribution over sample space $[k]$.and the global universe is defined as $\tilde{\mathcal{P}}_{\textsf{global}} = \mathcal{P}([k])= \tilde{\mathcal{P}}_{0, k, \mu_k}$. We define the local neighborhood as $\tilde{\mathcal{P}}_{\textsf{local}} = \mathcal{N}_\gamma(\mu_k) = \tilde{\mathcal{P}}_{\frac{1}{\gamma},\, \gamma,\, \mu_k}$ for the same reference meausre $\mu_k$ and for $ \gamma = \frac{k}{2} -1$. For these definitions, so that \( \mathcal{N}_\gamma(\mu_k) \subseteq \mathcal{P}([k]) \) and all neighborhood assumptions are satisfied.

% In this setting, both the global and local worst-case $f$-divergences, $\mathcal{R}\big(\mathcal{Q}_{[k], \tilde{\mathcal{P}}, \varepsilon}, \tilde{\mathcal{P}}, f\big)$ and $\mathcal{R}\big(\mathcal{Q}_{[k], \tilde{\mathcal{P}}, \varepsilon}, \mathcal{N}_\gamma(\mu_k), f\big)$, can be computed in closed form. The global minimax risk is derived from~\citet{park2024exactly}, while the local minimax risk for neighborhood $\mathcal{N}_\gamma(\mu_k)$ is obtained from a finite-sample-space version of Theorem~\ref{thm: local pure} as detailed in Appendix \ref{appendix: finite space}.
%  % We define the global and local universes respectively as:
% % \[
% % \tilde{\mathcal{P}}_{\textsf{local}} = \mathcal{N}_\gamma(\mu_k) = \tilde{\mathcal{P}}_{\frac{1}{\gamma},\, \gamma,\, \mu_k} \quad \text{for } \gamma = \frac{k}{2} -1,
% % \qquad
% % \tilde{\mathcal{P}}_{\textsf{global}} = \mathcal{P}([k])= \tilde{\mathcal{P}}_{0, k, \mu_k},
% % \]
% We present the plot for \( k = 20 \) in Figure~\ref{fig: finite worst-case k = 20}; results for other values of \( k \) are provided in Appendix~\ref{appendix: finite space}. As shown in Figure~\ref{fig: finite worst-case k = 20}, the local minimax sampler from Theorem~\ref{thm: local pure} consistently outperforms the global minimax sampler of~\citet{park2024exactly} across all \( \varepsilon \) values and all $f$-divergences.

\setlength{\textfloatsep}{0.4cm}
\begin{figure}[ht]
  \centering
  \includegraphics[width=0.8\linewidth]{figures/result_finite_k20.pdf}
  \caption{Theoretical worst-case 
$f$-divergences of global and local minimax samplers under the pure LDP setting with uniform reference distribution $\mu_k$ over finite space ($k = 20$).
% (Left: KL divergence; Center: Total variation distance; Right: Squared Hellinger distance.)
}
  \label{fig: finite worst-case k = 20}
\end{figure}

% Commented Fig
\subsection{Continuous sample space}\label{subsec: experiment 1d continuous}

In the continuous setting with \( \mathcal{X} = \mathbb{R} \), we fix the universe \( \tilde{\mathcal{P}}_{\textsf{local}} \) and evaluate the empirical worst‐case \( f \)-divergence over 100 randomly generated client distributions \( P_1, \ldots, P_{100} \in \tilde{\mathcal{P}}_{\textsf{local}} \).  Each \( P_j \) represents a client and is constructed as a mixture of a random number of one–dimensional Laplace components with scale parameter \( b = 1 \); the complete procedure for generating these distributions is described in detail in Appendix~\ref{appendix: details of continuous}.

We define the local and global universes as  
\( \tilde{\mathcal{P}}_{\textsf{local}} = \tilde{\mathcal{P}}_{1/3,\, 3,\, h_{\mathcal{L}}} \) and  
\( \tilde{\mathcal{P}}_{\textsf{global}} = \tilde{\mathcal{P}}_{1/9,\, 9,\, h_{\mathcal{L}}} \), 
where \( h_{\mathcal{L}} \) is the density of a Laplace distribution with mean zero and scale \( b = 1 \).  
Since \( \tilde{\mathcal{P}}_{\textsf{local}} \subseteq \tilde{\mathcal{P}}_{\textsf{global}} \), every client distribution \( P_j \) also belongs to the global universe; thus the input distributions are identical for both samplers.
We evaluate the empirical worst‐case divergence of each sampler as  
\( \max_{j \in [100]} D_f\!\bigl(P_j \,\|\, \mathbf{Q}(P_j)\bigr) \).  
The local minimax sampler follows Theorem~\ref{thm: local pure}, while the global sampler is the optimal sampler from~\citep[Theorem 3.3]{park2024exactly}.  
As illustrated in Figure~\ref{fig: continuous worst-case}, the local sampler consistently achieves lower worst‐case \( f \)-divergence than its global counterpart across all \( f \)-divergences and privacy levels \( \varepsilon \).








\setlength{\textfloatsep}{0.4cm}
\begin{figure}[ht]
  \centering
  \includegraphics[width=0.8\linewidth]{figures/result_1DLaplaceMix_1.pdf}
  \caption{Empirical worst-case 
$f$-divergences of global and local minimax samplers under the pure LDP setting, over 100 experiments on a 1-D Laplace mixture.
% (Left: KL divergence; Center: Total variation distance; Right: Squared Hellinger distance.)
}
  \label{fig: continuous worst-case}
\end{figure}


% \vspace{0.1cm}
\section{Limitations and future direction}\label{sec: conclusion}
\vspace{0.1cm}
% In this paper, we extended the locally private sampling framework of \citep{park2024exactly} in two directions. First, we extend their optimal pure LDP samplers to optimal samplers with functional LDP guarantee.  Second, we extend their global minimax formulation to the local one and apply it to reason about sampling in the presence of public data. Our results can potentially be applied to fine-tuning large models (such as language models) and unit testing.
%studied private sampling under different variants of LDP in a local minimax setting, where distributions lie in a neighborhood around a known non-private distribution \( P_0 \). We first characterize the global minimax risk and proposed optimal samplers under functional LDP for both continuous and discrete settings. Building on this, we define the local minimax problem under functional and pure LDP, showing that the local minimax risk is fully determined by the global minimax risk when the distribution class  is restricted to a neighborhood around $P_0$. We derive optimal samplers for both cases and demonstrate their effectiveness through numerical comparisons against global minimax baselines.


% \textbf{Limitation and future work.} 

% As noted in Section~\ref{sec: local functional}, our local minimax formulation defines neighborhoods via the condition \( E_\gamma(P \| P_0) = E_\gamma(P_0 \| P) = 0 \), following \citet{feldman2024instance}. A natural extension is the approximate neighborhood \( N_{\gamma,\delta}(P_0) \coloneqq \{ P \in \mathcal{P}(\mathcal{X}) : E_\gamma(P \| P_0) \leq \delta \text{ and } E_\gamma(P_0 \| P) \leq \delta \} \), which can further be generalized to neighborhoods defined based on  \( f \)-divergences for twice differentiable \( f \), defined as \( N_{f,\delta}(P_0) \coloneqq \{ P \in \mathcal{P}(\mathcal{X}) : D_f(P \| P_0) \leq \delta \text{ and } D_f(P_0 \| P) \leq \delta \} \) \citep{polyanskiy2015dissipation,raginsky2016strong,zamanlooy2023strong,asoodeh2024contraction}. These generalizations highlight the flexibility of our formulation and provide a clear pathway for extending the local minimax framework. In addition, while this work focuses on releasing a single sample per client, extending to the multi-sample setting is a natural direction for future work.



Our main contribution is the development of a minimax-optimal private sampler, validated through a series of experiments on synthetic data. However, the absence of evaluations on high-dimensional real-world datasets limits our understanding of its practical utility. A key reason for this is the computational complexity of our sampler (in particular the clipping function), which poses challenges for scalability in high-dimensional settings. Addressing this limitation---by developing more efficient implementations of our samplers or approximate variants, potentially leveraging techniques such as MCMC--- is an important direction for future work.


Beyond computational considerations, another important direction for future work is to generalize the formulation of the local neighborhood. The neighborhood in our local minimax formulation is defined using the $E_\gamma$-divergence: $P\in N_\gamma(P_0)$ if  \( E_\gamma(P \| P_0) = E_\gamma(P_0 \| P) = 0 \), as in \citep{feldman2024instance}. A natural extension is 
$N_{\gamma,\zeta}(P_0) \coloneqq \left\{ P \in \mathcal{P}(\mathcal{X}) :\; E_\gamma(P \| P_0) \leq \zeta \;\text{and}\; E_\gamma(P_0 \| P) \leq \zeta \right\}$. 
The use of $E_\gamma$-divergence in our formulation is particularly useful, as it provides a foundation for broader neighborhood definitions. In particular, it can be extended to neighborhoods based on general \( f \)-divergences with twice differentiable \( f \), as  
$
N_{f,\zeta}(P_0) \coloneqq \left\{ P \in \mathcal{P}(\mathcal{X}) :\; D_f(P \| P_0) \leq \zeta \;\text{and}\; D_f(P_0 \| P) \leq \zeta \right\}
$ \citep{polyanskiy2015dissipation,raginsky2016strong,zamanlooy2023strong,asoodeh2024contraction}.  
This approach provides a principled strategy for characterizing the local minimax solution over the broad and flexible class of distributional neighborhoods $N_{f,\zeta}(P_0)$. We leave this generalization for future work. 
Furthermore, this work focuses on the setting where each client releases a single sample. Extending to the case of multiple samples per client is a natural direction for future work.

% Finally, the broader impacts of our work is presented in Appendix~\ref{appendix: broader impact}. 

 \newpage

 \section*{Acknowledgments} 
 This work was supported in part by the Natural Sciences and Engineering Research Council of Canada (NSERC).
% \section{Submission of papers to NeurIPS 2025}


% Please read the instructions below carefully and follow them faithfully.


% \subsection{Style}


% Papers to be submitted to NeurIPS 2025 must be prepared according to the
% instructions presented here. Papers may only be up to {\bf nine} pages long,
% including figures.
% % Additional pages \emph{containing only acknowledgments and references} are allowed.
% Additional pages \emph{containing references, checklist, and the optional technical appendices} do not count as content pages.
% Papers that exceed the page limit will not be
% reviewed, or in any other way considered for presentation at the conference.


% The margins in 2025 are the same as those in previous years.


% Authors are required to use the NeurIPS \LaTeX{} style files obtainable at the
% NeurIPS website as indicated below. Please make sure you use the current files
% and not previous versions. Tweaking the style files may be grounds for
% rejection.


% \subsection{Retrieval of style files}


% The style files for NeurIPS and other conference information are available on
% the website at
% \begin{center}
%   \url{https://neurips.cc}
% \end{center}
% The file \verb+neurips_2025.pdf+ contains these instructions and illustrates the
% various formatting requirements your NeurIPS paper must satisfy.


% The only supported style file for NeurIPS 2025 is \verb+neurips_2025.sty+,
% rewritten for \LaTeXe{}.  \textbf{Previous style files for \LaTeX{} 2.09,
%   Microsoft Word, and RTF are no longer supported!}


% The \LaTeX{} style file contains three optional arguments: \verb+final+, which
% creates a camera-ready copy, \verb+preprint+, which creates a preprint for
% submission to, e.g., arXiv, and \verb+nonatbib+, which will not load the
% \verb+natbib+ package for you in case of package clash.


% \paragraph{Preprint option}
% If you wish to post a preprint of your work online, e.g., on arXiv, using the
% NeurIPS style, please use the \verb+preprint+ option. This will create a
% nonanonymized version of your work with the text ``Preprint. Work in progress.''
% in the footer. This version may be distributed as you see fit, as long as you do not say which conference it was submitted to. Please \textbf{do
%   not} use the \verb+final+ option, which should \textbf{only} be used for
% papers accepted to NeurIPS.


% At submission time, please omit the \verb+final+ and \verb+preprint+
% options. This will anonymize your submission and add line numbers to aid
% review. Please do \emph{not} refer to these line numbers in your paper as they
% will be removed during generation of camera-ready copies.


% The file \verb+neurips_2025.tex+ may be used as a ``shell'' for writing your
% paper. All you have to do is replace the author, title, abstract, and text of
% the paper with your own.


% The formatting instructions contained in these style files are summarized in
% Sections \ref{gen_inst}, \ref{headings}, and \ref{others} below.


% \section{General formatting instructions}
% \label{gen_inst}


% The text must be confined within a rectangle 5.5~inches (33~picas) wide and
% 9~inches (54~picas) long. The left margin is 1.5~inch (9~picas).  Use 10~point
% type with a vertical spacing (leading) of 11~points.  Times New Roman is the
% preferred typeface throughout, and will be selected for you by default.
% Paragraphs are separated by \nicefrac{1}{2}~line space (5.5 points), with no
% indentation.


% The paper title should be 17~point, initial caps/lower case, bold, centered
% between two horizontal rules. The top rule should be 4~points thick and the
% bottom rule should be 1~point thick. Allow \nicefrac{1}{4}~inch space above and
% below the title to rules. All pages should start at 1~inch (6~picas) from the
% top of the page.


% For the final version, authors' names are set in boldface, and each name is
% centered above the corresponding address. The lead author's name is to be listed
% first (left-most), and the co-authors' names (if different address) are set to
% follow. If there is only one co-author, list both author and co-author side by
% side.


% Please pay special attention to the instructions in Section \ref{others}
% regarding figures, tables, acknowledgments, and references.


% \section{Headings: first level}
% \label{headings}


% All headings should be lower case (except for first word and proper nouns),
% flush left, and bold.


% First-level headings should be in 12-point type.


% \subsection{Headings: second level}


% Second-level headings should be in 10-point type.


% \subsubsection{Headings: third level}


% Third-level headings should be in 10-point type.


% \paragraph{Paragraphs}


% There is also a \verb+\paragraph+ command available, which sets the heading in
% bold, flush left, and inline with the text, with the heading followed by 1\,em
% of space.


% \section{Citations, figures, tables, references}
% \label{others}


% These instructions apply to everyone.


% \subsection{Citations within the text}


% The \verb+natbib+ package will be loaded for you by default.  Citations may be
% author/year or numeric, as long as you maintain internal consistency.  As to the
% format of the references themselves, any style is acceptable as long as it is
% used consistently.


% The documentation for \verb+natbib+ may be found at
% \begin{center}
%   \url{http://mirrors.ctan.org/macros/latex/contrib/natbib/natnotes.pdf}
% \end{center}
% Of note is the command \verb+\citet+, which produces citations appropriate for
% use in inline text.  For example,
% \begin{verbatim}
%    \citet{hasselmo} investigated\dots
% \end{verbatim}
% produces
% \begin{quote}
%   Hasselmo, et al.\ (1995) investigated\dots
% \end{quote}


% If you wish to load the \verb+natbib+ package with options, you may add the
% following before loading the \verb+neurips_2025+ package:
% \begin{verbatim}
%    \PassOptionsToPackage{options}{natbib}
% \end{verbatim}


% If \verb+natbib+ clashes with another package you load, you can add the optional
% argument \verb+nonatbib+ when loading the style file:
% \begin{verbatim}
%    \usepackage[nonatbib]{neurips_2025}
% \end{verbatim}


% As submission is double blind, refer to your own published work in the third
% person. That is, use ``In the previous work of Jones et al.\ [4],'' not ``In our
% previous work [4].'' If you cite your other papers that are not widely available
% (e.g., a journal paper under review), use anonymous author names in the
% citation, e.g., an author of the form ``A.\ Anonymous'' and include a copy of the anonymized paper in the supplementary material.


% \subsection{Footnotes}


% Footnotes should be used sparingly.  If you do require a footnote, indicate
% footnotes with a number\footnote{Sample of the first footnote.} in the
% text. Place the footnotes at the bottom of the page on which they appear.
% Precede the footnote with a horizontal rule of 2~inches (12~picas).


% Note that footnotes are properly typeset \emph{after} punctuation
% marks.\footnote{As in this example.}


% \subsection{Figures}


% \begin{figure}
%   \centering
%   \fbox{\rule[-.5cm]{0cm}{4cm} \rule[-.5cm]{4cm}{0cm}}
%   \caption{Sample figure caption.}
% \end{figure}


% All artwork must be neat, clean, and legible. Lines should be dark enough for
% purposes of reproduction. The figure number and caption always appear after the
% figure. Place one line space before the figure caption and one line space after
% the figure. The figure caption should be lower case (except for first word and
% proper nouns); figures are numbered consecutively.


% You may use color figures.  However, it is best for the figure captions and the
% paper body to be legible if the paper is printed in either black/white or in
% color.


% \subsection{Tables}


% All tables must be centered, neat, clean and legible.  The table number and
% title always appear before the table.  See Table~\ref{sample-table}.


% Place one line space before the table title, one line space after the
% table title, and one line space after the table. The table title must
% be lower case (except for first word and proper nouns); tables are
% numbered consecutively.


% Note that publication-quality tables \emph{do not contain vertical rules.} We
% strongly suggest the use of the \verb+booktabs+ package, which allows for
% typesetting high-quality, professional tables:
% \begin{center}
%   \url{https://www.ctan.org/pkg/booktabs}
% \end{center}
% This package was used to typeset Table~\ref{sample-table}.


% \begin{table}
%   \caption{Sample table title}
%   \label{sample-table}
%   \centering
%   \begin{tabular}{lll}
%     \toprule
%     \multicolumn{2}{c}{Part}                   \\
%     \cmidrule(r){1-2}
%     Name     & Description     & Size ($\mu$m) \\
%     \midrule
%     Dendrite & Input terminal  & $\sim$100     \\
%     Axon     & Output terminal & $\sim$10      \\
%     Soma     & Cell body       & up to $10^6$  \\
%     \bottomrule
%   \end{tabular}
% \end{table}

% \subsection{Math}
% Note that display math in bare TeX commands will not create correct line numbers for submission. Please use LaTeX (or AMSTeX) commands for unnumbered display math. (You really shouldn't be using \$\$ anyway; see \url{https://tex.stackexchange.com/questions/503/why-is-preferable-to} and \url{https://tex.stackexchange.com/questions/40492/what-are-the-differences-between-align-equation-and-displaymath} for more information.)

% \subsection{Final instructions}

% Do not change any aspects of the formatting parameters in the style files.  In
% particular, do not modify the width or length of the rectangle the text should
% fit into, and do not change font sizes (except perhaps in the
% \textbf{References} section; see below). Please note that pages should be
% numbered.


% \section{Preparing PDF files}


% Please prepare submission files with paper size ``US Letter,'' and not, for
% example, ``A4.''


% Fonts were the main cause of problems in the past years. Your PDF file must only
% contain Type 1 or Embedded TrueType fonts. Here are a few instructions to
% achieve this.


% \begin{itemize}


% \item You should directly generate PDF files using \verb+pdflatex+.


% \item You can check which fonts a PDF files uses.  In Acrobat Reader, select the
%   menu Files$>$Document Properties$>$Fonts and select Show All Fonts. You can
%   also use the program \verb+pdffonts+ which comes with \verb+xpdf+ and is
%   available out-of-the-box on most Linux machines.


% \item \verb+xfig+ "patterned" shapes are implemented with bitmap fonts.  Use
%   "solid" shapes instead.


% \item The \verb+\bbold+ package almost always uses bitmap fonts.  You should use
%   the equivalent AMS Fonts:
% \begin{verbatim}
%    \usepackage{amsfonts}
% \end{verbatim}
% followed by, e.g., \verb+\mathbb{R}+, \verb+\mathbb{N}+, or \verb+\mathbb{C}+
% for $\mathbb{R}$, $\mathbb{N}$ or $\mathbb{C}$.  You can also use the following
% workaround for reals, natural and complex:
% \begin{verbatim}
%    \newcommand{\RR}{I\!\!R} %real numbers
%    \newcommand{\Nat}{I\!\!N} %natural numbers
%    \newcommand{\CC}{I\!\!\!\!C} %complex numbers
% \end{verbatim}
% Note that \verb+amsfonts+ is automatically loaded by the \verb+amssymb+ package.


% \end{itemize}


% If your file contains type 3 fonts or non embedded TrueType fonts, we will ask
% you to fix it.


% \subsection{Margins in \LaTeX{}}


% Most of the margin problems come from figures positioned by hand using
% \verb+\special+ or other commands. We suggest using the command
% \verb+\includegraphics+ from the \verb+graphicx+ package. Always specify the
% figure width as a multiple of the line width as in the example below:
% \begin{verbatim}
%    \usepackage[pdftex]{graphicx} ...
%    \includegraphics[width=0.8\linewidth]{myfile.pdf}
% \end{verbatim}
% See Section 4.4 in the graphics bundle documentation
% (\url{http://mirrors.ctan.org/macros/latex/required/graphics/grfguide.pdf})


% A number of width problems arise when \LaTeX{} cannot properly hyphenate a
% line. Please give LaTeX hyphenation hints using the \verb+\-+ command when
% necessary.

% \begin{ack}
% Use unnumbered first level headings for the acknowledgments. All acknowledgments
% go at the end of the paper before the list of references. Moreover, you are required to declare
% funding (financial activities supporting the submitted work) and competing interests (related financial activities outside the submitted work).
% More information about this disclosure can be found at: \url{https://neurips.cc/Conferences/2025/PaperInformation/FundingDisclosure}.


% Do {\bf not} include this section in the anonymized submission, only in the final paper. You can use the \texttt{ack} environment provided in the style file to automatically hide this section in the anonymized submission.
% \end{ack}

% \section*{References}

% \bibliographystyle{unsrtnat}
\bibliography{reference}

% References follow the acknowledgments in the camera-ready paper. Use unnumbered first-level heading for
% the references. Any choice of citation style is acceptable as long as you are
% consistent. It is permissible to reduce the font size to \verb+small+ (9 point)
% when listing the references.
% Note that the Reference section does not count towards the page limit.
% \medskip


% {
% \small


% [1] Alexander, J.A.\ \& Mozer, M.C.\ (1995) Template-based algorithms for
% connectionist rule extraction. In G.\ Tesauro, D.S.\ Touretzky and T.K.\ Leen
% (eds.), {\it Advances in Neural Information Processing Systems 7},
% pp.\ 609--616. Cambridge, MA: MIT Press.


% [2] Bower, J.M.\ \& Beeman, D.\ (1995) {\it The Book of GENESIS: Exploring
%   Realistic Neural Models with the GEneral NEural SImulation System.}  New York:
% TELOS/Springer--Verlag.


% [3] Hasselmo, M.E., Schnell, E.\ \& Barkai, E.\ (1995) Dynamics of learning and
% recall at excitatory recurrent synapses and cholinergic modulation in rat
% hippocampal region CA3. {\it Journal of Neuroscience} {\bf 15}(7):5249-5262.
% }


%%%%%%%%%%%%%%%%%%%%%%%%%%%%%%%%%%%%%%%%%%%%%%%%%%%%%%%%%%%%
\newpage




% \section*{NeurIPS Paper Checklist}

% %%% BEGIN INSTRUCTIONS %%%
% % The checklist is designed to encourage best practices for responsible machine learning research, addressing issues of reproducibility, transparency, research ethics, and societal impact. Do not remove the checklist: {\bf The papers not including the checklist will be desk rejected.} The checklist should follow the references and follow the (optional) supplemental material.  The checklist does NOT count towards the page
% % limit. 

% % Please read the checklist guidelines carefully for information on how to answer these questions. For each question in the checklist:
% % \begin{itemize}
% %     \item You should answer \answerYes{}, \answerNo{}, or \answerNA{}.
% %     \item \answerNA{} means either that the question is Not Applicable for that particular paper or the relevant information is Not Available.
% %     \item Please provide a short (1–2 sentence) justification right after your answer (even for NA). 
% %    % \item {\bf The papers not including the checklist will be desk rejected.}
% % \end{itemize}

% % {\bf The checklist answers are an integral part of your paper submission.} They are visible to the reviewers, area chairs, senior area chairs, and ethics reviewers. You will be asked to also include it (after eventual revisions) with the final version of your paper, and its final version will be published with the paper.

% % The reviewers of your paper will be asked to use the checklist as one of the factors in their evaluation. While "\answerYes{}" is generally preferable to "\answerNo{}", it is perfectly acceptable to answer "\answerNo{}" provided a proper justification is given (e.g., "error bars are not reported because it would be too computationally expensive" or "we were unable to find the license for the dataset we used"). In general, answering "\answerNo{}" or "\answerNA{}" is not grounds for rejection. While the questions are phrased in a binary way, we acknowledge that the true answer is often more nuanced, so please just use your best judgment and write a justification to elaborate. All supporting evidence can appear either in the main paper or the supplemental material, provided in appendix. If you answer \answerYes{} to a question, in the justification please point to the section(s) where related material for the question can be found.

% % IMPORTANT, please:
% % \begin{itemize}
% %     \item {\bf Delete this instruction block, but keep the section heading ``NeurIPS Paper Checklist"},
% %     \item  {\bf Keep the checklist subsection headings, questions/answers and guidelines below.}
% %     \item {\bf Do not modify the questions and only use the provided macros for your answers}.
% % \end{itemize} 
 

% %%% END INSTRUCTIONS %%%


% \begin{enumerate}

% \item {\bf Claims}
%     \item[] Question: Do the main claims made in the abstract and introduction accurately reflect the paper's contributions and scope?
%     \item[] Answer: \answerYes{} % Replace by \answerYes{}, \answerNo{}, or \answerNA{}.
%     \item[] Justification: The main contributions and scope of our paper are clearly outlined in the abstract and in introduction (Section~\ref{sec: introduction}).
%     \item[] Guidelines:
%     \begin{itemize}
%         \item The answer NA means that the abstract and introduction do not include the claims made in the paper.
%         \item The abstract and/or introduction should clearly state the claims made, including the contributions made in the paper and important assumptions and limitations. A No or NA answer to this question will not be perceived well by the reviewers. 
%         \item The claims made should match theoretical and experimental results, and reflect how much the results can be expected to generalize to other settings. 
%         \item It is fine to include aspirational goals as motivation as long as it is clear that these goals are not attained by the paper. 
%     \end{itemize}

% \item {\bf Limitations}
%     \item[] Question: Does the paper discuss the limitations of the work performed by the authors?
%     \item[] Answer: \answerYes{} % Replace by \answerYes{}, \answerNo{}, or \answerNA{}.
%     \item[] Justification: The limitations of our work is stated in Section~\ref{sec: conclusion}. 
%     \item[] Guidelines:
%     \begin{itemize}
%         \item The answer NA means that the paper has no limitation while the answer No means that the paper has limitations, but those are not discussed in the paper. 
%         \item The authors are encouraged to create a separate "Limitations" section in their paper.
%         \item The paper should point out any strong assumptions and how robust the results are to violations of these assumptions (e.g., independence assumptions, noiseless settings, model well-specification, asymptotic approximations only holding locally). The authors should reflect on how these assumptions might be violated in practice and what the implications would be.
%         \item The authors should reflect on the scope of the claims made, e.g., if the approach was only tested on a few datasets or with a few runs. In general, empirical results often depend on implicit assumptions, which should be articulated.
%         \item The authors should reflect on the factors that influence the performance of the approach. For example, a facial recognition algorithm may perform poorly when image resolution is low or images are taken in low lighting. Or a speech-to-text system might not be used reliably to provide closed captions for online lectures because it fails to handle technical jargon.
%         \item The authors should discuss the computational efficiency of the proposed algorithms and how they scale with dataset size.
%         \item If applicable, the authors should discuss possible limitations of their approach to address problems of privacy and fairness.
%         \item While the authors might fear that complete honesty about limitations might be used by reviewers as grounds for rejection, a worse outcome might be that reviewers discover limitations that aren't acknowledged in the paper. The authors should use their best judgment and recognize that individual actions in favor of transparency play an important role in developing norms that preserve the integrity of the community. Reviewers will be specifically instructed to not penalize honesty concerning limitations.
%     \end{itemize}

% \item {\bf Theory assumptions and proofs}
%     \item[] Question: For each theoretical result, does the paper provide the full set of assumptions and a complete (and correct) proof?
%     \item[] Answer:  \answerYes{} % Replace by \answerYes{}, \answerNo{}, or \answerNA{}.
%     \item[] Justification: All theoretical results in the paper, including theorems, propositions, and lemmas  explicitly state the necessary assumptions and are supported by either detailed proofs in the appendices or appropriate references.
%     \item[] Guidelines:
%     \begin{itemize}
%         \item The answer NA means that the paper does not include theoretical results. 
%         \item All the theorems, formulas, and proofs in the paper should be numbered and cross-referenced.
%         \item All assumptions should be clearly stated or referenced in the statement of any theorems.
%         \item The proofs can either appear in the main paper or the supplemental material, but if they appear in the supplemental material, the authors are encouraged to provide a short proof sketch to provide intuition. 
%         \item Inversely, any informal proof provided in the core of the paper should be complemented by formal proofs provided in appendix or supplemental material.
%         \item Theorems and Lemmas that the proof relies upon should be properly referenced. 
%     \end{itemize}

%     \item {\bf Experimental result reproducibility}
%     \item[] Question: Does the paper fully disclose all the information needed to reproduce the main experimental results of the paper to the extent that it affects the main claims and/or conclusions of the paper (regardless of whether the code and data are provided or not)?
%     \item[] Answer: \answerYes{} % Replace by \answerYes{}, \answerNo{}, or \answerNA{}.
%     \item[] Justification: The detailed explanation of all of the experiments is provide in corresponding section or appendix. Moreover, the instructions for reproducing all figures and experimental results are provided in Appendix~\ref{appendix: reproduce}.
%     \item[] Guidelines:
%     \begin{itemize}
%         \item The answer NA means that the paper does not include experiments.
%         \item If the paper includes experiments, a No answer to this question will not be perceived well by the reviewers: Making the paper reproducible is important, regardless of whether the code and data are provided or not.
%         \item If the contribution is a dataset and/or model, the authors should describe the steps taken to make their results reproducible or verifiable. 
%         \item Depending on the contribution, reproducibility can be accomplished in various ways. For example, if the contribution is a novel architecture, describing the architecture fully might suffice, or if the contribution is a specific model and empirical evaluation, it may be necessary to either make it possible for others to replicate the model with the same dataset, or provide access to the model. In general. releasing code and data is often one good way to accomplish this, but reproducibility can also be provided via detailed instructions for how to replicate the results, access to a hosted model (e.g., in the case of a large language model), releasing of a model checkpoint, or other means that are appropriate to the research performed.
%         \item While NeurIPS does not require releasing code, the conference does require all submissions to provide some reasonable avenue for reproducibility, which may depend on the nature of the contribution. For example
%         \begin{enumerate}
%             \item If the contribution is primarily a new algorithm, the paper should make it clear how to reproduce that algorithm.
%             \item If the contribution is primarily a new model architecture, the paper should describe the architecture clearly and fully.
%             \item If the contribution is a new model (e.g., a large language model), then there should either be a way to access this model for reproducing the results or a way to reproduce the model (e.g., with an open-source dataset or instructions for how to construct the dataset).
%             \item We recognize that reproducibility may be tricky in some cases, in which case authors are welcome to describe the particular way they provide for reproducibility. In the case of closed-source models, it may be that access to the model is limited in some way (e.g., to registered users), but it should be possible for other researchers to have some path to reproducing or verifying the results.
%         \end{enumerate}
%     \end{itemize}


% \item {\bf Open access to data and code}
%     \item[] Question: Does the paper provide open access to the data and code, with sufficient instructions to faithfully reproduce the main experimental results, as described in supplemental material?
%     \item[] Answer: \answerYes{} % Replace by \answerYes{}, \answerNo{}, or \answerNA{}.
%     \item[] Justification: We will attach an anonymized version of the complete code for generating all experimental results as part of the supplementary material at the time of submission, in accordance with the supplementary material deadline. The supplementary material will be made publicly available after the review period. Additionally, instructions for reproducing all figures and experimental results are provided in Appendix~\ref{appendix: reproduce}.
%     \item[] Guidelines:
%     \begin{itemize}
%         \item The answer NA means that paper does not include experiments requiring code.
%         \item Please see the NeurIPS code and data submission guidelines (\url{https://nips.cc/public/guides/CodeSubmissionPolicy}) for more details.
%         \item While we encourage the release of code and data, we understand that this might not be possible, so “No” is an acceptable answer. Papers cannot be rejected simply for not including code, unless this is central to the contribution (e.g., for a new open-source benchmark).
%         \item The instructions should contain the exact command and environment needed to run to reproduce the results. See the NeurIPS code and data submission guidelines (\url{https://nips.cc/public/guides/CodeSubmissionPolicy}) for more details.
%         \item The authors should provide instructions on data access and preparation, including how to access the raw data, preprocessed data, intermediate data, and generated data, etc.
%         \item The authors should provide scripts to reproduce all experimental results for the new proposed method and baselines. If only a subset of experiments are reproducible, they should state which ones are omitted from the script and why.
%         \item At submission time, to preserve anonymity, the authors should release anonymized versions (if applicable).
%         \item Providing as much information as possible in supplemental material (appended to the paper) is recommended, but including URLs to data and code is permitted.
%     \end{itemize}


% \item {\bf Experimental setting/details}
%     \item[] Question: Does the paper specify all the training and test details (e.g., data splits, hyperparameters, how they were chosen, type of optimizer, etc.) necessary to understand the results?
%     \item[] Answer: \answerYes{} % Replace by \answerYes{}, \answerNo{}, or \answerNA{}.
%     \item[] Justification: We provide sufficient details to implement our proposed sampler and include the code for both the proposed method and the baseline in the supplementary material. 
%     \item[] Guidelines:
%     \begin{itemize}
%         \item The answer NA means that the paper does not include experiments.
%         \item The experimental setting should be presented in the core of the paper to a level of detail that is necessary to appreciate the results and make sense of them.
%         \item The full details can be provided either with the code, in appendix, or as supplemental material.
%     \end{itemize}

% \item {\bf Experiment statistical significance}
%     \item[] Question: Does the paper report error bars suitably and correctly defined or other appropriate information about the statistical significance of the experiments?
%     \item[] Answer: \answerNo{} % Replace by \answerYes{}, \answerNo{}, or \answerNA{}.
%     \item[] Justification: Since the main contribution lies in theoretically characterizing the \emph{worst-case} utility, our experiments are designed to extract and evaluate this worst-case behavior, rather than to assess average-case performance or variability. 
%     \item[] Guidelines:
%     \begin{itemize}
%         \item The answer NA means that the paper does not include experiments.
%         \item The authors should answer "Yes" if the results are accompanied by error bars, confidence intervals, or statistical significance tests, at least for the experiments that support the main claims of the paper.
%         \item The factors of variability that the error bars are capturing should be clearly stated (for example, train/test split, initialization, random drawing of some parameter, or overall run with given experimental conditions).
%         \item The method for calculating the error bars should be explained (closed form formula, call to a library function, bootstrap, etc.)
%         \item The assumptions made should be given (e.g., Normally distributed errors).
%         \item It should be clear whether the error bar is the standard deviation or the standard error of the mean.
%         \item It is OK to report 1-sigma error bars, but one should state it. The authors should preferably report a 2-sigma error bar than state that they have a 96\% CI, if the hypothesis of Normality of errors is not verified.
%         \item For asymmetric distributions, the authors should be careful not to show in tables or figures symmetric error bars that would yield results that are out of range (e.g. negative error rates).
%         \item If error bars are reported in tables or plots, The authors should explain in the text how they were calculated and reference the corresponding figures or tables in the text.
%     \end{itemize}

% \item {\bf Experiments compute resources}
%     \item[] Question: For each experiment, does the paper provide sufficient information on the computer resources (type of compute workers, memory, time of execution) needed to reproduce the experiments?
%     \item[] Answer: \answerYes{} % Replace by \answerYes{}, \answerNo{}, or \answerNA{}.
%     \item[] Justification: Details on the computational resources and running times for the experiments are provided in Appendix~\ref{appendix: reproduce}.
%     \item[] Guidelines:
%     \begin{itemize}
%         \item The answer NA means that the paper does not include experiments.
%         \item The paper should indicate the type of compute workers CPU or GPU, internal cluster, or cloud provider, including relevant memory and storage.
%         \item The paper should provide the amount of compute required for each of the individual experimental runs as well as estimate the total compute. 
%         \item The paper should disclose whether the full research project required more compute than the experiments reported in the paper (e.g., preliminary or failed experiments that didn't make it into the paper). 
%     \end{itemize}
    
% \item {\bf Code of ethics}
%     \item[] Question: Does the research conducted in the paper conform, in every respect, with the NeurIPS Code of Ethics \url{https://neurips.cc/public/EthicsGuidelines}?
%     \item[] Answer: \answerYes{} % Replace by \answerYes{}, \answerNo{}, or \answerNA{}.
%     \item[] Justification: Our research conforms with the NeurIPS Code of Ethics.
%     \item[] Guidelines:
%     \begin{itemize}
%         \item The answer NA means that the authors have not reviewed the NeurIPS Code of Ethics.
%         \item If the authors answer No, they should explain the special circumstances that require a deviation from the Code of Ethics.
%         \item The authors should make sure to preserve anonymity (e.g., if there is a special consideration due to laws or regulations in their jurisdiction).
%     \end{itemize}


% \item {\bf Broader impacts}
%     \item[] Question: Does the paper discuss both potential positive societal impacts and negative societal impacts of the work performed?
%     \item[] Answer: \answerYes{} % Replace by \answerYes{}, \answerNo{}, or \answerNA{}.
%     \item[] Justification: We discussed both the positive and negative societal impacts of out work in Appendix \ref{appendix: broader impact}.
%     \item[] Guidelines:
%     \begin{itemize}
%         \item The answer NA means that there is no societal impact of the work performed.
%         \item If the authors answer NA or No, they should explain why their work has no societal impact or why the paper does not address societal impact.
%         \item Examples of negative societal impacts include potential malicious or unintended uses (e.g., disinformation, generating fake profiles, surveillance), fairness considerations (e.g., deployment of technologies that could make decisions that unfairly impact specific groups), privacy considerations, and security considerations.
%         \item The conference expects that many papers will be foundational research and not tied to particular applications, let alone deployments. However, if there is a direct path to any negative applications, the authors should point it out. For example, it is legitimate to point out that an improvement in the quality of generative models could be used to generate deepfakes for disinformation. On the other hand, it is not needed to point out that a generic algorithm for optimizing neural networks could enable people to train models that generate Deepfakes faster.
%         \item The authors should consider possible harms that could arise when the technology is being used as intended and functioning correctly, harms that could arise when the technology is being used as intended but gives incorrect results, and harms following from (intentional or unintentional) misuse of the technology.
%         \item If there are negative societal impacts, the authors could also discuss possible mitigation strategies (e.g., gated release of models, providing defenses in addition to attacks, mechanisms for monitoring misuse, mechanisms to monitor how a system learns from feedback over time, improving the efficiency and accessibility of ML).
%     \end{itemize}
    
% \item {\bf Safeguards}
%     \item[] Question: Does the paper describe safeguards that have been put in place for responsible release of data or models that have a high risk for misuse (e.g., pretrained language models, image generators, or scraped datasets)?
%     \item[] Answer: \answerNA{} % Replace by \answerYes{}, \answerNo{}, or \answerNA{}.
%     \item[] Justification: We do not identify any foreseeable risk of misuse.
%     \item[] Guidelines:
%     \begin{itemize}
%         \item The answer NA means that the paper poses no such risks.
%         \item Released models that have a high risk for misuse or dual-use should be released with necessary safeguards to allow for controlled use of the model, for example by requiring that users adhere to usage guidelines or restrictions to access the model or implementing safety filters. 
%         \item Datasets that have been scraped from the Internet could pose safety risks. The authors should describe how they avoided releasing unsafe images.
%         \item We recognize that providing effective safeguards is challenging, and many papers do not require this, but we encourage authors to take this into account and make a best faith effort.
%     \end{itemize}

% \item {\bf Licenses for existing assets}
%     \item[] Question: Are the creators or original owners of assets (e.g., code, data, models), used in the paper, properly credited and are the license and terms of use explicitly mentioned and properly respected?
%     \item[] Answer: \answerYes{} % Replace by \answerYes{}, \answerNo{}, or \answerNA{}.
%     \item[] Justification: We cite the original paper associated with the source code and ensure that its license and terms of use are properly respected.
%     \item[] Guidelines:
%     \begin{itemize}
%         \item The answer NA means that the paper does not use existing assets.
%         \item The authors should cite the original paper that produced the code package or dataset.
%         \item The authors should state which version of the asset is used and, if possible, include a URL.
%         \item The name of the license (e.g., CC-BY 4.0) should be included for each asset.
%         \item For scraped data from a particular source (e.g., website), the copyright and terms of service of that source should be provided.
%         \item If assets are released, the license, copyright information, and terms of use in the package should be provided. For popular datasets, \url{paperswithcode.com/datasets} has curated licenses for some datasets. Their licensing guide can help determine the license of a dataset.
%         \item For existing datasets that are re-packaged, both the original license and the license of the derived asset (if it has changed) should be provided.
%         \item If this information is not available online, the authors are encouraged to reach out to the asset's creators.
%     \end{itemize}

% \item {\bf New assets}
%     \item[] Question: Are new assets introduced in the paper well documented and is the documentation provided alongside the assets?
%     \item[] Answer: \answerNA{} % Replace by \answerYes{}, \answerNo{}, or \answerNA{}.
%     \item[] Justification: The paper does not release new assets.
%     \item[] Guidelines:
%     \begin{itemize}
%         \item The answer NA means that the paper does not release new assets.
%         \item Researchers should communicate the details of the dataset/code/model as part of their submissions via structured templates. This includes details about training, license, limitations, etc. 
%         \item The paper should discuss whether and how consent was obtained from people whose asset is used.
%         \item At submission time, remember to anonymize your assets (if applicable). You can either create an anonymized URL or include an anonymized zip file.
%     \end{itemize}

% \item {\bf Crowdsourcing and research with human subjects}
%     \item[] Question: For crowdsourcing experiments and research with human subjects, does the paper include the full text of instructions given to participants and screenshots, if applicable, as well as details about compensation (if any)? 
%     \item[] Answer: \answerNA{} % Replace by \answerYes{}, \answerNo{}, or \answerNA{}.
%     \item[] Justification: Our paper does not involve crowdsourcing nor research with human subjects.
%     \item[] Guidelines:
%     \begin{itemize}
%         \item The answer NA means that the paper does not involve crowdsourcing nor research with human subjects.
%         \item Including this information in the supplemental material is fine, but if the main contribution of the paper involves human subjects, then as much detail as possible should be included in the main paper. 
%         \item According to the NeurIPS Code of Ethics, workers involved in data collection, curation, or other labor should be paid at least the minimum wage in the country of the data collector. 
%     \end{itemize}

% \item {\bf Institutional review board (IRB) approvals or equivalent for research with human subjects}
%     \item[] Question: Does the paper describe potential risks incurred by study participants, whether such risks were disclosed to the subjects, and whether Institutional Review Board (IRB) approvals (or an equivalent approval/review based on the requirements of your country or institution) were obtained?
%     \item[] Answer: \answerNA{} % Replace by \answerYes{}, \answerNo{}, or \answerNA{}.
%     \item[] Justification: The paper does not involve crowdsourcing nor research with human subjects.
%     \item[] Guidelines:
%     \begin{itemize}
%         \item The answer NA means that the paper does not involve crowdsourcing nor research with human subjects.
%         \item Depending on the country in which research is conducted, IRB approval (or equivalent) may be required for any human subjects research. If you obtained IRB approval, you should clearly state this in the paper. 
%         \item We recognize that the procedures for this may vary significantly between institutions and locations, and we expect authors to adhere to the NeurIPS Code of Ethics and the guidelines for their institution. 
%         \item For initial submissions, do not include any information that would break anonymity (if applicable), such as the institution conducting the review.
%     \end{itemize}

% \item {\bf Declaration of LLM usage}
%     \item[] Question: Does the paper describe the usage of LLMs if it is an important, original, or non-standard component of the core methods in this research? Note that if the LLM is used only for writing, editing, or formatting purposes and does not impact the core methodology, scientific rigorousness, or originality of the research, declaration is not required.
%     %this research? 
%     \item[] Answer: \answerNA{} % Replace by \answerYes{}, \answerNo{}, or \answerNA{}.
%     \item[] Justification: The core method development in this research does not involve LLMs as any important, original, or non-standard components.
%     \item[] Guidelines:
%     \begin{itemize}
%         \item The answer NA means that the core method development in this research does not involve LLMs as any important, original, or non-standard components.
%         \item Please refer to our LLM policy (\url{https://neurips.cc/Conferences/2025/LLM}) for what should or should not be described.
%     \end{itemize}

% \end{enumerate}

\newpage
\ifshowappendix

%%%%\input{neurIPS_2025_Hrad_suppl}
%%%%\input{neurips_2025_Hrad}
\appendix 

\textbf{Roadmap of Appendix:}
The appendix is organized as follows. A more complete notation table—complementing the abbreviated summary in the main body—is provided in Section~\ref{appendix: notation}. Supplementary definitions and theorems used in the proofs of our results appear in Section~\ref{appendix: Supplementary Definitions and Theorems}. Section~\ref{sec: proofs prilim} presents preliminary results that support the proofs of our main theorems, which are then detailed in Section~\ref{sec: proofs main}. Section~\ref{sec: appendix experimental setup and edditional experiments} describes our experimental setup and provides additional numerical results. 
% Finally, instructions for reproducing our experiments and a discussion of the broader impacts of our work are given in Sections~\ref{appendix: reproduce} and~\ref{appendix: broader impact}, respectively.
Finally, instructions for reproducing our experiments is given in Sections~\ref{appendix: reproduce}.

\section{Notation table}\label{appendix: notation}
Table~\ref{tab: notation-table} provides a more comprehensive summary of the notation used throughout the paper, complementing the abbreviated overview in the main body for the reader’s convenience.
\begin{table}[ht]
  \caption{Important notations used in the paper}
  \label{tab: notation-table}
  \centering
  \begin{tabular}{ll}
    \toprule
    Notation & Description \\
    \midrule
    $\mathcal{P}(\mathcal{X})$ & Set of all probability distributions over sample space $\mathcal{X}$ \\
    $\mathcal{C}(\mathbb{R}^n)$ & Set of all continuous probability distributions on $\mathbb{R}^n$ \\
    $p$ & Probability density function (PDF) of $P \in \mathcal{C}(\mathbb{R}^n)$ \\
    $[k]$ & Shorthand for $\{1, 2, \dots, k\}$ for $k \in \mathbb{N}$ \\
    $\mu_k$ & uniform distribution over finite sample space $[k]$ \\
    $\mathrm{clip}(x; s_1, s_2)$ & Clipping function: $\max\{s_1, \min\{s_2, x\}\}$ \\
    $g^*$ & Convex conjugate of a function $g$ \\
    $\mathcal{L}(m, b)$ & $n$-dimensional Laplace distribution with mean $m \in \mathbb{R}^n$ and scale $b > 0$ \\
    $\mathcal{L}(x \mid m, b)$ & Density of the Laplace distribution at $x \in \mathbb{R}^n$ \\
    $\mathcal{N}(m, \Sigma)$ & $n$-dimensional Gaussian with mean $m \in \mathbb{R}^n$ and covariance matrix $\Sigma$ \\

    $\mathcal{N}_{m, \Sigma}[x]$ & Density of the Gaussian distribution at $x \in \mathbb{R}^n$ \\
    $\mathcal{Q}_{\mathcal{X},\tilde{\mathcal{P}},\varepsilon}$ & Set of $\varepsilon$-LDP samplers: $\bQ: \tilde{\mathcal{P}} \to \mathcal{P}(\mathcal{X})$ \\
    $\mathcal{Q}_{\mathcal{X},\tilde{\mathcal{P}},\varepsilon,\delta}$ & Set of $(\varepsilon,\delta)$-LDP samplers: $\bQ: \tilde{\mathcal{P}} \to \mathcal{P}(\mathcal{X})$ \\
    $\mathcal{Q}_{\mathcal{X},\tilde{\mathcal{P}},g}$ & Set of $g$-FLDP samplers: $\bQ: \tilde{\mathcal{P}} \to \mathcal{P}(\mathcal{X})$ \\

    $D_f^B(\lambda_1 \,\|\, \lambda_2)$ & $f$-divergence between Bernoulli distributions with $\Pr(1) = \lambda_1$ and $\lambda_2$ \\
    \bottomrule
  \end{tabular}
\end{table}


\section{Supplementary definitions and theorems
}\label{appendix: Supplementary Definitions and Theorems}

\subsection{\texorpdfstring{General definition and special cases of $f$-divergence}{f-divergence}}\label{subsec: f-div} 

In this appendix, we review the general definition and special cases of $f$-divergences which are important in this work. 

\begin{definition}[General case of $f$-divergence \citep{sason2016f}]\label{def: general f-div}
    Let \( P \) and \( Q \) be probability measures, and let \( \mu \) be a dominating measure of \( P \) and \( Q \) (i.e., \( P, Q \ll \mu \); e.g., \( \mu = P + Q \)), and let \( p := \frac{dP}{d\mu} \) and \( q := \frac{dQ}{d\mu} \). The \textit{f}-divergence from \( P \) to \( Q \) is given, independently of \( \mu \), by
\begin{equation}
    D_f(P\|Q) := \int q f\left( \frac{p}{q} \right) d\mu,
\end{equation}
\noindent where
\begin{align}
    f(0) &:= \lim_{t \to 0^+} f(t), \\
    0 f\left( \frac{0}{0} \right) &:= 0, \\
    0 f\left( \frac{a}{0} \right) &:= \lim_{t \to 0^+} t f\left( \frac{a}{t} \right) = a \lim_{u \to \infty} \frac{f(u)}{u}, \quad a > 0.
\end{align}
\end{definition}


Popular examples of $f$-divergences include KL divergence ($f(t) = t\log t$), total variation distance ($f(t) =\frac{1}{2} |t-1|$), squared Hellinger distance ($f(t) = (\sqrt{t} - 1)^2$), and $\chi^2$-divergence ($f(t) = (t-1)^2$). We also formally define $E_\gamma$-divergence for  \( f(t) = \max\{t - \gamma, 0\} \) for \( \gamma \geq 1 \).

\begin{definition}[$E_\gamma$-divergence \citep{sharma2013fundamental}]
\label{def:Egamma-divergence}
Given $\gamma \geq 1$, for two probability measures $P,Q \in \mathcal{P}(\mathcal{X})$ with $P \ll Q$, the
\emph{$E_\gamma$-divergence} is
\[
    E_\gamma(P \,\|\, Q)
    \;:=\;
    \mathbb{E}_Q\!\left[
        \max\!\Bigl\{\tfrac{dP}{dQ}-\gamma,\,0\Bigr\}
    \right]= \frac{1}{2} \int \left| \mathrm{d}P - \gamma\,\mathrm{d}Q \right| - \frac{1}{2} (\gamma - 1).
\]
% \[
% E_\gamma(P \| Q) := \frac{1}{2} \int \left| \mathrm{d}P - \gamma\,\mathrm{d}Q \right| - \frac{1}{2} (\gamma - 1).
% \]
\end{definition}

For a more extensive list of 
$f$-divergences and their properties, we refer readers to~\citet{sason2016f}.  

\subsection{Measure-theoretic assumptions}

The appendices assume familiarity with basic measure theory and real analysis. 
Throughout the main paper and the appendices, we adopt the following conventions. For every sample space \( \mathcal{X} \), a \(\sigma\)-algebra on \( \mathcal{X} \) is assumed to be given implicitly. Unless stated otherwise, we use the discrete \(\sigma\)-algebra when \( \mathcal{X} \) is finite, and the Borel \(\sigma\)-algebra when \( \mathcal{X} = \mathbb{R}^n \). 
Whenever we refer to a “subset” of \( \mathcal{X} \), we mean a measurable subset. Similarly, the notation \( A \subseteq \mathcal{X} \) always indicates that \( A \) is measurable.
Finally, the term “continuous distribution” refers specifically to a distribution that is absolutely continuous with respect to the Lebesgue measure.


\subsection{Supplementary theorems}

This subsection presents two well-known results that we will use in proving our main results.


\begin{theorem}[Data processing inequality]\label{thm: dpi}
   Let \( M \) be a conditional distribution (Markov kernel) mapping from \( \mathcal{X} \) to \( \mathcal{Y} \). Given distributions \( P_1, P_2 \in \mathcal{P}(\mathcal{X}) \), and their push-forward measures \( Q_1, Q_2 \in \mathcal{P}(\mathcal{Y}) \) through \( M \), the following inequality holds for any \( f \)-divergence:
\[
D_f(P_1 \parallel P_2) \geq D_f(Q_1 \parallel Q_2).
\]
\end{theorem}

\begin{proposition}[Proposition 1 of \citet{dong2022gaussian}]\label{prop: trade-off}
    A function $g : [0,1] \to [0,1]$ is a trade-off function if and only if $g$ is convex, continuous, non-increasing, and $g(x) \leq 1 - x$ for $x \in [0,1]$.
\end{proposition}

\section{Preliminary results
}\label{sec: proofs prilim}

\begin{lemma}\label{lem:decrease-lambda}
Let $\mu$ be a (fixed) distribution over some output space $\mathcal{X}$. For each $0 \leq \lambda \leq 1$, define a sampler $\bQ_\lambda$ by
\[
  \bQ_\lambda(A \mid P) \;=\; \lambda\,P(A) \;+\; (1-\lambda)\,\mu(A),
  \quad
  \text{for any distribution }P \text{ on } \mathcal{X} \text{ and event } A \subseteq \mathcal{X}.
\]
Suppose there exists $0 \leq \lambda_1 \leq 1$ such that $\bQ_{\lambda_1}$ is $(\varepsilon,\delta)$-LDP. Then for any $\lambda_2 \in [0,\lambda_1]$, the sampler $\bQ_{\lambda_2}$ is also $(\varepsilon,\delta)$-LDP.
\end{lemma}

\begin{proof}
Fix $\lambda_2 \in [0,\lambda_1]$, and let $P,P'$ be any two distributions on $\mathcal{X}$. We need to show that for every event $A$,
\[
  \bQ_{\lambda_2}(A \mid P)
  \;\le\;
  e^\varepsilon\;\bQ_{\lambda_2}(A \mid P')
  \;+\;
  \delta.
\]
Because $\lambda_2 \le \lambda_1$, there is a $\beta \in [0,1]$ such that
\[
  \lambda_2 \;=\;\beta \,\lambda_1.
\]
Then, for all events $A$,
\[
  \bQ_{\lambda_2}(A \mid P)
  \;=\;
  \lambda_2\,P(A) \;+\; (1-\lambda_2)\,\mu(A)
  \;=\;
  \beta\,\lambda_1\,P(A) + \bigl[1 - \beta\,\lambda_1\bigr]\,\mu(A).
\]
Recognizing that $\lambda_1\,P(A) + (1-\lambda_1)\,\mu(A)$ is $\bQ_{\lambda_1}(A \mid P)$, we rewrite:
\[
  \bQ_{\lambda_2}(A \mid P)
  \;=\;
  \beta\,\bQ_{\lambda_1}(A \mid P)
  \;+\;
  \bigl[\,1-\beta\,\bigr]\mu(A).
\]
Similarly,
\[
  \bQ_{\lambda_2}(A \mid P')
  \;=\;
  \beta\,\bQ_{\lambda_1}(A \mid P') \;+\; \bigl[1-\beta\bigr]\mu(A).
\]

Since $\bQ_{\lambda_1}$ is $(\varepsilon,\delta)$-LDP, we know
\[
  \bQ_{\lambda_1}(A \mid P)
  \;\le\;
  e^\varepsilon\, \bQ_{\lambda_1}(A \mid P') \;+\;\delta.
\]
Multiplying both sides by $\beta \ge 0$ and then adding $(1-\beta)\,\mu(A)$, we obtain
\[
  \beta\,\bQ_{\lambda_1}(A \mid P) \;+\; (1-\beta)\,\mu(A)
  \;\le\;
  \beta\,e^\varepsilon\,\bQ_{\lambda_1}(A \mid P')
  \;+\;
  \beta\,\delta
  \;+\;
  (1-\beta)\,\mu(A).
\]
But the left-hand side above is exactly $\bQ_{\lambda_2}(A \mid P)$, so
\[
  \bQ_{\lambda_2}(A \mid P)
  \;\le\;
  \beta\,e^\varepsilon\,\bQ_{\lambda_1}(A \mid P')
  \;+\;
  \beta\,\delta
  \;+\;
  (1-\beta)\,\mu(A).
\]
On the other hand,
\[
  e^\varepsilon\,\bQ_{\lambda_2}(A \mid P')
  \;=\;
  e^\varepsilon\,\Bigl[\beta\,\bQ_{\lambda_1}(A \mid P') + (1-\beta)\,\mu(A)\Bigr]
  \;=\;
  \beta\,e^\varepsilon\,\bQ_{\lambda_1}(A \mid P') + e^\varepsilon\,(1-\beta)\,\mu(A).
\]
Hence,
\[
  e^\varepsilon\,\bQ_{\lambda_2}(A \mid P') \;+\;\delta
  \;=\;
  \beta\,e^\varepsilon\,\bQ_{\lambda_1}(A \mid P')
  \;+\;
  (1-\beta)\,e^\varepsilon\,\mu(A)
  \;+\;\delta.
\]
It is enough to check
\[
  \beta\,e^\varepsilon\,\bQ_{\lambda_1}(A \mid P')
  \;+\;
  \beta\,\delta
  \;+\;
  (1-\beta)\,\mu(A)
  \;\;\le\;\;
  \beta\,e^\varepsilon\,\bQ_{\lambda_1}(A \mid P')
  \;+\;
  (1-\beta)\,e^\varepsilon\,\mu(A)
  \;+\;
  \delta,
\]
which simplifies to
\[
  \beta\,\delta + (1-\beta)\,\mu(A)
  \;\le\;
  (1-\beta)\,e^\varepsilon\,\mu(A) + \delta.
\]
Rearranging,
\[
  (1-\beta)\,\mu(A)\bigl[1 - e^\varepsilon\bigr]
  \;\;\le\;\;
  \delta \bigl[\,1 - \beta\bigr].
\]
Since $1 - e^\varepsilon \le 0$ for $\varepsilon \ge 0$, and $\delta \ge 0$, this inequality holds trivially (the left-hand side is at most $0$, while the right-hand side is nonnegative). Thus the overall chain of inequalities is valid, and we conclude
\[
  \bQ_{\lambda_2}(A \mid P)
  \;\le\;
  e^\varepsilon\,\bQ_{\lambda_2}(A \mid P') \;+\;\delta,
  \quad
  \forall P,P',\quad \text{event } A.
\]
Hence $\bQ_{\lambda_2}$ is also $(\varepsilon,\delta)$-LDP.
\end{proof}

Now, we first restate the normalization condition and introduce the non-triviality assumption for $(\eps,\delta)$-LDP before stating the next proposition.

\textbf{Normalization condition:}
\begin{equation}\label{eqn: normalization condition}
    \mu(\mathcal{X}) = 1, \quad 0 \le c_1 < 1 < c_2.
\end{equation}
% \begin{equation}\label{eqn: continuous normalization condition}
%     \int_{\R^n}h(x)dx = 1, \quad 0 \leq c_1 < 1 < c_2.
% \end{equation}
\textbf{Non-triviality condition:}
% \begin{equation}\label{eqn: non-trivial condition condition pure}
%     c_2 > c_1 e^\eps.
% \end{equation}
\begin{equation}\label{eqn: non-trivial condition condition approximate}
    \delta \leq \frac{(c_2-c_1e^\eps)(1-c_1)}{c_2-c_1}.
\end{equation}
\medskip
\begin{proposition}\label{prop: global appx linear}
    Consider the following optimization problem:
    \begin{equation}\label{eqn: global appx}
        \mathcal{R}\big(\mathcal{Q}_{\mathcal{X}, \tilde{\mathcal{P}}, \varepsilon, \delta}, \tilde{\mathcal{P}}, f\big) 
        \coloneqq 
        \inf_{\bQ \in \mathcal{Q}_{\mathcal{X}, \tilde{\mathcal{P}}, \varepsilon, \delta}} 
        \;\sup_{P \in \tilde{\mathcal{P}}} 
        \;D_f\big(P \, \| \, \bQ(P)\big).
    \end{equation}

    Let $\tilde{\mathcal{P}} = \tilde{\mathcal{P}}_{c_1, c_2, \mu}$ such that the normalization condition~\eqref{eqn: normalization condition} and the non-triviality condition (\ref{eqn: non-trivial condition condition approximate}) hold and $\frac{c_2 - c_1}{1 - c_1} \in \mathbb{N}$. Suppose that $\mu$ is $(\alpha,\frac{1}{\alpha},1)$-decomposable with $\alpha = \frac{1 - c_1}{c_2 - c_1}$. Define
    \[
    r_1 := \frac{c_1}{c_2 - c_1} \cdot \frac{(1 - c_1)e^{\varepsilon} + c_2 - 1}{1 - \delta},
    \qquad
    r_2 := \frac{c_2}{c_2 - c_1} \cdot \frac{(1 - c_1)e^{\varepsilon} + c_2 - 1}{e^{\varepsilon} + \frac{c_2 - 1}{1 - c_1} \delta}.
    \]

    Then, the optimal value of the problem~\eqref{eqn: global appx} is given by
    \begin{equation}\label{eqn: global appx optimal value}
        \mathcal{R}\big(\mathcal{Q}_{\mathcal{X}, \tilde{\mathcal{P}}, \varepsilon, \delta}, \tilde{\mathcal{P}}, f\big)
        = \frac{1 - r_1}{r_2 - r_1} \, f(r_2) + \frac{r_2 - 1}{r_2 - r_1} \, f(r_1).
    \end{equation}

    Moreover, the sampler
    $\bQ^\star_{\varepsilon, \delta}(P) \coloneqq \lambda^\star_{\varepsilon, \delta} P + (1 - \lambda^\star_{\varepsilon, \delta}) \mu,
    $
    with
    $
    \lambda^\star_{\varepsilon, \delta} 
    = \frac{e^{\varepsilon} + \frac{c_2 - c_1}{1 - c_1} \delta - 1}
           {(1 - c_1) e^{\varepsilon} + c_2 - 1}$,
    belongs to the class $\mathcal{Q}_{\mathcal{X}, \tilde{\mathcal{P}}, \varepsilon, \delta}$ and is minimax-optimal under any $f$-divergence $D_f$. That is,
    \begin{equation}\label{eqn: global appx linear is optimal}
        \sup_{P \in \tilde{\mathcal{P}}} D_f\big(P \, \| \, \bQ^\star_{\varepsilon, \delta}(P)\big)
        = \mathcal{R}\big(\mathcal{Q}_{\mathcal{X}, \tilde{\mathcal{P}}, \varepsilon, \delta}, \tilde{\mathcal{P}}, f\big).
    \end{equation}
\end{proposition}
\begin{proof}
We first show that if the non-triviality constraint (\ref{eqn: non-trivial condition condition approximate}) does not hold, then the identity sampler $\bQ^I(P) = P$ satisfies $(\eps,\delta)$-LDP and has the trivial minimax risk zero.  

Suppose we have $\delta>\frac{(c_2-c_1e^\eps)(1-c_1)}{c_2-c_1}$. Let $A\subseteq\mathcal{X}$ be a measurable subset and $P,P^\prime\in\mathcal{P}_{c_1,c_2,\mu}$. If $0\leq\mu(A)\leq\frac{1-c_1}{c_2-c_1}$,
\begin{eqnarray*}
    e^\eps P(A)+\delta-P^\prime(A)&\geq&c_1e^\eps\mu(A)+\delta-c_2\mu(A)\\
    &\geq&\min\left\{\frac{(c_1e^\eps-c_2)(1-c_1)}{c_2-c_1}+\delta,\delta\right\}\\
    &\geq&0.
\end{eqnarray*}
If $\mu(A)\geq\frac{1-c_1}{c_2-c_1}$,
\begin{eqnarray*}
    e^\eps P(A)+\delta-P^\prime(A)&\geq&c_1e^\eps\mu(A)+\delta-(1-P^\prime(A^c))\\
    &\geq&c_1e^\eps\mu(A)+\delta-(1-c_1(1-\mu(A)))\\
    &=&(c_1e^\eps-c_1)\mu(A)+\delta-1+c_1\\
    &\geq&\frac{(c_1e^\eps-c_1)(1-c_1)}{c_2-c_1}-1+c_1+\delta\\
    &=&\frac{(c_1e^\eps-c_2)(1-c_1)}{c_2-c_1}+\delta\\
    &>&0.
\end{eqnarray*}

Therefore, for $\delta>\frac{(c_2-c_1e^\eps)(1-c_1)}{c_2-c_1}$, $\bQ^I(P) \in \mathcal{Q}_{\mathcal{X}, \tilde{\mathcal{P}}, \varepsilon, \delta}$ and the minimax risk is trivially zero. In the remainder of the proof, we consider the non-trivial case of $\delta \leq \frac{(c_2-c_1e^\eps)(1-c_1)}{c_2-c_1}$.

   \medskip
   
    \textbf{Step 1: Proof of $\bQ^\star_{\varepsilon, \delta}$ is a probability measure} 

    In this proof, we show that $\bQ^\star_{\varepsilon, \delta}$  defined by $
\bQ^\star_{\varepsilon, \delta}(P) 
:= \lambda^\star_{\varepsilon, \delta} P + \bigl(1 - \lambda^\star_{\varepsilon, \delta}\bigr) \mu
$ is a probability measure.
Since $P$ and $\mu$ are probability measures, we have
\[
\bQ^\star_{\varepsilon, \delta}(\mathcal{X})
= \lambda^\star_{\varepsilon, \delta} P(\mathcal{X}) 
+ \bigl(1 - \lambda^\star_{\varepsilon, \delta}\bigr)\mu(\mathcal{X})
= \lambda^\star_{\varepsilon, \delta} + \bigl(1 - \lambda^\star_{\varepsilon, \delta}\bigr)
= 1.
\]
Assumptions $0\leq c_1 < 1 < c_2$ and $\delta \leq \frac{(c_2-c_1e^\eps)(1-c_1)}{c_2-c_1}$ imply
\[\lambda^\star_{\varepsilon, \delta} 
    = \frac{e^{\varepsilon} + \frac{c_2 - c_1}{1 - c_1} \delta - 1}
           {(1 - c_1) e^{\varepsilon} + c_2 - 1} \leq \frac{e^{\varepsilon} + \frac{c_2 - c_1}{1 - c_1} \frac{(c_2-c_1e^\eps)(1-c_1)}{c_2-c_1} - 1}
           {(1 - c_1) e^{\varepsilon} + c_2 - 1} = 1.\]
Therefore, for any $A \subseteq \X$, we have $\bQ^\star_{\varepsilon, \delta}(A) \geq 0$. It suffices to show that for every measurable subset $A \subseteq \mathcal{X}$, we also have
\[
\bQ^\star_{\varepsilon, \delta}(A) \le 1,
\]
% since it then follows that
% \[
% \bQ^\star_{\varepsilon, \delta}(A) 
% = 1 - \bQ^\star_{\varepsilon, \delta}(A^c) \ge 0,
% \]
which implies $\bQ^\star_{\varepsilon, \delta}$ is a probability measure. We proceed by bounding:
\begin{align*}
\bQ^\star_{\varepsilon, \delta}(A) 
& = \lambda^\star_{\varepsilon, \delta} P(A) 
  + \bigl(1 - \lambda^\star_{\varepsilon, \delta}\bigr)\mu(A) \\
& \le \lambda^\star_{\varepsilon, \delta} 
  \Bigl(1 - c_1\bigl(1 - \mu(A)\bigr)\Bigr)
  + \bigl(1 - \lambda^\star_{\varepsilon, \delta}\bigr)\mu(A) \\
& = \Bigl(1 - \bigl(1 - c_1\bigr)\lambda^\star_{\varepsilon, \delta}\Bigr)\mu(A)
  + \lambda^\star_{\varepsilon, \delta}(1 - c_1) \\
& \le \frac{(c_2 - c_1)\,(1 - \delta)}{\bigl(1 - c_1\bigr)e^{\varepsilon} + c_2 - 1}
  \;+\; \frac{\bigl(1 - c_1\bigr)e^{\varepsilon} + (c_2 - c_1)\delta 
  - \bigl(1 - c_1\bigr)}{\bigl(1 - c_1\bigr)e^{\varepsilon} + c_2 - 1} \\
& = \frac{\bigl(1 - c_1\bigr)e^{\varepsilon} + c_2 - 1}
  {\bigl(1 - c_1\bigr)e^{\varepsilon} + c_2 - 1} \\
& = 1.
\end{align*}
The inequality in the second line uses the fact that
\[
P(A) 
\;\le\; 1 - c_1\bigl(1 - \mu(A)\bigr).
\]
This follows from the definition
\[
\tilde{\mathcal{P}}_{c_1, c_2, \mu}
:= \Bigl\{\,P \in \mathcal{P}(\mathcal{X}) :\; P \ll \mu,\quad 
  c_1 \le \tfrac{dP}{d\mu} \le c_2 \;\;\mu\text{-a.e.}\Bigr\},
\]
and the fact that for $P \in \tilde{\mathcal{P}}_{c_1, c_2, \mu}$:
\[
c_1 \,\mu(A^c) \;\le\; P(A^c) 
\;\;\implies\;\;
1 - c_1\,\mu(A^c) 
\;\ge\; 1 - P(A^c) 
\;\;\implies\;\;
1 - c_1\bigl(1 - \mu(A)\bigr) 
\;\ge\; P(A).
\]
And the inequality in the fourth line comes by plugging $\lambda^\star_{\varepsilon, \delta} 
    = \frac{e^{\varepsilon} + \frac{c_2 - c_1}{1 - c_1} \delta - 1}
           {(1 - c_1) e^{\varepsilon} + c_2 - 1}$ into the equation and the fact that $\mu(A) \leq 1$. Hence, $\bQ^\star_{\varepsilon, \delta}$ is indeed a probability measure.


    \medskip

    \textbf{Step 2: Proof of $\bQ^\star_{\varepsilon, \delta} \in \mathcal{Q}_{\mathcal{X}, \tilde{\mathcal{P}}, \varepsilon, \delta}$ } 

Now, we want to show that $\bQ^\star_{\varepsilon, \delta}$ is an $(\eps,\delta)$-LDP sampler. For this purpose, we have to show that for each measurable set $A \subseteq \mathcal{X}$ and any pair of probability measures $P$ and $P'$, we have:
\[e^\eps\big(\bQ^\star_{\varepsilon, \delta}(P)(A)\big) + \delta - \bQ^\star_{\varepsilon, \delta}(P')(A) \geq 0.\]

First, suppose \(\mu(A) \le \frac{1 - c_1}{c_2 - c_1}\). We then have:
\begin{align*}
 e^\varepsilon \Bigl(\lambda^\star_{\varepsilon,\delta}& P(A) 
   +   \bigl(1 - \lambda^\star_{\varepsilon,\delta}\bigr)\mu(A)\Bigr)
 \;+\;\delta
 \;-\;\Bigl(\lambda^\star_{\varepsilon,\delta} P'(A) 
   + \bigl(1 - \lambda^\star_{\varepsilon,\delta}\bigr)\mu(A)\Bigr)\\
&\ge 
  e^\varepsilon \Bigl(\lambda^\star_{\varepsilon,\delta}\,c_1\,\mu(A) 
  + \bigl(1 - \lambda^\star_{\varepsilon,\delta}\bigr)\mu(A)\Bigr)
 \;+\;\delta
 \;-\;
 \Bigl(\lambda^\star_{\varepsilon,\delta}\,c_2\,\mu(A)
   + \bigl(1 - \lambda^\star_{\varepsilon,\delta}\bigr)\mu(A)\Bigr)\\
&=\Big(\lambda^\star_{\varepsilon,\delta}\big( e^\eps c_1 - e^\eps - c_2 + 1\big) + \big(e^\eps - 1\big)\Big)
   \;\mu(A)
 \;+\;\delta\\
&= -\,\frac{c_2 - c_1}{1 - c_1}\,\delta\,\mu(A) 
   \;+\;\delta\\
&\ge 0.
\end{align*}

Next, suppose \(\mu(A) \ge \frac{1 - c_1}{c_2 - c_1}\). Then:
\begin{align*}
e^{\varepsilon} \Big(&\lambda^\star_{\varepsilon,\delta}  P(A)  +  \big(1 - \lambda^\star_{\varepsilon,\delta}\big) \mu(A)\Big) + \delta - \Big(\lambda^\star_{\varepsilon,\delta} P'(A) + \big(1 - \lambda^\star_{\varepsilon,\delta}\big) \mu(A)\Big) \\ 
& \geq e^{\varepsilon} \Big(\lambda^\star_{\varepsilon,\delta} c_1 \mu(A) + \big(1 - \lambda^\star_{\varepsilon,\delta}\big) \mu(A)\Big) + \delta - \Big(\lambda^\star_{\varepsilon,\delta} \big(1 - c_1 + c_1 \mu(A)\big) + \big(1 - \lambda^\star_{\varepsilon,\delta}\big) \mu(A)\Big) \\
& = \Big(e^{\varepsilon} \big(1 - (1 - c_1) \lambda^\star_{\varepsilon,\delta}\big) - \big(1 - (1 - c_1) \lambda^\star_{\varepsilon,\delta}\big)\Big) \mu(A) - \lambda^\star_{\varepsilon,\delta}(1 - c_1) + \delta \\
& = (e^{\varepsilon} - 1)\big(1 - (1 - c_1) \lambda^\star_{\varepsilon,\delta}\big) \mu(A) - \lambda^\star_{\varepsilon,\delta}(1 - c_1) + \delta \\
& \geq (e^{\varepsilon} - 1)\big(1 - (1 - c_1) \lambda^\star_{\varepsilon,\delta}\big) \Big(\frac{1 - c_1}{c_2 - c_1}\Big) - \lambda^\star_{\varepsilon,\delta}(1 - c_1) + \delta \\
& = \frac{(e^{\varepsilon} - 1)(1 - \delta)(1 - c_1)}{(1 - c_1)e^{\varepsilon} + c_2 - 1} - \frac{(1 - c_1)e^{\varepsilon} + (c_2 - c_1)\delta - (1 - c_1)}{(1 - c_1)e^{\varepsilon} + c_2 - 1} + \delta \\
& = \frac{-\delta(1 - c_1)(e^{\varepsilon} - 1) - (c_2 - c_1)\delta}{(1 - c_1)e^{\varepsilon} + c_2 - 1} + \delta \\
& = 0.
\end{align*}

Therefore, in both cases:
\begin{align*}
e^\varepsilon \Bigl(\lambda^\star_{\varepsilon,\delta} P(A) 
   + \bigl(1 - \lambda^\star_{\varepsilon,\delta}\bigr)\mu(A)\Bigr)
\;+\;\delta
\;-\;
\Bigl(\lambda^\star_{\varepsilon,\delta} P'(A) 
   + \bigl(1 - \lambda^\star_{\varepsilon,\delta}\bigr)\mu(A)\Bigr)
\;\;\ge\;0,
\end{align*}


which establishes the \((\varepsilon,\delta)\)-LDP criterion for every measurable set \(A\). Hence the sampler is \((\varepsilon,\delta)\)-LDP.


    \textbf{Step 3: Proof of optimality, lower bound}

    In this part, we show that:
    \[\frac{1 - r_1}{r_2 - r_1} \, f(r_2) + \frac{r_2 - 1}{r_2 - r_1} \, f(r_1) \leq \mathcal{R}\big(\mathcal{Q}_{\mathcal{X}, \tilde{\mathcal{P}}, \varepsilon, \delta}, \tilde{\mathcal{P}}, f\big).\]
     To this end, we need to show that for each $\bQ \in \mathcal{Q}_{\mathcal{X}, \tilde{\mathcal{P}}, \varepsilon, \delta}$, we have
     \[\frac{1 - r_1}{r_2 - r_1} \, f(r_2) + \frac{r_2 - 1}{r_2 - r_1} \, f(r_1) \le \sup_{P \in \tilde{\mathcal{P}}} 
        \;D_f\big(P \, \| \, \bQ(P)\big)\]
       
When $\frac{c_2 - c_1}{1 - c_1} \in \mathbb{N}$, let $\bQ$ be an $(\varepsilon, \delta)$-LDP sampler, i.e, $\bQ \in \mathcal{Q}_{\mathcal{X}, \tilde{\mathcal{P}}, \varepsilon, \delta}$ Define
\[
\alpha := \frac{1 - c_1}{c_2 - c_1}
\quad ,
\quad t := \alpha^{-1} \in \mathbb{N}.
\]
Since $\mu$ is $(\alpha,\frac{1}{\alpha},1)$-decomposable with $\frac{1}{\alpha} \in \mathbb{N}$, we can find disjoint sets $A_1, \ldots, A_t \subseteq \mathcal{X}$ with $\mu(A_i) = \alpha$. Define $P_i$ as a probability measure by $\frac{dP_i}{d\mu} = p_i$ defined as:
\[
p_i(x) = c_2 \mathbf{1}_{A_i}(x) + c_1 \mathbf{1}_{A_i^c}(x),
\]
where \[
\mathbf{1}_{A_i}(x) =
\begin{cases}
1 & \text{if } x \in A_i, \\
0 & \text{otherwise}.
\end{cases}
\]
Let $\bQ_i := \bQ(P_i)$. We aim to show:
\[
\min_i \bQ_i(A_i) \leq \frac{e^\varepsilon + (t-1)\delta}{e^\varepsilon + t - 1}.
\]

By the $(\varepsilon, \delta)$-LDP property, for every $i > 1$, we have
\[
e^\varepsilon \bQ_1(A_i) + \delta \geq \bQ_i(A_i).
\]
Thus,
\begin{equation*}
    \sum_{i=2}^t \Big( e^\varepsilon \bQ_1(A_i) + \delta \Big) \geq \sum_{i=2}^t \bQ_i(A_i).
\end{equation*}



Note that
\begin{align*}
    \sum_{i=2}^t \Big( e^\varepsilon \bQ_1(A_i) + \delta\Big)
& = e^\varepsilon \bQ_1\left(\bigcup_{i=2}^t A_i\right) + (t-1)\delta \\ &
\leq e^\varepsilon \left(1 - \bQ_1(A_1)\right) + (t-1)\delta
\\ & \leq e^\varepsilon \left(1 - \min_i \bQ_i(A_i)\right) + (t-1)\delta,
\end{align*}
and
\begin{equation*}
    \sum_{i=2}^t \bQ_i(A_i) \geq (t-1) \min_i \bQ_i(A_i).
\end{equation*}


Therefore,
\[
\min_i \bQ_i(A_i) \leq \frac{e^\varepsilon + (t-1)\delta}{e^\varepsilon + t-1}.
\]

Rewriting in terms of $\alpha$, we have
\begin{equation}\label{eqn: min Q_i}
\min_i \bQ_i(A_i) \leq \frac{\alpha e^\varepsilon + (1-\alpha)\delta}{\alpha e^\varepsilon + 1-\alpha}.    
\end{equation}



\medskip

Next, we lower bound the worst-case $f$-divergence:
\[
\sup_i D_f(P_i \,\|\, \bQ(P_i))
\geq \sup_i D_f^B\left(c_2 \alpha \,\|\, \bQ_i(A_i)\right),
\]
where $D_f^B(\lambda_1 \,\|\, \lambda_2)$ denotes the $f$-divergence between Bernoulli distributions with $\Pr(1) = \lambda_1$ and $\lambda_2$:
\[
D^B_f(\lambda_1 \parallel \lambda_2) = \lambda_2 f \left( \frac{\lambda_1}{\lambda_2} \right) + (1 - \lambda_2) f \left( \frac{1 - \lambda_1}{1 - \lambda_2} \right).
\]

For each $A_i$, the push-forward measures of $P_i$ and $\bQ_i$ by the indicator function $\mathbf{1}_{A_i}$ are Bernoulli distributions with $\Pr(1) = c_2 \alpha$ and $\bQ_i(A_i)$, respectively. By the data processing inequality (Theorem \ref{thm: dpi}), we have:
\[
D_f ( P_i\| \bQ_i) ) \geq D_f^B \left( c_2 \alpha \| \bQ_i(A_i) \right).
\]

Therefore, we have:
\begin{align*}
    \sup_{P \in \tilde{\mathcal{P}}} 
        \;D_f\big(P \, \| \, \bQ(P)\big) & \geq \sup_{i \in [t]} 
        \; D_f ( P_i \, \| \, \bQ_i) \\ & \geq  \sup_{i \in [t]} 
        \; D_f^B \left( c_2 \alpha \, \| \, \bQ_i(A_i) \right) \\ & \geq 
        D_f^B \left( c_2 \alpha \, \| \, \frac{\alpha e^\varepsilon + (1-\alpha)\delta}{\alpha e^\varepsilon + 1-\alpha} \right) 
\end{align*}

where the inequality in the last line follows from (\ref{eqn: min Q_i}) and the fact that if $\delta \leq \frac{(c_2-c_1e^\eps)(1-c_1)}{c_2-c_1}$, then $ \frac{\alpha e^\varepsilon + (1-\alpha)\delta}{\alpha e^\varepsilon + 1-\alpha}  \leq c_2 \alpha$ and $D_f^{\mathrm{B}}(\lambda_1 \,\|\, \lambda_2)$ is decreasaing in $\lambda_2 \in [0 , \lambda_1]$.

Here is the proof that if $\delta \leq \frac{(c_2-c_1e^\eps)(1-c_1)}{c_2-c_1}$, then $ \frac{\alpha e^\varepsilon + (1-\alpha)\delta}{\alpha e^\varepsilon + 1-\alpha}  \leq c_2 \alpha$.

\begin{eqnarray*}
    \frac{\alpha e^\eps+(1-\alpha)\delta}{\alpha e^\eps+1-\alpha}&=&\frac{(1-c_1) e^\eps+(c_2-1)\delta}{(1-c_1) e^\eps+c_2-1}\\
    &\leq&\frac{(1-c_1) e^\eps+(c_2-1)\left(\frac{(c_2-c_1e^\eps)(1-c_1)}{c_2-c_1}\right)}{(1-c_1) e^\eps+c_2-1}\\
    &=&\frac{1-c_1}{c_2-c_1}\left(\frac{(c_2-c_1)e^\eps+(c_2-1)(c_2-c_1e^\eps)}{(1-c_1)e^\eps+c_2-1}\right)\\
    &=&\frac{c_2(1-c_1)}{c_2-c_1}\\
    &=&c_2\alpha.
\end{eqnarray*}

Therefore, it suffices to compute $D_f^B \left( c_2 \alpha \, \| \, \frac{\alpha e^\varepsilon + (1-\alpha)\delta}{\alpha e^\varepsilon + 1-\alpha} \right)$. It can be shown that
\begin{align*}
  D_f^B\left(c_2 \alpha \,\Big\|\, \frac{\alpha e^\varepsilon + (1-\alpha)\delta}{\alpha e^\varepsilon + 1-\alpha}\right) &
  = \frac{\alpha e^\varepsilon + (1-\alpha)\delta}{\alpha e^\varepsilon + 1-\alpha}
f\left( c_2 \alpha \frac{\alpha e^\varepsilon + 1-\alpha}{\alpha e^\varepsilon + (1-\alpha)\delta} \right) \\
& + \frac{(1-\alpha)(1-\delta)}{\alpha e^\varepsilon + 1-\alpha}
f\left( c_1 (1-\alpha) \frac{\alpha e^\varepsilon + 1-\alpha}{(1-\alpha)(1-\delta)} \right) \\
&= \frac{(1-c_1)e^\varepsilon + (c_2-1)\delta}{(1-c_1)e^\varepsilon + c_2-1}
f\left( \frac{c_2}{c_2-c_1} \frac{(1-c_1)e^\varepsilon + c_2-1}{e^\varepsilon + \frac{c_2-1}{1-c_1}\delta} \right) \\
&+ \frac{(c_2-1)(1-\delta)}{(1-c_1)e^\varepsilon + c_2-1}
f\left( \frac{c_1}{c_2-c_1} \frac{(1-c_1)e^\varepsilon + c_2-1}{1-\delta} \right).
\end{align*}


% Thus,
% \begin{align*}
% \sup_i D_f(P_i \,\|\, \bQ(P_i))
% &\geq D_f^B\left(c_2\alpha \,\Big\|\, \frac{\alpha e^\varepsilon + (1-\alpha)\delta}{\alpha e^\varepsilon + 1 - \alpha}\right) \\
% &= \frac{\alpha e^\varepsilon + (1-\alpha)\delta}{\alpha e^\varepsilon + 1-\alpha} f\left(c_2\alpha \left( \frac{\alpha e^\varepsilon + 1-\alpha}{\alpha e^\varepsilon + (1-\alpha)\delta} \right)\right) \\
% &\quad + \frac{(1-\alpha)(1-\delta)}{\alpha e^\varepsilon + 1-\alpha} f\left(c_1(1-\alpha) \left( \frac{\alpha e^\varepsilon + 1-\alpha}{(1-\alpha)(1-\delta)} \right)\right) \\
% &= \frac{(1-c_1)e^\varepsilon + (c_2-1)\delta}{(1-c_1)e^\varepsilon + c_2-1} f\left( \frac{c_2}{c_2-c_1} \frac{(1-c_1)e^\varepsilon + c_2-1}{e^\varepsilon + \frac{c_2-1}{1-c_1}\delta} \right) \\
% &\quad + \frac{(c_2-1)(1-\delta)}{(1-c_1)e^\varepsilon + c_2-1} f\left( \frac{c_1}{c_2-c_1} \frac{(1-c_1)e^\varepsilon + c_2-1}{1-\delta} \right).
% \end{align*}
% The first inequality 

Defining
\[
    r_1 := \frac{c_1}{c_2 - c_1} \cdot \frac{(1 - c_1)e^{\varepsilon} + c_2 - 1}{1 - \delta},
    \qquad
    r_2 := \frac{c_2}{c_2 - c_1} \cdot \frac{(1 - c_1)e^{\varepsilon} + c_2 - 1}{e^{\varepsilon} + \frac{c_2 - 1}{1 - c_1} \delta},
    \]
 we have   
\begin{align*}
1 - r_1 &= \frac{(1-c_1)\left(c_2 - c_1 e^\varepsilon - \frac{c_2-c_1}{1-c_1}\delta\right)}{(c_2-c_1)(1-\delta)}, \\
r_2 - 1 &= \frac{(c_2-1)\left(c_2 - c_1 e^\varepsilon - \frac{c_2-c_1}{1-c_1}\delta\right)}{(c_2-c_1)\left(e^\varepsilon + \frac{c_2-1}{1-c_1}\delta\right)}.
\end{align*}
Thus
\begin{align*}
\frac{1 - r_1}{r_2 - r_1} &= \frac{(1-c_1)e^\varepsilon + (c_2-1)\delta}{(1-c_1)e^\varepsilon + c_2-1}, \\
\frac{r_2 - 1}{r_2 - r_1} &= \frac{(c_2-1)(1-\delta)}{(1-c_1)e^\varepsilon + c_2-1}.
\end{align*}
Hence, for each $\bQ \in \mathcal{Q}_{\mathcal{X}, \tilde{\mathcal{P}}, \varepsilon, \delta}$, we have
     \[\frac{1 - r_1}{r_2 - r_1} \, f(r_2) + \frac{r_2 - 1}{r_2 - r_1} \, f(r_1) \le \sup_{P \in \tilde{\mathcal{P}}} 
        \;D_f\big(P \, \| \, \bQ(P)\big).\]
        
    \textbf{Step 4: Proof of optimality, upper bound (achievability part)}

    In this part, we show that:
    \[\mathcal{R}\big(\mathcal{Q}_{\mathcal{X}, \tilde{\mathcal{P}}, \varepsilon, \delta}, \tilde{\mathcal{P}}, f\big)
       \le \frac{1 - r_1}{r_2 - r_1} \, f(r_2) + \frac{r_2 - 1}{r_2 - r_1} \, f(r_1).\]
        To this end, we need to show that for each $P \in  \tilde{\mathcal{P}}$, we have:
        \[D_f\big(P \, \| \, \bQ^\star_{\varepsilon, \delta}(P)\big) \le \frac{1 - r_1}{r_2 - r_1} \, f(r_2) + \frac{r_2 - 1}{r_2 - r_1} \, f(r_1).\]

    Fix \( P \in \tilde{\mathcal{P}} \). Let \( p = \frac{dP}{d\mu} \) and \( q = \frac{d\bQ^\star_{\eps,\delta}(P)}{d\mu} \). We first claim that
\[
r_1 \leq \frac{p(x)}{q(x)} \leq r_2
\]
for \(\mu\)-almost every \( x \in \mathcal{X} \). Indeed, we have
\[
\frac{p(x)}{q(x)}
= \frac{p(x)}{\lambda^\star_{\varepsilon,\delta} p(x) + (1 - \lambda^\star_{\varepsilon,\delta})}
,\]
which is an increasing function of \( p(x) \geq 0 \) as long as $\lambda^\star_{\eps,\delta} \leq 1$. Now we show that if $\delta \leq \frac{(c_2-c_1e^\eps)(1-c_1)}{c_2-c_1}$, then $\lambda^\star_{\varepsilon, \delta} 
    = \frac{e^{\varepsilon} + \frac{c_2 - c_1}{1 - c_1} \delta - 1}
           {(1 - c_1) e^{\varepsilon} + c_2 - 1} \leq 1$. 
           \begin{align*}
               \lambda^\star_{\varepsilon, \delta} 
   & = \frac{e^{\varepsilon} + \frac{c_2 - c_1}{1 - c_1} \delta - 1}
           {(1 - c_1) e^{\varepsilon} + c_2 - 1} \\
           &\leq  \frac{e^{\varepsilon} + \frac{c_2 - c_1}{1 - c_1} \Big( \frac{(c_2-c_1e^\eps)(1-c_1)}{c_2-c_1}\Big) - 1}
           {(1 - c_1) e^{\varepsilon} + c_2 - 1}.\\
           & = \frac{e^{\varepsilon} +  c_2-c_1e^\eps - 1}
           {(1 - c_1) e^{\varepsilon} + c_2 - 1} \\
           & = 1.
           \end{align*}
           Since \( p(x) \) is bounded between \( c_1 \) and \( c_2 \), it follows that
\begin{equation}\label{eqn: p/q ratio}
    \frac{c_1}{\lambda^\star_{\varepsilon,\delta} c_1 + (1 - \lambda^\star_{\varepsilon,\delta})}
\leq \frac{p(x)}{q(x)}
\leq \frac{c_2}{\lambda^\star_{\varepsilon,\delta} c_2 + (1 - \lambda^\star_{\varepsilon,\delta})}
\end{equation}
for \(\mu\)-almost every \( x \in \mathcal{X} \).
From the premise of the proposition, we have
\begin{equation*}
    \lambda^\star_{\varepsilon, \delta} 
    = \frac{e^{\varepsilon} + \frac{c_2 - c_1}{1 - c_1} \delta - 1}
           {(1 - c_1) e^{\varepsilon} + c_2 - 1}
\end{equation*}
It can be shown that:
\[
\frac{c_1}{c_2 -c_1}\frac{  (1-c_1) e^{\varepsilon} + c_2 - 1 }{(1-\delta)}
\leq \frac{p(x)}{q(x)}
\leq \frac{c_2}{c_2 - c_1} \cdot \frac{(1 - c_1)e^{\varepsilon} + c_2 - 1}{e^{\varepsilon} + \frac{c_2 - 1}{1 - c_1} \delta}
\]
which concludes
\[
r_1 \leq \frac{p(x)}{q(x)} \leq r_2
\]
for \(\mu\)-almost every \( x \in \mathcal{X} \).
To finalize the proof of the achievability part, we need the following lemma.

\begin{lemma}[{[Theorem 2.1]{\citep{rukhin1997information}}}]\label{lem: bounded f-div}
Let \( P, Q \in \mathcal{P}(\mathcal{X}) \). Suppose that \( P \) and \( Q \) are both absolutely continuous with respect to a reference measure \( \mu \) on \( \mathcal{X} \). Assume that there exist \( r_1, r_2 \in \mathbb{R} \) with \( 0 \leq r_1 < 1 < r_2 \) such that the corresponding densities \( p = \frac{dP}{d\mu} \) and \( q = \frac{dQ}{d\mu} \) satisfy \( q(x) > 0 \) and
\[
r_1 \leq \frac{p(x)}{q(x)} \leq r_2
\quad \text{for } \mu\text{-almost every } x \in \mathcal{X}.
\]
Then, for any \( f \)-divergence \( D_f \), it holds that
\[
D_f(P\|Q) \leq \frac{1 - r_1}{r_2 - r_1} f(r_2) + \frac{r_2 - 1}{r_2 - r_1} f(r_1).
\]
\end{lemma}

Lemma \ref{lem: bounded f-div}, directly shows that for each $P \in  \tilde{\mathcal{P}}$, we have:
        \[D_f\big(P \, \| \, \bQ^\star_{\varepsilon, \delta}(P)\big) \le \frac{1 - r_1}{r_2 - r_1} \, f(r_2) + \frac{r_2 - 1}{r_2 - r_1} \, f(r_1)\]
for
\[
    r_1 = \frac{c_1}{c_2 - c_1} \cdot \frac{(1 - c_1)e^{\varepsilon} + c_2 - 1}{1 - \delta},
    \qquad
    r_2 = \frac{c_2}{c_2 - c_1} \cdot \frac{(1 - c_1)e^{\varepsilon} + c_2 - 1}{e^{\varepsilon} + \frac{c_2 - 1}{1 - c_1} \delta}.
    \]
\end{proof}

\begin{corollary}[$g_{\eps,\delta}$-FLDP as a special case]\label{cor: special case of f-LDP approximate}

Let $\tilde{\mathcal{P}} = \tilde{\mathcal{P}}_{c_1, c_2, \mu}$. Under Assumption \ref{assumption: norm}, suppose the reference measure $\mu$ is $(\alpha,\frac{1}{\alpha},1)$-decomposable with $\alpha = \frac{1 - c_1}{c_2 - c_1}$. Then the sampler
\begin{equation}
    \bQ^\star_{c_1,c_2,\mu,g_{\varepsilon,\delta}}(P) \coloneqq \lambda^\star_{c_1,c_2,g_{\varepsilon,\delta}} P + (1 - \lambda^\star_{c_1,c_2,g_{\varepsilon,\delta}})\mu
\end{equation}
belongs to the class $\mathcal{Q}_{\mathcal{X}, \tilde{\mathcal{P}}, g_{\varepsilon,\delta}}$ and is minimax-optimal with respect to any $f$-divergence $D_f$, that is,
\begin{equation}
    \sup_{P \in \tilde{\mathcal{P}}} D_f\big(P \, \| \, \bQ^\star_{c_1,c_2,\mu,g_{\varepsilon,\delta}}(P)\big)
    = \mathcal{R}\big(\mathcal{Q}_{\mathcal{X}, \tilde{\mathcal{P}}, g_{\varepsilon,\delta}}, \tilde{\mathcal{P}}, f\big),
\end{equation}
where 
\[
\lambda^\star_{c_1,c_2,g_{\varepsilon,\delta}} = \frac{e^{\varepsilon} + \frac{c_2 - c_1}{1 - c_1} \, \delta - 1}{(1 - c_1) e^{\varepsilon} + c_2 - 1}.
\]

\end{corollary}


% \subsection{Proof of Corollary \ref{cor: special case of f-LDP approximate}}
% \HG{Add the proof of the 
% parts corresponding to the optimal value}
\begin{proof}
% \HG{to be checked}
From Theorem \ref{thm: global functional}, we know that the sampler \[\bQ^\star_{c_1,c_2,\mu,g_{\varepsilon,\delta}}(P) \coloneqq \lambda^\star_{c_1,c_2,g_{\varepsilon,\delta}} P + (1 - \lambda^\star_{c_1,c_2,g_{\varepsilon,\delta}})\mu\]  for $\lambda^\star_{c_1,c_2,g_{\varepsilon,\delta}}$ defined as \[
    \lambda^\star_{c_1,c_2,g_{\varepsilon,\delta}} 
    \coloneq \inf_{\beta \geq 0} 
    \frac{e^\beta + \frac{c_2 - c_1}{1 - c_1} \big( 1 + g_{\varepsilon,\delta}^*(-e^\beta) \big) - 1}
         {(1 - c_1)e^\beta + c_2 - 1}.
    \] 
    belongs to the class $\mathcal{Q}_{\mathcal{X}, \tilde{\mathcal{P}}, g_{\varepsilon,\delta}}$ and is minimax-optimal with respect to any $f$-divergence $D_f$. Therefore, it suffices to compute $g_{\varepsilon,\delta}^*(-e^\varepsilon)$ and then solve the minimization problem.
    

   
A direct comparison of the three affine pieces shows

\begin{equation}
    g_{\varepsilon,\delta}(\theta)=
\begin{cases}
1-\delta-e^{\varepsilon}\theta, & 0\le\theta\le \frac{1-\delta}{e^{\varepsilon}+1},\\[4pt]
e^{-\varepsilon}(1-\delta-\theta), & \frac{1-\delta}{e^{\varepsilon}+1}\le\theta\le 1-\delta,\\[4pt]
0, & 1-\delta\le\theta\le 1,\\[4pt]
+\infty, & \text{otherwise.}
\end{cases}\label{eq: piecewise-f}
\end{equation}

For any \(y\in\mathbb{R}\),
\(
g_{\varepsilon,\delta}^{*}(y)
   =\sup\limits_{0\le\theta\le1}(\theta y-g_{\varepsilon,\delta}(\theta))
\)
splits naturally into the supremum over the three intervals of~\eqref{eq: piecewise-f}.
Denote the corresponding optimized values by
\[
\begin{aligned}
S_1(y)&:=\sup_{0\le\theta\le  \frac{1-\delta}{e^{\varepsilon}+1}}\bigl\{\theta y-(1-\delta-e^{\varepsilon}\theta)\bigr\},\\
S_2(y)&:=\sup_{ \frac{1-\delta}{e^{\varepsilon}+1}\le\theta\le 1 - \delta}\bigl\{\theta y-e^{-\varepsilon}(1-\delta-\theta)\bigr\},\\
S_3(y)&:=\sup_{1 - \delta\le\theta\le 1}\bigl\{\theta y\bigr\}.
\end{aligned}
\]

\vspace{4pt}
\noindent
We now optimize over each sub-interval.

\emph{(i) Interval \(0\le\theta\le  \frac{1-\delta}{e^{\varepsilon}+1}\).}  
Writing the objective as \((y+e^{\varepsilon})\theta-(1-\delta)\), it is linear in \(\theta\).
Hence  
\[
S_1(y)=
\begin{cases}
-(1-\delta), & y\le -e^{\varepsilon},\\[4pt]
\bigl(y+e^{\varepsilon}\bigr) \frac{1-\delta}{e^{\varepsilon}+1}-(1-\delta)
      =\dfrac{1-\delta}{e^{\varepsilon}+1}(y-1),
      & y\ge -e^{\varepsilon}.
\end{cases}
\]

\emph{(ii) Interval \(\frac{1-\delta}{e^{\varepsilon}+1}\le\theta\le 1-\delta\).}  
Here the objective equals \((y+e^{-\varepsilon})\theta-e^{-\varepsilon}(1-\delta)\).
Thus  
\[
S_2(y)=
\begin{cases}
\dfrac{1-\delta}{e^{\varepsilon}+1}(y-1), & y\le -e^{-\varepsilon},\\[6pt]
(1-\delta)y, & y\ge -e^{-\varepsilon}.
\end{cases}
\]

\emph{(iii) Interval \(1 - \delta\le\theta\le 1\).}  
The objective is simply \(y\theta\), so  
\[
S_3(y)=
\begin{cases}
(1-\delta)\,y, & y< 0,\\[4pt]
y, & y\ge 0.
\end{cases}
\]

We now take the overall supremum.

Comparing \(S_1,S_2,S_3\) on the four regimes  
\(\bigl(-\infty,-e^{\varepsilon}\bigr)\),
\(\bigl[-e^{\varepsilon},-e^{-\varepsilon}\bigr]\),
\(\bigl[-e^{-\varepsilon},0\bigr]\),
\([0,\infty)\)
gives
\[
g_{\varepsilon,\delta}^{*}(y)
   =\max\{S_1(y),S_2(y),S_3(y)\}
   =
\begin{cases}
-(1-\delta),
   & y< -e^{\varepsilon},\\[6pt]
\dfrac{1-\delta}{e^{\varepsilon}+1}(y-1),
   & -e^{\varepsilon}\le y\le -e^{-\varepsilon},\\[10pt]
(1-\delta)\,y,
   & -e^{-\varepsilon}\le y\le 0,\\[6pt]
y,
   & y\ge 0.
\end{cases}
\]

%     \[g_{\varepsilon,\delta}^*(\theta)=
% \begin{cases}
% -(1-\delta), 
%     & \text{if }\theta \le -e^{\varepsilon}, \\[6pt]
% \displaystyle
% (1-\delta)\,
% \frac{(\theta-1)}{e^{\varepsilon}+1},
%     & \text{if } -e^{\varepsilon}\le \theta \le -e^{-\varepsilon},\\[6pt]
% (1-\delta)\,\theta,
%     & \text{if } -e^{-\varepsilon}\le \theta \le 0,\\[6pt]
% \theta,
%     & \text{if } \theta \ge 0.
% \end{cases}\]

\textbf{Step 2: Solve the infimum to obtain $\lambda^\star_{c_1,c_2,g_{\varepsilon,\delta}}$}

 From   \[
    \lambda^\star_{c_1,c_2,g_{\varepsilon,\delta}} 
    = \inf_{\beta \geq 0} 
    \frac{e^\beta + \frac{c_2 - c_1}{1 - c_1} \big( 1 + g_{\eps,\delta}^*(-e^\beta) \big) - 1}
         {(1 - c_1)e^\beta + c_2 - 1}
    \] 
    Hence, we only need to know the value of $g_{\eps,\delta}^*(-e^\beta )$. We know $-e^{\beta} \le -1$ for $\beta \geq 0$. Therefore, $-e^{\beta} \leq - e^{-\eps}$ for given $\eps \geq 0$. i.e.,
    \[g_{\varepsilon,\delta}^{*}(-e^\beta) = \begin{cases}
        -(1-\delta), &  -e^\beta \le -e^\eps,\\[6pt]
        \dfrac{1-\delta}{e^{\varepsilon}+1}(-e^\beta-1), & -e^\eps \le -e^\beta \leq -1.
        
    \end{cases}\]
    



It follows that:
\begin{align*}\lambda^\star_{c_1,c_2,g_{\varepsilon,\delta}} & =  
    % \min\Bigg\{ \inf_{\beta \geq \eps} 
    % \frac{e^\beta + \frac{c_2 - c_1}{1 - c_1} \big( 1 - (1 - \delta) \big) - 1}
    %      {(1 - c_1)e^\beta + c_2 - 1}  , \inf_{\beta \in [0,\eps]} 
    % \frac{e^\beta + \frac{c_2 - c_1}{1 - c_1} \Big( 1 - \frac{1-\delta}{e^{\varepsilon}+1}(e^\beta+1) \Big) - 1}
    %      {(1 - c_1)e^\beta + c_2 - 1}\Bigg\} \\ & = 
         \min\Bigg\{ \inf_{\beta \geq \eps} 
    \frac{e^\beta + \frac{c_2 - c_1}{1 - c_1}  \delta  - 1}
         {(1 - c_1)e^\beta + c_2 - 1} \, , \, \inf_{\beta \in [0,\eps]} 
    \frac{e^\beta + \frac{c_2 - c_1}{1 - c_1} \bigg( 1 - \frac{e^\beta+1}{e^{\varepsilon}+1}(1 - \delta) \bigg) - 1}
         {(1 - c_1)e^\beta + c_2 - 1}\Bigg\} \\
         & = \min\bigg\{\frac{e^\eps + \frac{c_2 - c_1}{1 - c_1}\delta - 1}{(1 - c_1)e^\eps + c_2 - 1} \, , \, \frac{1 - (1-\delta)\frac{2}{e^\eps + 1}}{1 - c_1}\bigg\}.
\end{align*}
Since $\frac{c_2 - c_1}{1 - c_1} \in \mathbb{N}$ and $c_2 > 1$, it follows that $\frac{c_2 - c_1}{1 - c_1} \geq 2$, which in turn implies  $c_1 + c_2 \geq 2$. Now we want to show that if $c_1 + c_2 \geq 2$, then:
\[\frac{e^\eps + \frac{c_2 - c_1}{1 - c_1}\delta - 1}{(1 - c_1)e^\eps + c_2 - 1} \leq \frac{1 - (1-\delta)\frac{2}{e^\eps + 1}}{1 - c_1}.\]

Because $c_1 < 1 < c_2$ and $e^\varepsilon \ge 1$, the denominator
\(
(1-c_1)e^\varepsilon + c_2 - 1
\)
is strictly positive, so we can multiply both sides of the
inequality by it without changing the sign.

Set
\[
A \coloneqq (1 - c_1)e^\varepsilon + c_2 - 1 \;>\; 0.
\]
The claim is equivalent to
\[
(1 - c_1)(e^\varepsilon - 1) + (c_2 - c_1)\,\delta
\;\le\;
\Bigl(1 - (1-\delta)\,\tfrac{2}{e^\varepsilon + 1}\Bigr)A .
\] 
Bring the left–hand side to the right and factor out $(1-\delta)$:  
\begin{align*}
0
&\;\le\;
A - (1-\delta)\,\frac{2A}{e^\varepsilon + 1}
      -\bigl[(1 - c_1)(e^\varepsilon - 1) + (c_2 - c_1)\delta\bigr]\\[2pt]
&\;=\;
(c_2 - c_1)(1-\delta) - (1-\delta)\,\frac{2A}{e^\varepsilon + 1}\\[2pt]
&\;=\;
(1-\delta)\Bigl[(c_2 - c_1) - \frac{2A}{e^\varepsilon + 1}\Bigr].
\end{align*}

Since $1-\delta \ge 0$, we only need to prove  
\[
(c_2 - c_1)(e^\varepsilon + 1) \;\ge\; 2A .
\]

Substituting $A$ gives
\[
(c_2 - c_1)e^\varepsilon + (c_2 - c_1)
\;\ge\;
2(1 - c_1)e^\varepsilon + 2(c_2 - 1).
\]
Rearranging, we have:
\[
\bigl(c_2 + c_1 - 2\bigr)(e^\varepsilon - 1)\;\ge\;0.
\]

Because $e^\varepsilon - 1 \ge 0$ for every $\varepsilon \ge 0$,
the last inequality holds precisely when $c_1 + c_2 - 2 \ge 0$,
i.e.\ when $c_1 + c_2 \ge 2$.  This is exactly the hypothesis,
so the desired inequality is proved and we have:
\[\lambda^\star_{c_1,c_2,g_{\varepsilon,\delta}} = \frac{e^{\varepsilon} + \frac{c_2 - c_1}{1 - c_1} \, \delta - 1}{(1 - c_1) e^{\varepsilon} + c_2 - 1}.\]

\textbf{Step 3: Optimality proof}

Following Theorem \ref{thm: global functional}, for the obtained value of $\lambda_{c_1,c_2,g_{\eps,\delta}}$, 
the sampler \begin{equation*}
    \bQ^\star_{c_1,c_2,\mu,g_{\varepsilon,\delta}}(P) \coloneqq \lambda^\star_{c_1,c_2,g_{\varepsilon,\delta}} P + (1 - \lambda^\star_{c_1,c_2,g_{\varepsilon,\delta}})\mu
\end{equation*}
belongs to the class $\mathcal{Q}_{\mathcal{X}, \tilde{\mathcal{P}}, g_{\varepsilon,\delta}}$ and is minimax-optimal with respect to any $f$-divergence $D_f$, that is,
\begin{equation*}
    \sup_{P \in \tilde{\mathcal{P}}} D_f\big(P \, \| \, \bQ^\star_{c_1,c_2,\mu,g_{\varepsilon,\delta}}(P)\big)
    = \mathcal{R}\big(\mathcal{Q}_{\mathcal{X}, \tilde{\mathcal{P}}, g_{\varepsilon,\delta}}, \tilde{\mathcal{P}}, f\big).
\end{equation*}
\end{proof}


\section{Proofs of the main results
}\label{sec: proofs main}

\subsection{Proof of Proposition \ref{prop: triviality functional}}\label{appendix: proof of triviality functional}
\begin{proof}
    Based on Proposition~6 in~\citet{dong2022gaussian}, for a  trade-off function $g$, a sampler is \emph{$g$-FLDP} if and only if it is $\bigl(\varepsilon,\,1 + g^*(-\,e^\varepsilon)\bigr)$-LDP for all $\varepsilon \ge 0$. In the proof of Proposition \ref{prop: global appx linear}, we showed that if $\delta > \frac{(c_2-c_1e^\eps)(1-c_1)}{c_2-c_1}$, then the identity sampler $\bQ^I(P) = P$ satisfies $(\eps,\delta)$-LDP and has the trivial minimax risk zero.

Suppose the non-triviality condition (\ref{eqn: non-trivial condition condition functional}) does not hold, i.e., for any $\eps \geq 0$, we have:
\[1 + g^*(-e^\eps) > \frac{(c_2-c_1e^\eps)(1-c_1)}{c_2-c_1}.\]

In this case, the trivial sampler $\bQ^I(P) = P$ satisfies $(\eps,1 + g^*(-e^\eps))$-LDP for any $\eps \geq 0$ and therefore satisfies $g$-FLDP and has the trivial minimax risk zero. 
\end{proof}


\subsection{Proof of Theorem \ref{thm: global functional}}\label{proof: theorem global functional}
% \HG{Explain the technical reason for $ratio \in \N$}

\begin{proof}
In this proof, we assume that the non-triviality condition (\ref{eqn: non-trivial condition condition functional}) is satisfied.



The proof of this theorem is directly based on Proposition~\ref{prop: global appx linear}. Accordingly, the global minimax-optimal sampler in Theorem~\ref{thm: global functional} builds upon the minimax-optimal sampler proposed in Proposition~\ref{prop: global appx linear}. The lower bound in the optimality proof of Proposition~\ref{prop: global appx linear} relies on the existence of disjoint sets \( A_1, \ldots, A_t \), each with measure \( \mu(A_i) = \alpha \), where \( t = \frac{1}{\alpha} \). In fact, the converse part of the optimality proof hinges on the \( (\alpha, \frac{1}{\alpha}, 1) \)-decomposability of \( \mu \), where \( \frac{1}{\alpha} = \frac{c_2 - c_1}{1 - c_1} \in \mathbb{N} \). This technical requirement motivates the assumption \( \frac{1}{\alpha} = \frac{c_2 - c_1}{1 - c_1} \in \mathbb{N} \) in Assumption~\ref{assumption: norm}.

\textbf{Step 1: To propose a $g$-FLDP sampler $\bQ^\star_{c_1,c_2,\mu,g}$} 

    From Proposition~\ref{prop: global appx linear}, we know that for each pair $(\varepsilon, \delta)$, the sampler
\[
\bQ^\star_{\varepsilon, \delta}(P) 
\;=\; 
\lambda^\star_{\varepsilon, \delta}\,P \;+\; \bigl(1 \;-\; \lambda^\star_{\varepsilon, \delta}\bigr)\mu,
\]
with 
\[
\lambda^\star_{\varepsilon, \delta} 
\;=\;
\frac{e^{\varepsilon} \;+\; \tfrac{c_2 - c_1}{1 - c_1}\,\delta \;-\; 1}
     {\,(1 - c_1)\,e^{\varepsilon} \;+\; c_2 \;-\; 1},
\]
is minimax-optimal in the class $\mathcal{Q}_{\mathcal{X}, \tilde{\mathcal{P}}, \varepsilon, \delta}$.

Next, based on Proposition~6 in~\citet{dong2022gaussian}, for a trade-off function $g$, a sampler is \emph{$g$-FLDP} if and only if it is $\bigl(\varepsilon,\,1 + g^*(-\,e^\varepsilon)\bigr)$-LDP for all $\varepsilon \ge 0$. Consequently, if we set $\delta(\varepsilon) = 1 + g^*\!\bigl(-\,e^\varepsilon\bigr)$ and define 
\[
\bQ^\star_{\varepsilon,\,\delta(\varepsilon)}(P) 
\;\coloneqq\; 
\lambda^\star_{\varepsilon,\,\delta(\varepsilon)}\,P \;+\; \Bigl(1 - \lambda^\star_{\varepsilon,\,\delta(\varepsilon)}\Bigr)\mu,
\]
then by Proposition~\ref{prop: global appx linear}, we have $\bQ^\star_{\varepsilon,\,\delta(\varepsilon)} \in \mathcal{Q}_{\mathcal{X}, \tilde{\mathcal{P}}, \varepsilon, \,\delta(\varepsilon)}$. 
% Together with~\citet{dong2022gaussian}, this implies $\bQ^\star_{\varepsilon,\,\delta(\varepsilon)}$ is $g$-FLDP if and only if it is $\bigl(\varepsilon, \,\delta(\varepsilon)\bigr)$-LDP for all $\varepsilon \ge 0$.

We then use Lemma~\ref{lem:decrease-lambda}, which states that if $\bQ_{\lambda_1}(P) = \lambda_1\,P + (1 - \lambda_1)\mu$ is $\bigl(\varepsilon,\delta\bigr)$-LDP for $0 \leq \lambda_1 \leq 1$, then for \emph{any} $\lambda_2 \in [0,\,\lambda_1]$, the sampler $\bQ_{\lambda_2}(P) = \lambda_2\,P + (1 - \lambda_2)\mu$ is also $\bigl(\varepsilon,\delta\bigr)$-LDP.  

Therefore, if we define 
\[
\lambda^\star_{{c_1,c_2,g}} 
\;=\;
\inf_{\beta \,\ge\, 0}\;
\frac{\;e^{\beta } \;+\; \tfrac{c_2 - c_1}{\,1 - c_1\,}\,\Bigl[\,1 + g^*(-\,e^\beta )\Bigr] \;-\; 1}
     {\,(1 - c_1)\,e^\beta  + \bigl(c_2 - 1\bigr)},
\]
this guarantees that the sampler
\[
\bQ^\star_{{c_1,c_2,\mu,g}}(P) 
\;\coloneqq\; 
\lambda^\star_{c_1,c_2,g}\,P \;+\; \bigl(1 - \lambda^\star_{c_1,c_2,g}\bigr)\,\mu
\]
is $\bigl(\varepsilon,\,1 + g^*(-e^\varepsilon)\bigr)$-LDP for \emph{all} $\varepsilon \ge 0$. Note that the non-triviality constraint (\ref{eqn: non-trivial condition condition functional}) guarantees that there exists an $\eps \geq 0$ for which $1 + g^*(-e^\eps) \leq \frac{(c_2-c_1e^\eps)(1-c_1)}{c_2-c_1}$ and therefore $\lambda^\star_{c_1,c_2,g} \leq 1$. By Proposition~6 of~\citet{dong2022gaussian}, being $\bigl(\varepsilon,\,1 + g^*(-e^\varepsilon)\bigr)$-LDP for all $\varepsilon \ge 0$ is exactly the definition of being $g$-FLDP. 
Hence, $\bQ^\star_{c_1,c_2,\mu,g}$ is indeed a $g$-FLDP sampler and therefore belongs to $\mathcal{Q}_{\mathcal{X},\,\tilde{\mathcal{P}},\,g}$.  This completes the proof of the first step.


           \textbf{Step 2: To prove the optimality of $\bQ^\star_{c_1,c_2,\mu,g}$ } 

           Proposition~6 of~\citet{dong2022gaussian} tells us that if we set
\[
\delta(\varepsilon) \;=\; 1 \;+\; g^*\!\bigl(-\,e^\varepsilon\bigr),
\]
then for every $\varepsilon \geq 0$ we have
\begin{equation}\label{eqn: functional lower bounded by appx}
    \inf_{\bQ \,\in\, \mathcal{Q}_{\mathcal{X}, \tilde{\mathcal{P}}, g}} 
    \;\sup_{P \,\in\, \tilde{\mathcal{P}}}
    \,D_f\bigl(P \,\|\, \bQ(P)\bigr) 
    \;\;\ge\;\; 
    \inf_{\bQ \,\in\, \mathcal{Q}_{\mathcal{X}, \tilde{\mathcal{P}}, \varepsilon, \delta(\varepsilon)}} 
    \;\sup_{P \,\in\, \tilde{\mathcal{P}}}
    \,D_f\bigl(P \,\|\, \bQ(P)\bigr).
\end{equation}
On the other hand, define the sampler 
\[
\bQ^\star_{c_1,c_2,\mu,g}(P) 
\;\coloneqq\; 
\lambda^\star_{c_1,c_2,g}\,P \;+\; \bigl(1 - \lambda^\star_{c_1,c_2,g}\bigr)\mu,
\]
where
\[
\lambda^\star_{c_1,c_2,g} 
\;=\;
\inf_{\eps  \,\ge\, 0}
\;\;
\frac{\,e^\eps  \;+\; \tfrac{c_2 - c_1}{1 - c_1}\,\Bigl(1 + g^*\!\bigl(-\,e^\eps \bigr)\Bigr) \;-\; 1}
     {\,(1 - c_1)\,e^\eps  + \bigl(c_2 - 1\bigr)}.
\]
Let $\varepsilon^\star_g$ be the value of $\varepsilon$ that achieves this infimum. The infimum is attained at a finite value of $\varepsilon$ because, if we define
\[
h_g(\varepsilon) \coloneqq \frac{\,e^\varepsilon \;+\; \tfrac{c_2 - c_1}{1 - c_1}\,\Bigl(1 + g^*\!\bigl(-\,e^\varepsilon\bigr)\Bigr) \;-\; 1}
{(1 - c_1)\,e^\varepsilon + (c_2 - 1)},
\]
then we have
\[
h_g(0) \leq \lim_{\varepsilon \to \infty} h_g(\varepsilon).
\]
Together with the continuity of $h_g(\varepsilon)$, this implies that the infimum is achieved at some finite $\varepsilon_g^\star$.

To establish the continuity of $h_g(\varepsilon)$, we note that both the numerator and denominator are continuous functions of $\varepsilon$, and the denominator is never zero. To verify that the numerator is continuous, it suffices to show that $g^*$ is continuous. This follows directly from the definition of the convex conjugate and the fact that the original function $g$ is continuous on the compact interval $[0,1]$ and takes values in $[0,1]$ (see Proposition \ref{prop: trade-off}). Hence,
\[
\lambda^\star_{c_1,c_2,g}
\;=\;
\frac{\,e^{\varepsilon^\star_g} \;+\; \tfrac{c_2 - c_1}{1 - c_1}\,\Bigl(1 + g^*\!\bigl(-\,e^{\varepsilon^\star_g}\bigr)\Bigr) \;-\; 1}
     {\,(1 - c_1)\,e^{\varepsilon^\star_g} + \bigl(c_2 - 1\bigr)}.
\]
By construction, $\bQ^\star_{c_1,c_2,\mu,g}$ belongs to the set 
\(\mathcal{Q}_{\mathcal{X}, \tilde{\mathcal{P}}, \varepsilon^\star_g, \delta(\varepsilon^\star_g)}\).  
From Proposition~\ref{prop: global appx linear}, we know $\bQ^\star_{c_1,c_2,\mu,g}$ is minimax-optimal in 
\(\mathcal{Q}_{\mathcal{X}, \tilde{\mathcal{P}}, \varepsilon^\star_g, \delta(\varepsilon^\star_g)}\).  
Hence,
\begin{equation}\label{eqn: functional coincides appx optimal}
    \sup_{P \,\in\, \tilde{\mathcal{P}}}
    \,D_f\Bigl(P \,\bigl\|\, \bQ^\star_{c_1,c_2,\mu,g}(P)\Bigr)
    \;=\;
    \inf_{\bQ \,\in\, \mathcal{Q}_{\mathcal{X}, \tilde{\mathcal{P}}, \varepsilon^\star_g, \delta(\varepsilon^\star_g)}}
    \;\sup_{P \,\in\, \tilde{\mathcal{P}}}
    \,D_f\bigl(P \,\|\, \bQ(P)\bigr).
\end{equation}
Finally, combining \eqref{eqn: functional lower bounded by appx}, \eqref{eqn: functional coincides appx optimal}, and the trivial inequality 
\[
\inf_{\bQ \,\in\, \mathcal{Q}_{\mathcal{X}, \tilde{\mathcal{P}}, g}}
\;\sup_{P \,\in\, \tilde{\mathcal{P}}}
\,D_f\bigl(P \,\|\, \bQ(P)\bigr)
\;\;\le\;\;
\sup_{P \,\in\, \tilde{\mathcal{P}}}
\,D_f\Bigl(P \,\bigl\|\, \bQ^\star_{c_1,c_2,\mu,g}(P)\Bigr),
\]
we conclude:
\begin{align*}
    \sup_{P \,\in\, \tilde{\mathcal{P}}}
\,D_f\Bigl(P \,\bigl\|\, \bQ^\star_{c_1,c_2,\mu,g}(P)\Bigr) &\;=\;
    \inf_{\bQ \,\in\, \mathcal{Q}_{\mathcal{X}, \tilde{\mathcal{P}}, \varepsilon^\star_g, \delta(\varepsilon^\star_g)}}
    \;\sup_{P \,\in\, \tilde{\mathcal{P}}}
    \,D_f\bigl(P \,\|\, \bQ(P)\bigr) 
\\ & \;=\;
    \mathcal{R}\Bigl(\mathcal{Q}_{\mathcal{X}, \tilde{\mathcal{P}}, \varepsilon^\star_g, \delta(\varepsilon^\star_g)}, 
                    \tilde{\mathcal{P}}, f\Bigr)
                      \\
    & \;=\; \inf_{\bQ \,\in\, \mathcal{Q}_{\mathcal{X}, \tilde{\mathcal{P}}, g}}
    \;\sup_{P \,\in\, \tilde{\mathcal{P}}}
    \,D_f\bigl(P \,\|\, \bQ(P)\bigr) \\
    & \;=\; \mathcal{R}\big(\mathcal{Q}_{\mathcal{X}, \tilde{\mathcal{P}}, g}, \tilde{\mathcal{P}}, f\big),
\end{align*}
which completes the proof of the second step.

\textbf{Step 3: Computing the optimal value of $\mathcal{R}\big(\mathcal{Q}_{\mathcal{X}, \tilde{\mathcal{P}}, g}, \tilde{\mathcal{P}}, f\big)$}

It follows from the previous step that in order to compute $\mathcal{R}\big(\mathcal{Q}_{\mathcal{X}, \tilde{\mathcal{P}}, g}, \tilde{\mathcal{P}}, f\big)$, it suffices to compute $\mathcal{R}\Bigl(\mathcal{Q}_{\mathcal{X}, \tilde{\mathcal{P}}, \varepsilon^\star_g, \delta(\varepsilon^\star_g)}, 
                    \tilde{\mathcal{P}}, f\Bigr)$.

                    From Proposition \ref{prop: global appx linear}, we know that
\[ \mathcal{R}\Bigl(\mathcal{Q}_{\mathcal{X}, \tilde{\mathcal{P}}, \varepsilon^\star_g, \delta(\varepsilon^\star_g)}, 
                    \tilde{\mathcal{P}}, f\Bigr) = \frac{1 - r_1}{r_2 - r_1} \, f(r_2) + \frac{r_2 - 1}{r_2 - r_1} \, f(r_1)\]
for
\[
    r_1 = \frac{c_1}{c_2 - c_1} \cdot \frac{(1 - c_1)e^{\varepsilon^\star_g} + c_2 - 1}{1 - \delta(\varepsilon^\star_g)},
    \qquad
    r_2 = \frac{c_2}{c_2 - c_1} \cdot \frac{(1 - c_1)e^{\varepsilon^\star_g} + c_2 - 1}{e^{\varepsilon^\star_g} + \frac{c_2 - 1}{1 - c_1} \delta(\varepsilon^\star_g)}.
    \]

    We now compute $r_1$ and $r_2$ in terms of $\lambda_{c_1,c_2,g}^\star$. For \( \delta(\varepsilon^\star_g) = 1 + g^*(-e^{\varepsilon^\star_g}) \), we have
\[
r_1 = \frac{c_1}{c_2 - c_1} \cdot 
      \frac{(1 - c_1)e^{\varepsilon^\star_g} + c_2 - 1}
           {-g^*(-e^{\varepsilon^\star_g})},
\quad
r_2 = \frac{c_2}{c_2 - c_1} \cdot 
      \frac{(1 - c_1)e^{\varepsilon^\star_g} + c_2 - 1}
           {e^{\varepsilon^\star_g} + \left(\frac{c_2 - 1}{1 - c_1}\right)\left(1 + g^*(-e^{\varepsilon^\star_g})\right)}.
\]
Moreover, we have
\[
\lambda_{c_1,c_2,g}^\star 
\;=\;
\frac{\,e^{\varepsilon^\star_g} \;+\; \tfrac{c_2 - c_1}{1 - c_1}\,\Bigl(1 + g^*\!\bigl(-\,e^{\varepsilon^\star_g}\bigr)\Bigr) \;-\; 1}
     {\,(1 - c_1)\,e^{\varepsilon^\star_g} + \bigl(c_2 - 1\bigr)}.
\]
Let $\theta \coloneqq e^{\varepsilon_g^\star},
\quad
\delta(\varepsilon_g^\star) \coloneqq 1 + g^{\!*}\!\bigl(-\theta\bigr)$, we have
\begin{align}\label{eqn: lambda_g}
&\lambda_{c_1,c_2,g}^\star
      = \frac{\theta
              + \dfrac{c_2-c_1}{1-c_1}\,\delta(\varepsilon_g^\star)
              - 1}
             {(1-c_1)\theta + (c_2-1)}.
\end{align}

Therefore,
\[
r_1
  = \frac{c_1}{c_2-c_1}
    \frac{(1-c_1)\theta + c_2 - 1}{1-\delta(\varepsilon_g^\star)}.
\]
From \eqref{eqn: lambda_g},
\begin{align*}
1-(1-c_1)\lambda_{c_1,c_2,g}^\star
  &=\frac{(c_2-c_1)\bigl[1-\delta(\varepsilon_g^\star)\bigr]}
          {(1-c_1)\theta + (c_2-1)} .
\end{align*}
Hence
\[
r_1
 = \frac{c_1}{1-(1-c_1)\lambda_{c_1,c_2,g}^\star} .
\]

Similarly,
\[
r_2
  = \frac{c_2}{c_2-c_1}\,
    \frac{(1-c_1)\theta + c_2 - 1}
         {\theta + \dfrac{c_2-1}{1-c_1}\,\delta(\varepsilon_g^\star)} .
\]
Using \eqref{eqn: lambda_g} again,
\begin{align*}
(c_2-1)\lambda_{c_1,c_2,g}^\star + 1
  &= \frac{c_2-c_1}{(1-c_1)\theta + (c_2-1)}
     \Bigl[\theta + \tfrac{c_2-1}{1-c_1}\,\delta(\varepsilon_g^\star)\Bigr].
\end{align*}
Thus
\[
r_2
 = \frac{c_2}{(c_2-1)\lambda_{c_1,c_2,g}^\star + 1 }.
\]
In conclusion,
\[\mathcal{R}\Bigl(\mathcal{Q}_{\mathcal{X}, \tilde{\mathcal{P}}, \varepsilon^\star_g, \delta(\varepsilon^\star_g)}, 
                    \tilde{\mathcal{P}}, f\Bigr) = \mathcal{R}\big(\mathcal{Q}_{\mathcal{X}, \tilde{\mathcal{P}}, g}, \tilde{\mathcal{P}}, f\big) = \frac{1 - r_1}{r_2 - r_1} f(r_2) + \frac{r_2 - 1}{r_2 - r_1} f(r_1),
\]

for\[r_1=\frac{c_1}{1-(1-c_1)\lambda_{c_1,c_2,g}^\star} \quad \text{and } \quad r_2=\frac{c_2}{(c_2-1)\lambda_{c_1,c_2,g}^\star+1}.\]
\end{proof}

\subsection{Proof of Theorem \ref{thm: continuous global functional}}

\begin{proof}
  
Based on Theorem~\ref{thm: global functional}, we define \( \tilde{\mathcal{P}} = \tilde{\mathcal{P}}_{c_1, c_2, \mu} \) under Assumption~\ref{assumption: norm}. Moreover, suppose \( \mu \) is \( \left(\alpha, \frac{1}{\alpha}, 1\right) \)-decomposable with \( \alpha = \frac{1 - c_1}{c_2 - c_1} \). Then, the sampler defined as
\[
\bQ^\star_{c_1,c_2,\mu,g}(P) = \lambda^\star_{c_1,c_2,g} P + \left(1 - \lambda^\star_{c_1,c_2,g}\right) \mu
\]
satisfies $g$-FLDP and is minimax-optimal with respect to any $f$-divergence, where \( \lambda_{c_1,c_2,g}^\star \) is defined as
\[
\lambda^\star_{c_1,c_2,g} = \inf_{\beta \geq 0} 
\tfrac{e^\beta + \frac{c_2 - c_1}{1 - c_1} ( 1 + g^*(-e^\beta)) - 1}
     {(1 - c_1)e^\beta + c_2 - 1}.
\]

Theorem~\ref{thm: continuous global functional} is a special case of Theorem~\ref{thm: global functional}. Below, we demonstrate this reduction explicitly.

Recall that in the setup of Theorem~\ref{thm: global functional}, for a general sample space \( \mathcal{X} \), the universe \( \tilde{\mathcal{P}} \) is defined as
\[
\tilde{\mathcal{P}}_{c_1, c_2, \mu}
:= \left\{P\in \mathcal{P}(\mathcal{X}) : P \ll \mu,\ c_1 \leq \frac{dP}{d\mu} \leq c_2,\ \mu\text{-a.e.} \right\}.
\]

In the setup of Theorem~\ref{thm: continuous global functional}, \( \mathcal{X} = \mathbb{R}^n \), and \( \tilde{\mathcal{P}} \) is defined as
\[
\tilde{\mathcal{P}}_{c_1, c_2, h} := \left\{ P \in \mathcal{C}(\mathbb{R}^n) : c_1 h(x) \leq p(x) \leq c_2 h(x),\quad \forall x \in \mathbb{R}^n \right\}.
\]

% Hence, it suffices to show that for \( \mathcal{X} = \mathbb{R}^n \) and \( \mu = h \), we have \( \tilde{\mathcal{P}}_{c_1, c_2, \mu} = \tilde{\mathcal{P}}_{c_1, c_2, h} \), and \( \mu \) is \( \left(\alpha, \frac{1}{\alpha}, 1\right) \)-decomposable for \( \alpha = \frac{1 - c_1}{c_2 - c_1} \).

Let \( \lambda \) denote the Lebesgue measure on \( \mathbb{R}^n \). We have \( \mathcal{X} = \mathbb{R}^n \), and \( \tilde{\mathcal{P}} = \tilde{\mathcal{P}}_{c_1, c_2, h} = \tilde{\mathcal{P}}_{c_1, c_2, \mu} \), where \( \mu \ll \lambda \) and \( \frac{d\mu}{d\lambda} = h \). Therefore, \( \bQ^\star_{c_1,c_2,\mu,g} = \bQ^\star_{c_1,c_2,h,g} \).

This is because for each \( P \in \tilde{\mathcal{P}} \) with corresponding PDF \( p(x) \), the chain rule of the Radon-Nikodym derivative gives:
\[
p(x) = \frac{dP}{d\lambda} = \frac{dP}{d\mu}(x) \cdot \frac{d\mu}{d\lambda}(x) = \frac{dP}{d\mu}(x) h(x).
\]

It remains to show that \( \mu \) is \( \left(\alpha, \frac{1}{\alpha}, 1\right) \)-decomposable for \( \alpha = \frac{1 - c_1}{c_2 - c_1} \). From Assumption~\ref{assumption: norm}, we know that \( t = \frac{c_2 - c_1}{1 - c_1} \in \mathbb{N} \).  Therefore we need to show that  \( \mu \) is \( \left(\frac{1}{t}, t, 1\right) \)-decomposable.

Since \( \mu \ll \lambda \), the function \( s \mapsto \mu((-\infty, s] \times \mathbb{R}^{n-1}) \) is continuous. As \( s \to -\infty \), this measure tends to \( 0 \), and as \( s \to \infty \), it tends to \( 1 \). By the intermediate value theorem, for each \( i \in [t] \), there exists a threshold \( s_i \in \mathbb{R} \) such that
\[
\mu((-\infty, s_i] \times \mathbb{R}^{n-1}) = \alpha i.
\]
We then define the sets \( A_1 = (-\infty, s_1] \times \mathbb{R}^{n-1} \), and for \( i \geq 2 \), set \( A_i = (s_{i-1}, s_i] \times \mathbb{R}^{n-1} \). These sets satisfy the requirements of \( \left(\frac{1}{t}, t, 1\right) \)-decomposability. Therefore, \( \mu \) is indeed \( \left(\frac{1}{t}, t, 1\right) \)-decomposable, completing the proof.

Therefore, Theorem~\ref{thm: global functional} contains Theorem~\ref{thm: continuous global functional} as a special case. Under Assumption~\ref{assumption: norm}, the sampler \( \mathbf{Q}_{c_1,c_2,h,g}^\star \), defined as a continuous distribution whose density is given by
\begin{equation}
   q^\star_g(P)(x) \coloneqq \lambda^\star_{c_1,c_2,g} p(x) + \left(1 - \lambda^\star_{c_1,c_2,g}\right) h(x), \quad \lambda^\star_{c_1,c_2,g} = \inf_{\beta \geq 0} 
\tfrac{e^\beta + \frac{c_2 - c_1}{1 - c_1} ( 1 + g^*(-e^\beta)) - 1}
     {(1 - c_1)e^\beta + c_2 - 1},
\end{equation}
satisfies \( g \)-FLDP and is minimax-optimal under any \( f \)-divergence. That is,
\begin{equation*}
    \sup_{P \in \tilde{\mathcal{P}}} D_f\big(P \,\|\, \bQ^\star_{c_1,c_2,h,g}(P)\big)
    = \mathcal{R}\big(\mathcal{Q}_{\mathbb{R}^n, \tilde{\mathcal{P}}, g}, \tilde{\mathcal{P}}, f\big)
    = \frac{1 - r_1}{r_2 - r_1} f(r_2) + \frac{r_2 - 1}{r_2 - r_1} f(r_1),
\end{equation*}
for $r_1=\frac{c_1}{1-(1-c_1)\lambda_{c_1,c_2,g}^\star}$ and $r_2=\frac{c_2}{(c_2-1)\lambda_{c_1,c_2,g}^\star+1}$.
\end{proof}
\subsection{Proof of Theorem \ref{thm: discrete global functional}}
\begin{proof}
    From Theorem \ref{thm: global functional}, we know that if $\tilde{\mathcal{P}} = \tilde{\mathcal{P}}_{c_1, c_2, \mu}$ is defined under Assumption \ref{assumption: norm} and  $\mu$ is $(\alpha,\frac{1}{\alpha},1)$-decomposable with $\alpha = \frac{1 - c_1}{c_2 - c_1}$, then the sampler
    $\bQ^\star_{c_1,c_2,\mu,g}(P) \coloneqq \lambda^\star_{c_1,c_2,g} P + (1 - \lambda^\star_{c_1,c_2,g})\mu$
    belongs to the class $\mathcal{Q}_{\mathcal{X}, \tilde{\mathcal{P}}, g}$ and is minimax-optimal under any $f$-divergence $D_f$; that is,
    \begin{equation*}
        \sup_{P \in \tilde{\mathcal{P}}} D_f\big(P \, \| \, \bQ^\star_{c_1,c_2,\mu,g}(P)\big)
        = \mathcal{R}\big(\mathcal{Q}_{\mathcal{X}, \tilde{\mathcal{P}}, g}, \tilde{\mathcal{P}}, f\big),
    \end{equation*}
In this case, 
    \[
    \lambda^\star_{c_1,c_2,g} 
    = \inf_{\beta \geq 0} 
    \frac{e^\beta + \frac{c_2 - c_1}{1 - c_1} \big( 1 + g^*(-e^\beta) \big) - 1}
         {(1 - c_1)e^\beta + c_2 - 1}.
    \]


 It can be shown that:
\[\mathcal{P}([k]) = \tilde{\mathcal{P}}_{0, k, \mu_k}
:= \left\{ P \in \mathcal{P}([k]) : \quad P \ll \mu_k,\quad 0 \leq \frac{dP}{d\mu_k} \leq k \quad \mu_k \text{-a.e.} \right\}
\]
Consider the following setting:  
\[
\tilde{\mathcal{P}} = \mathcal{P}([k]), \quad \mathcal{X} = [k], \quad c_1 = 0, \quad c_2 = k, \quad \mu = \mu_k,
\]
where \( \mu_k \) is the uniform distribution on \([k]\). In this case, we have \( \frac{c_2 - c_1}{1 - c_1} = k \in \mathbb{N} \), and \( \mu_k \) is \( \left(\alpha, \frac{1}{\alpha}, 1 \right) \)-decomposable with \( \alpha = \frac{1}{k} \). 
Moreover, for this specific instantiation, it can be shown that for all \( \varepsilon \geq 0 \),
\[
1 + g^*(-e^\varepsilon) \leq \frac{(c_2 - c_1 e^\varepsilon)(1 - c_1)}{c_2 - c_1} 
= \frac{(k)(1)}{k} = 1.
\]
This comes from the definition of the convex conjugate and the properties of the function \( g \):
\[
\forall \varepsilon \geq 0: \quad -1 \leq g^*(-e^\varepsilon) \leq 0,
\]
and hence,
\[
1 + g^*(-e^\varepsilon) \leq 1.
\]
Therefore, all conditions of Theorem \ref{thm: global functional} are satisfied and  Theorem \ref{thm: discrete global functional} is a special case of the more general case Theorem \ref{thm: global functional}. Hence, the sampler $
    \bQ^\star_{k,g}(P) \coloneqq \lambda^\star_{k,g} P + (1 - \lambda^\star_{k,g})\mu_k$
    belongs to the class $\mathcal{Q}_{[k], \tilde{\mathcal{P}}, g}$ and is minimax-optimal under any $f$-divergence $D_f$; that is,
    \begin{equation*}
        \sup_{P \in \tilde{\mathcal{P}}} D_f\big(P \, \| \, \bQ^\star_{k,g}(P)\big)
        = \mathcal{R}\big(\mathcal{Q}_{[k], \tilde{\mathcal{P}}, g}, \tilde{\mathcal{P}}, f\big),
    \end{equation*}
    where $\lambda^\star_{k,g}$ is the optimal solution to the following optimization problem:
    \[
    \lambda^\star_{k,g} 
    = \inf_{\beta \geq 0} 
    \frac{e^\beta + k \big( 1 + g^*(-e^\beta) \big) - 1}
         {e^\beta + k - 1}.
    \]
\end{proof}

\subsection{Proof of Corollary \ref{cor: special case of f-LDP}}\label{appendix: proof of g_eps}
\begin{proof}

Corollary~\ref{cor: special case of f-LDP} follows immediately from Corollary~\ref{cor: special case of f-LDP approximate} by setting \(\delta = 0\). With this choice, we have \(g_\varepsilon = g_{\varepsilon,0}\), thus the proof is identical.  Moreover, as in the proof of Theorem~\ref{thm: continuous global functional}, the continuous case of Corollary~\ref{cor: special case of f-LDP} can be shown to be a special instance of the more general result stated in Corollary~\ref{cor: special case of f-LDP approximate}. For completeness, a brief proof sketch is provided below.


\textbf{Step 1: Find $g^*_{\eps}$}
\begin{equation}
    g^*_\varepsilon(y) =
    \begin{cases}
        -1, & \text{if } y < -e^{\varepsilon}, \\[6pt]
        \dfrac{y - 1}{e^{\varepsilon} + 1}, & \text{if } -e^{\varepsilon} \leq y \leq -e^{-\varepsilon}, \\[6pt]
        y, & \text{if } y > -e^{-\varepsilon}.
    \end{cases}
\end{equation}




\textbf{Step 2: Solve the infimum to obtain $\lambda^\star_{c_1,c_2,g_\eps}$}

\begin{equation}
    \lambda^\star_{c_1,c_2,g_\eps} =
    \begin{cases}
        \frac{e^\eps - 1}{(e^\eps + 1)(1 - c_1)}, & \text{if } c_1 + c_2 < 2, \\[6pt]
        \frac{e^\eps - 1}{(1 - c_1)e^\eps + c_2 - 1}, & \text{if } c_1 + c_2 \geq 2.
    \end{cases}
\end{equation}

Since $\frac{c_2 - c_1}{1 - c_1} \in \mathbb{N}$ and $c_2 > 1$, it follows that $\frac{c_2 - c_1}{1 - c_1} \geq 2$, which in turn implies  $c_1 + c_2 \geq 2$. Hence:

\[\lambda^\star_{c_1,c_2,g_\eps} = \frac{e^\eps - 1}{(1 - c_1)e^\eps + c_2 - 1}.\]

\textbf{Step 3: Optimality proof}



Following Theorem \ref{thm: global functional}, for the obtained value of $\lambda_{c_1,c_2,g_{\eps}}$, 
the sampler \begin{equation*}
    \bQ^\star_{c_1,c_2,\mu,g_{\varepsilon}}(P) \coloneqq \lambda^\star_{c_1,c_2,g_{\varepsilon}} P + (1 - \lambda^\star_{c_1,c_2,g_{\varepsilon}})\mu
\end{equation*}
belongs to the class $\mathcal{Q}_{\mathcal{X}, \tilde{\mathcal{P}}, g_{\varepsilon,\delta}}$ and is minimax-optimal with respect to any $f$-divergence $D_f$, that is,
\begin{equation*}
    \sup_{P \in \tilde{\mathcal{P}}} D_f\big(P \, \| \, \bQ^\star_{c_1,c_2,\mu,g_{\varepsilon}}(P)\big)
    = \mathcal{R}\big(\mathcal{Q}_{\mathcal{X}, \tilde{\mathcal{P}}, g_{\varepsilon}}, \tilde{\mathcal{P}}, f\big).
\end{equation*}
Therefore, we have:
\begin{equation*}
    \sup_{P \in \tilde{\mathcal{P}}} D_f\big(P \,\|\, \bQ^\star_{c_1,c_2,h,g_\eps}(P)\big)
    = \mathcal{R}\big(\mathcal{Q}_{\mathbb{R}^n, \tilde{\mathcal{P}}, g_\eps}, \tilde{\mathcal{P}}, f\big)
    = \frac{1 - r_1}{r_2 - r_1} f(r_2) + \frac{r_2 - 1}{r_2 - r_1} f(r_1),
\end{equation*}
for $
r_1 = c_1 \cdot \frac{(1 - c_1)e^{\eps} + c_2 - 1}{c_2 - c_1 }$, and $
r_2 = \frac{c_2}{c_2 - c_1} \cdot \frac{(1 - c_1)e^{\eps} + c_2 - 1}{e^{\eps}}$.

The optimal value of 
\(\mathcal{R}\bigl(\mathcal{Q}_{\mathbb{R}^n,\tilde{\mathcal{P}},g_\varepsilon},
                    \tilde{\mathcal{P}},f\bigr)\)
is obtained directly from
Theorem~\ref{thm: continuous global functional}
by substituting
\(
\lambda_{c_1,c_2,g_\varepsilon}^{\star}
      =\dfrac{e^{\varepsilon}-1}{(1-c_1)e^{\varepsilon}+c_2-1}.
\)

\end{proof}




\subsection{Proof of Corollary \ref{cor: special case GDP}}\label{proof: GDP}

\begin{proof}
    From Theorem \ref{thm: global functional} we know that for
    \[
    \lambda^\star_{G_\nu} 
    = \inf_{\beta \geq 0} 
    \frac{e^\beta + \frac{c_2 - c_1}{1 - c_1} \big( 1 + G_\nu^*(-e^\beta) \big) - 1}
         {(1 - c_1)e^\beta + c_2 - 1}.
    \]
   the sampler  \eqref{eqn: continuous global functional opt mech} for $g=G_\nu$ belongs to $\mathcal{Q}_{\R^n, \tilde{\mathcal{P}}, G_\nu}$ and is minimax-optimal with respect to any $f$-divergence, that is,
\begin{equation*}
    \sup_{P \in \tilde{\mathcal{P}}} D_f\big(P \, \| \, \bQ^\star_{G_\nu}(P)\big)
    = \mathcal{R}\big(\mathcal{Q}_{\R^n, \tilde{\mathcal{P}},G_\nu}, \tilde{\mathcal{P}}, f\big).
\end{equation*}
    Now, we compute $G_\nu^*(-e^\beta)$. It is shown in Corollary 1 of \citet{dong2022gaussian} that 
\[
G_\nu^*(y)
  = y\,\Phi\Bigl(-\frac{\nu}{2}-\frac{1}{\nu}\log(-y)\Bigr)\;-\;
  \Phi\Bigl(-\frac{\nu}{2}+\frac{1}{\nu}\log(-y)\Bigr).
\]
When $y = -e^\beta$,
\[G^*_{\nu}\bigl(-e^{\beta}\bigr)
  \;=\;
  -\,e^{\beta}\,
    \Phi\Bigl(-\frac{\nu}{2}-\frac{\beta}{\nu}\Bigr)
  \;-\;
    \Phi\Bigl(-\frac{\nu}{2}+\frac{\beta}{\nu}\Bigr)\]
Therefore, we have:
\begin{align*}
\lambda^\star_{G_\nu} & =\inf_{\beta \geq 0} 
    \frac{e^\beta + \Big(\frac{c_2 - c_1}{1 - c_1}\Big) \bigg( 1 -e^\beta \Phi\left(-\frac{\nu}{2} - \frac{\beta}{\nu}\right) - \Phi\left(-\frac{\nu}{2} + \frac{\beta}{\nu}\right) \bigg) - 1}
         {(1 - c_1)e^\beta + c_2 - 1}\\
         & = \inf_{\beta \geq 0} 
    \frac{e^\beta + \Big(\frac{c_2 - c_1}{1 - c_1}\Big) \bigg( \Phi\left(\frac{\nu}{2} - \frac{\beta}{\nu}\right)  -e^\beta \Phi\left(-\frac{\nu}{2} - \frac{\beta}{\nu}\right)  \bigg) - 1}
         {(1 - c_1)e^\beta + c_2 - 1}.    
\end{align*}

\end{proof}

\subsection{Proof of Theorem \ref{thm: local functional}}
% We first assume $\eps \le 2\log(\gamma)$ and provide the proof of the optimal sampler and optimal value. Then in the end we also consider the case that $2 \log(\gamma) > \eps$.
 
\begin{proof}
    \textbf{Step 1: Lower bound}

In the first step of the proof, we establish a lower bound for the local minimax objective function 
\(\mathcal{R}\big(\mathcal{Q}_{\mathcal{X},\tilde{\mathcal{P}},g},N_\gamma(P_0),f\big)\) as follows:
\begin{align} \label{eqn: lower bound local functional}
\mathcal{R}\big(\mathcal{Q}_{\mathcal{X},\tilde{\mathcal{P}},g},N_\gamma(P_0),f\big) 
&= \inf_{\bQ \in \mathcal{Q}_{\mathcal{X},\tilde{\mathcal{P}},g}} \, 
\sup_{P \in N_\gamma(P_0)} D_f\big(P \, \| \, \bQ(P)\big) \nonumber \\
&\geq \inf_{\bQ \in \mathcal{Q}_{\mathcal{X},N_\gamma(P_0),g}} \, 
\sup_{P \in N_\gamma(P_0)} D_f\big(P \, \| \, \bQ(P)\big) \nonumber \\
&= \mathcal{R}\big(\mathcal{Q}_{\mathcal{X},N_\gamma(P_0),g},N_\gamma(P_0),f\big).
\end{align}

The inequality in the second line follows from the fact that any sampler 
\(\bQ : \tilde{\mathcal{P}} \to \mathcal{P}(\mathcal{X})\) in 
\(\mathcal{Q}_{\mathcal{X}, \tilde{\mathcal{P}}, g}\) also belongs to 
\(\mathcal{Q}_{\mathcal{X}, N_\gamma(P_0), g}\). 
The third line is by definition of the global minimax risk in~\eqref{eqn: global minimax functional problem} when the universe $\mathcal{\tilde{P}}$ is restricted to \(N_\gamma(P_0)\).

\textbf{Step 2: Upper bound}

We now construct a sampler that achieves the lower bound. Specifically, we aim to find a sampler 
\(\bQ^\star_{g, N_\gamma(P_0)} \in \mathcal{Q}_{\mathcal{X}, \tilde{\mathcal{P}}, g}\) such that
\[
\sup_{P \in N_\gamma(P_0)} D_f\big(P \, \| \, \bQ^\star_{g, N_\gamma(P_0)}(P)\big) 
= \mathcal{R}\big(\mathcal{Q}_{\mathcal{X}, N_\gamma(P_0), g}, N_\gamma(P_0), f\big).
\]

Observe that the neighborhood \(N_\gamma(P_0)\) aligns with the class 
\(\tilde{\mathcal{P}} = \tilde{\mathcal{P}}_{c_1, c_2, \mu}\) defined in Theorem~\ref{thm: global functional}:
\[
\tilde{\mathcal{P}}_{c_1, c_2, \mu}
:= \left\{ P \in \mathcal{P}(\mathcal{X}) : \quad P \ll \mu, \quad 
c_1 \leq \frac{dP}{d\mu} \leq c_2 \quad \mu\text{-a.e.} \right\},
\]
\[
N_\gamma(P_0) := \left\{ P \in \mathcal{P}(\mathcal{X}) : \quad 
\frac{1}{\gamma} \leq \frac{dP}{dP_0}(x) \leq \gamma \quad \forall x \in \mathcal{X} \right\}.
\]

By substituting \(c_1 = \tfrac{1}{\gamma}\), \(c_2 = \gamma\), and \(\mu = P_0\), we obtain 
\(\tilde{\mathcal{P}}_{c_1, c_2, \mu} = N_\gamma(P_0)\).

Additionally, since \(\gamma \in \mathbb{N}\), it follows that 
\(\tfrac{c_2 - c_1}{1 - c_1} = \tfrac{\gamma - \frac{1}{\gamma}}{1 - \frac{1}{\gamma}} \in \mathbb{N}\), 
satisfying the integer condition in Theorem~\ref{thm: global functional}. The \((\alpha,\frac{1}{\alpha},1)\)-decomposability 
of \(P_0\) with \(\alpha = \tfrac{1}{\gamma + 1}\) also matches the requirement that $\mu$ is $(\alpha,\frac{1}{\alpha},1)$-decomposable with  
\(\alpha = \tfrac{1 - c_1}{c_2 - c_1}\).

Therefore, all the conditions in Theorem~\ref{thm: global functional} are satisfied. Thus, the optimal 
sampler \(\bQ^\star_g\) corresponding to 
\((c_1, c_2, \mu) = \left(\tfrac{1}{\gamma}, \gamma, P_0\right)\) achieves
\[
\sup_{P \in N_\gamma(P_0)} D_f\big(P \, \| \, \bQ^\star_g(P)\big)
= \mathcal{R}\big(\mathcal{Q}_{\mathcal{X}, N_\gamma(P_0), g}, N_\gamma(P_0), f\big).
\]

Since \(\bQ^\star_g\) is defined only on \(N_\gamma(P_0)\), we extend it to the larger domain 
\(\tilde{\mathcal{P}}\) by defining:
\[
\bQ^\star_{g, N_\gamma(P_0)}(P) := 
\begin{cases}
\bQ^\star_g(P), & \text{if } P \in N_\gamma(P_0), \\[6pt]
\bQ^\star_g(\hat{P}), & \text{otherwise},
\end{cases}
\]
where \( \hat{P} \in N_\gamma(P_0) \) is a distribution that minimizes \(  D_f(P \,\|\, P') \) over all $P' \in N_\gamma(P_0)$.

By construction, \(\bQ^\star_{g, N_\gamma(P_0)} \in \mathcal{Q}_{\mathcal{X}, \tilde{\mathcal{P}}, g}\), 
and for all \(P \in N_\gamma(P_0)\), we have:
\begin{align*}
\sup_{P \in N_\gamma(P_0)} D_f\big(P \, \| \, \bQ^\star_{g, N_\gamma(P_0)}(P)\big) 
&= \sup_{P \in N_\gamma(P_0)} D_f\big(P \, \| \, \bQ^\star_g(P)\big) \\
&= \mathcal{R}\big(\mathcal{Q}_{\mathcal{X}, N_\gamma(P_0), g}, N_\gamma(P_0), f\big).
\end{align*}

Hence, we obtain the upper bound:
\begin{align} \label{eqn: upper bound local functional}
\mathcal{R}\big(\mathcal{Q}_{\mathcal{X}, \tilde{\mathcal{P}}, g}, N_\gamma(P_0), f\big)
&= \inf_{\bQ \in \mathcal{Q}_{\mathcal{X}, \tilde{\mathcal{P}}, g}} \,
\sup_{P \in N_\gamma(P_0)} D_f\big(P \, \| \, \bQ(P)\big) \nonumber \\ 
&\leq \sup_{P \in N_\gamma(P_0)} D_f\big(P \, \| \, \bQ^\star_{g, N_\gamma(P_0)}(P)\big)  \nonumber \\ & = \mathcal{R}\big(\mathcal{Q}_{\mathcal{X}, N_\gamma(P_0), g}, N_\gamma(P_0), f\big).
\end{align}

Combining equations~\eqref{eqn: lower bound local functional} and~\eqref{eqn: upper bound local functional}, 
we conclude that
\[
\mathcal{R}\Big(\mathcal{Q}_{\mathcal{X}, \tilde{\mathcal{P}}, g}, N_\gamma(P_0), f\Big) 
= \mathcal{R}\Big(\mathcal{Q}_{\mathcal{X}, N_\gamma(P_0), g}, N_\gamma(P_0), f\Big).
\]
Note that this proof is stated for a general space \(\mathcal{X}\) and a general universe \(\tilde{\mathcal{P}} = \Ptm\). It extends directly to the continuous case with \(\mathcal{X} = \mathbb{R}^n\) and \(\tilde{\mathcal{P}} = \Pth\), following the same reduction argument used in the proof of Theorem~\ref{thm: continuous global functional}.
% If $2 \log(\gamma) > \eps$, then for any $P_1, P_2 \in N_{\gamma}(P_0)$, we have 
% $p_1(x)/p_2(x) \leq \frac{\gamma}{\frac{1}{\gamma}} \leq e^{\varepsilon}$, hence we can easily observe that the sampler 
% $\mathbf{Q}_f$ defined as $\mathbf{Q}_f(P) = P$ for all $P \in N_{\gamma}(P_0)$ satisfies 
% $\eps$-LDP and  $\mathcal{R}\big(\mathcal{Q}_{\mathcal{X},N_\gamma(P_0),g},N_\gamma(P_0),f\big) = 0$. Since \(\bQ^I_f\) is defined only on \(N_\gamma(P_0)\), we extend it to the larger domain 
% \(\tilde{\mathcal{P}}\) by defining:
% \[
% \bQ^I_{f, N_\gamma(P_0)}(P) := 
% \begin{cases}
% \bQ^I_f(P), & \text{if } P \in N_\gamma(P_0), \\[6pt]
% \bQ^I_f(\hat{P}), & \text{otherwise},
% \end{cases}
% \]
% where \(\hat{P} \in N_\gamma(P_0)\) minimizes 
% \(\inf_{P' \in N_\gamma(P_0)} D_f(P \, \| \, P')\). With similar argument as the case of $2 \log(\gamma) \leq \eps$, we can conclude $\mathcal{R}\big(\mathcal{Q}_{\mathcal{X},\tilde{\mathcal{P}},g},N_\gamma(P_0),f\big) = 0$.
\end{proof}


\subsection{Proof of Theorem \ref{thm: local pure}} \label{proof: local pure}

\begin{proof}

let the global minimax formulation under $\eps$-LDP be defined as: 
\begin{equation}\label{eqn: global minimax pure problem}\mathcal{R}\big(\mathcal{Q}_{\mathcal{X},\tilde{\mathcal{P}},\eps},\tilde{\mathcal{P}},f\big) \coloneqq \inf_{\bQ \in \mathcal{Q}_{\mathcal{X},\tilde{\mathcal{P}},\eps}} \hspace{0.1cm} \sup_{P \in \tilde{\mathcal{P}}}\hspace{0.1cm} D_f\big(P \, \| \, \bQ(P)\big)
\end{equation}
Under certain assumptions, formally defined in Section~\ref{sec: global functional}, \citet{park2024exactly} derives the optimal sampler and optimal value for the optimization problem~\eqref{eqn: global minimax pure problem}.

\textbf{Step 1: Lower bound} 

In the first step of the proof, we establish a lower bound for the local minimax objective function 
\(\mathcal{R}\big(\mathcal{Q}_{\mathcal{X},\tilde{\mathcal{P}},\eps},N_\gamma(P_0),f\big)\) as follows:

\begin{align} \label{eqn: lower bound local pure}
\mathcal{R}\big(\mathcal{Q}_{\mathcal{X},\tilde{\mathcal{P}},\eps},N_\gamma(P_0),f\big) 
&= \inf_{\bQ \in \mathcal{Q}_{\mathcal{X},\tilde{\mathcal{P}},\eps}} \, 
\sup_{P \in N_\gamma(P_0)} D_f\big(P \, \| \, \bQ(P)\big) \nonumber \\
&\geq \inf_{\bQ \in \mathcal{Q}_{\mathcal{X},N_\gamma(P_0),\eps}} \, 
\sup_{P \in N_\gamma(P_0)} D_f\big(P \, \| \, \bQ(P)\big) \nonumber \\
&= \mathcal{R}\big(\mathcal{Q}_{\mathcal{X},N_\gamma(P_0),\eps},N_\gamma(P_0),f\big).
\end{align}

The inequality in the second line follows from the fact that any sampler 
\(\bQ : \tilde{\mathcal{P}} \to \mathcal{P}(\mathcal{X})\) in 
\(\mathcal{Q}_{\mathcal{X}, \tilde{\mathcal{P}}, \eps}\) also belongs to 
\(\mathcal{Q}_{\mathcal{X}, N_\gamma(P_0), \eps}\). 
The third line is by definition of the global minimax risk in~\eqref{eqn: global minimax pure problem} when the universe $\mathcal{\tilde{P}}$ is restricted to \(N_\gamma(P_0)\).

\textbf{Step 2: Upper bound} 

We now construct a sampler that achieves the lower bound. Specifically, we aim to find a sampler 
\(\bQ^\star_{\eps, N_\gamma(P_0)} \in \mathcal{Q}_{\mathcal{X}, \tilde{\mathcal{P}}, \eps}\) such that
\[
\sup_{P \in N_\gamma(P_0)} D_f\big(P \, \| \, \bQ^\star_{\eps, N_\gamma(P_0)}(P)\big) 
= \mathcal{R}\big(\mathcal{Q}_{\mathcal{X}, N_\gamma(P_0), \eps}, N_\gamma(P_0), f\big).
\]

% Observe that the neighborhood \(N_\gamma(P_0)\) aligns with the class 
% \(\tilde{\mathcal{P}} = \tilde{\mathcal{P}}_{c_1, c_2, \mu}\) defined in Theorem~\ref{thm: global functional}:
% \[
% \tilde{\mathcal{P}}_{c_1, c_2, \mu}
% := \left\{ P \in \mathcal{P}(\mathcal{X}) : \quad P \ll \mu, \quad 
% c_1 \leq \frac{dP}{d\mu} \leq c_2 \quad \mu\text{-a.e.} \right\},
% \]
% \[
% N_\gamma(P_0) := \left\{ P \in \mathcal{P}(\mathcal{X}) : \quad 
% \frac{1}{\gamma} \leq \frac{dP}{dP_0}(x) \leq \gamma \quad \forall x \in \mathcal{X} \right\}.
% \]

% By substituting \(c_1 = \tfrac{1}{\gamma}\), \(c_2 = \gamma\), and \(\mu = P_0\), we obtain 
% \(\tilde{\mathcal{P}}_{c_1, c_2, \mu} = N_\gamma(P_0)\).

% Additionally, since \(\gamma \in \mathbb{N}\), it follows that 
% \(\tfrac{c_2 - c_1}{1 - c_1} = \tfrac{\gamma - \frac{1}{\gamma}}{1 - \frac{1}{\gamma}} \in \mathbb{N}\), 
% satisfying the integer condition in Theorem~\ref{thm: global functional}. The \((\alpha,\frac{1}{\alpha},1)\)-decomposability 
% of \(P_0\) with \(\alpha = \tfrac{1}{\gamma + 1}\) also matches the requirement that $\mu$ is $(\alpha,\frac{1}{\alpha},1)$-decomposable with  
% \(\alpha = \tfrac{1 - c_1}{c_2 - c_1}\).

% Therefore, all the conditions in Theorem~\ref{thm: global functional} are satisfied. Thus, the optimal 
Let sampler \(\bQ^\star_\eps: \ngb \to \mathcal{P}(\mathcal{X})\) be an optimal sampler for the global minimax problem (\ref{eqn: global minimax pure problem}) where the universe $\mathcal{\tilde{P}}$ is restricted to \(N_\gamma(P_0)\). 
In other words, $\bQ^\star_\eps$ achieves
\[
\sup_{P \in N_\gamma(P_0)} D_f\big(P \, \| \, \bQ^\star_\eps(P)\big)
= \mathcal{R}\Big(\mathcal{Q}_{\mathcal{X}, N_\gamma(P_0), \eps}, N_\gamma(P_0), f\Big).
\]

Since \(\bQ^\star_\eps\) is defined only on \(N_\gamma(P_0)\), we extend it to the larger domain 
\(\tilde{\mathcal{P}}\) by defining:
\begin{equation}\label{eqn: local pure optimal mech}
    \bQ^\star_{\eps, N_\gamma(P_0)}(P) := 
\begin{cases}
\bQ^\star_\eps(P), & \text{if } P \in N_\gamma(P_0), \\[6pt]
\bQ^\star_\eps(\hat{P}), & \text{otherwise},
\end{cases}    
\end{equation}


where \( \hat{P} \in N_\gamma(P_0) \) is a distribution that minimizes \(  D_f(P \,\|\, P') \) over all $P' \in N_\gamma(P_0)$.

By construction, \(\bQ^\star_{\eps, N_\gamma(P_0)} \in \mathcal{Q}_{\mathcal{X}, \tilde{\mathcal{P}}, \eps}\), 
and for all \(P \in N_\gamma(P_0)\), we have:
\begin{align*}
\sup_{P \in N_\gamma(P_0)} D_f\big(P \, \| \, \bQ^\star_{\eps, N_\gamma(P_0)}(P)\big) 
&= \sup_{P \in N_\gamma(P_0)} D_f\big(P \, \| \, \bQ^\star_\eps(P)\big) \\
&= \mathcal{R}\big(\mathcal{Q}_{\mathcal{X}, N_\gamma(P_0), \eps}, N_\gamma(P_0), f\big).
\end{align*}

Hence, we obtain the upper bound:
\begin{align} \label{eqn: upper bound local pure}
\mathcal{R}\big(\mathcal{Q}_{\mathcal{X}, \tilde{\mathcal{P}}, \eps}, N_\gamma(P_0), f\big)
&= \inf_{\bQ \in \mathcal{Q}_{\mathcal{X}, \tilde{\mathcal{P}}, \eps}} \,
\sup_{P \in N_\gamma(P_0)} D_f\big(P \, \| \, \bQ(P)\big) \nonumber \\ 
&\leq \sup_{P \in N_\gamma(P_0)} D_f\big(P \, \| \, \bQ^\star_{\eps, N_\gamma(P_0)}(P)\big)  \nonumber \\ & = \mathcal{R}\big(\mathcal{Q}_{\mathcal{X}, N_\gamma(P_0), \eps}, N_\gamma(P_0), f\big).
\end{align}

Combining equations~\eqref{eqn: lower bound local pure} and~\eqref{eqn: upper bound local pure}, 
we conclude that
\[
\mathcal{R}\Big(\mathcal{Q}_{\mathcal{X}, \tilde{\mathcal{P}}, \eps}, N_\gamma(P_0), f\Big) 
= \mathcal{R}\Big(\mathcal{Q}_{\mathcal{X}, N_\gamma(P_0), \eps}, N_\gamma(P_0), f\Big).
\] 

\textbf{Step 3: Local minimax risk in a closed form} 


Under the same non-triviality assumption as \citet{park2024exactly}, for the pure LDP case we assume \(\varepsilon < 2\log(\gamma)\); otherwise, the identity sampler becomes the trivial local minimax-optimal sampler with zero minimax risk.


Consider the case where $2 \log(\gamma) \le \eps$. In this case, for any $P_1, P_2 \in N_{\gamma}(P_0)$, we have 
$p_1(x)/p_2(x) \leq \frac{\gamma}{\frac{1}{\gamma}} < e^{\varepsilon}$, hence we can easily observe that the sampler 
$\mathbf{Q}^I_\eps$ defined as $\mathbf{Q}^I_\eps(P) = P$ for all $P \in N_{\gamma}(P_0)$ satisfies 
$\eps$-LDP and  $\mathcal{R}\big(\mathcal{Q}_{\mathcal{X},N_\gamma(P_0),\eps},N_\gamma(P_0),f\big) = 0$. Since \(\bQ^I_\eps\) is defined only on \(N_\gamma(P_0)\), we extend it to the larger domain 
\(\tilde{\mathcal{P}}\) by defining:
\[
\bQ^I_{\eps, N_\gamma(P_0)}(P) := 
\begin{cases}
\bQ^I_\eps(P), & \text{if } P \in N_\gamma(P_0), \\[6pt]
\bQ^I_\eps(\hat{P}), & \text{otherwise},
\end{cases}
\]
where \( \hat{P} \in N_\gamma(P_0) \) is a distribution that minimizes \(  D_f(P \,\|\, P') \) over all $P' \in N_\gamma(P_0)$.

We can conclude $\mathcal{R}\big(\mathcal{Q}_{\mathcal{X},\tilde{\mathcal{P}},\eps},N_\gamma(P_0),f\big) = 0$. Therefore, to exclude trivial samplers we assume $\eps < 2\log(\gamma)$.

In order to obtain the local minimax risk under $\varepsilon$-LDP, we refer to Theorem C.4 in \citet{park2024exactly}.

\begin{theorem}[Theorem C.4 of \citet{park2024exactly}]\label{thm: global pure}
Let $\tilde{\mathcal{P}} = \tilde{\mathcal{P}}_{c_1, c_2, \mu}$ be a set of distributions such that the normalization condition~\eqref{eqn: normalization condition} holds, and suppose $\mu$ is $(\alpha,\frac{1}{\alpha},1)$-decomposable with $\alpha = \frac{1 - c_1}{c_2 - c_1}$. Define
\[
b \coloneqq \frac{c_2 - c_1}{(e^\varepsilon - 1)(1 - c_1) + c_2 - c_1}, \quad
r_1 \coloneqq \frac{c_1}{b}, \quad
r_2 \coloneqq \frac{c_2}{b e^\varepsilon}.
\]
Then, the optimal value of the problem~\eqref{eqn: global minimax pure problem} is given by
\[
\mathcal{R}\big(\mathcal{Q}_{\mathcal{X}, \tilde{\mathcal{P}}, \varepsilon}, \tilde{\mathcal{P}}, f\big)
= \frac{1 - r_1}{r_2 - r_1} \, f(r_2) + \frac{r_2 - 1}{r_2 - r_1} \, f(r_1).
\]

Furthermore, the sampler $\bQ^\star_{c_1, c_2, \mu, \varepsilon}$ defined below is an optimal solution to problem~\eqref{eqn: global minimax pure problem} under any $f$-divergence $D_f$. For each \(P \in \tilde{\mathcal{P}}\), let \(\bQ^\star_{c_1, c_2, \mu, \varepsilon}(P) = Q\) be a probability measure $Q \ll \mu$, such that
\[
\frac{dQ}{d\mu}(x) 
= 
\mathrm{clip}\left(
\frac{1}{r_P} \frac{dP}{d\mu}(x)
\;;\;
b,
\; b e^\varepsilon
\right),
\]
where \(r_P\) is the normalizing constant ensuring that \(\int \frac{dQ}{d\mu} \, d\mu(x) = 1\).
\end{theorem}

Observe that the neighborhood \(N_\gamma(P_0)\) aligns with the class 
\(\tilde{\mathcal{P}} = \tilde{\mathcal{P}}_{c_1, c_2, \mu}\) defined in Theorem~\ref{thm: global functional}:
\[
\tilde{\mathcal{P}}_{c_1, c_2, \mu}
:= \left\{ P \in \mathcal{P}(\mathcal{X}) : \quad P \ll \mu, \quad 
c_1 \leq \frac{dP}{d\mu} \leq c_2 \quad \mu\text{-a.e.} \right\},
\]
\[
N_\gamma(P_0) := \left\{ P \in \mathcal{P}(\mathcal{X}) : \quad 
\frac{1}{\gamma} \leq \frac{dP}{dP_0}(x) \leq \gamma \quad \forall x \in \mathcal{X} \right\}.
\]

By substituting \(c_1 = \tfrac{1}{\gamma}\), \(c_2 = \gamma\), and \(\mu = P_0\), we obtain 
\(\tilde{\mathcal{P}}_{c_1, c_2, \mu} = N_\gamma(P_0)\).

Additionally, the \((\alpha,\frac{1}{\alpha},1)\)-decomposability 
of \(P_0\) with \(\alpha = \tfrac{1}{\gamma + 1}\) also matches the requirement that $\mu$ is $(\alpha,\frac{1}{\alpha},1)$-decomposable with  
\(\alpha = \tfrac{1 - c_1}{c_2 - c_1}\). We now aim to compute the closed-form expression for
\[
\mathcal{R}\Big(\mathcal{Q}_{\mathcal{X}, N_\gamma(P_0), \varepsilon}, N_\gamma(P_0), f\Big).
\]
By the equivalence \(\tilde{\mathcal{P}}_{c_1, c_2, \mu} = N_\gamma(P_0)\) and Theorem~\ref{thm: global pure} applied with \((c_1, c_2, \mu) = \left(\tfrac{1}{\gamma}, \gamma, P_0\right)\), we have:
\[
\mathcal{R}\Big(\mathcal{Q}_{\mathcal{X}, N_\gamma(P_0), \varepsilon}, N_\gamma(P_0), f\Big)
= \frac{1 - r_1}{r_2 - r_1} \, f(r_2) + \frac{r_2 - 1}{r_2 - r_1} \, f(r_1),
\]
where
\[
b \coloneqq \frac{\gamma - \tfrac{1}{\gamma}}{(e^\varepsilon - 1)\left(1 - \tfrac{1}{\gamma}\right) + \gamma - \tfrac{1}{\gamma}} 
= \frac{\gamma + 1}{\gamma + e^\varepsilon}, \quad
r_1 \coloneqq \frac{1}{\gamma b} 
= \frac{e^\varepsilon + \gamma}{\gamma(\gamma + 1)}, \quad
r_2 \coloneqq \frac{\gamma}{b e^\varepsilon} 
= \frac{\gamma(e^\varepsilon + \gamma)}{e^\varepsilon(\gamma + 1)}.
\]
Since
\[
\mathcal{R}\Big(\mathcal{Q}_{\mathcal{X}, \tilde{\mathcal{P}}, \eps}, N_\gamma(P_0), f\Big) 
= \mathcal{R}\Big(\mathcal{Q}_{\mathcal{X}, N_\gamma(P_0), \eps}, N_\gamma(P_0), f\Big),
\] the optimal minimax risk is obtained.

\medskip

\textbf{Step 4: Optimal sampler} 

From the upper bound proof of the theorem, we know that the sampler 
\(\bQ^\star_{\varepsilon, N_\gamma(P_0)}\), defined in \eqref{eqn: local pure optimal mech}, is an optimal sampler, provided that
\[
\sup_{P \in N_\gamma(P_0)} D_f\big(P \, \| \, \bQ^\star_\varepsilon(P)\big)
= \mathcal{R}\Big(\mathcal{Q}_{\mathcal{X}, N_\gamma(P_0), \varepsilon}, N_\gamma(P_0), f\Big).
\]
Therefore, it suffices to construct a sampler 
\(\bQ^\star_\varepsilon : N_\gamma(P_0) \to \mathcal{P}(\mathcal{X})\) that belongs to $\mathcal{Q}_{\mathcal{X}, \ngb, \varepsilon}$
and satisfies the above equality.

By applying Theorem~\ref{thm: global pure} with 
\((c_1, c_2, \mu) = \left(\tfrac{1}{\gamma}, \gamma, P_0\right)\), 
we obtain the following definition for such a sampler. For each \(P \in N_\gamma(P_0)\), define \(\bQ^\star_\varepsilon(P) = Q\), where \(Q\) is a probability measure satisfying $Q \ll P_0$ and given by
\[
\frac{dQ}{dP_0}(x) 
= 
\mathrm{clip}\left(
\frac{1}{r_P} \frac{dP}{dP_0}(x)
\;;\;
\frac{\gamma + 1}{\gamma + e^\varepsilon},
\; \frac{\gamma + 1}{\gamma + e^\varepsilon} e^\varepsilon
\right),
\]
with \(r_P\) denoting the normalizing constant ensuring that 
\(\int \frac{dQ}{dP_0} \, dP_0(x) = 1\).

We extend this sampler to a function over all of \(\mathcal{P}(\mathcal{X})\) by defining
\[
\bQ^\star_{\varepsilon, N_\gamma(P_0)}(P) \coloneqq 
\begin{cases}
\bQ^\star_\varepsilon(P), & \text{if } P \in N_\gamma(P_0), \\[6pt]
\bQ^\star_\varepsilon(\hat{P}), & \text{otherwise},
\end{cases}
\]
where \( \hat{P} \in N_\gamma(P_0) \) is a distribution that minimizes \(  D_f(P \,\|\, P') \) over all $P' \in N_\gamma(P_0)$.

Then \(\bQ^\star_{\varepsilon, N_\gamma(P_0)} \in \mathcal{Q}_{\mathcal{X}, \tilde{\mathcal{P}}, \varepsilon}\), and the following holds:
\begin{equation*}
\sup_{P \in N_\gamma(P_0)} 
D_f\big(P \,\|\, \bQ^\star_{\varepsilon, N_\gamma(P_0)}(P)\big)
= 
\mathcal{R}\Big(\mathcal{Q}_{\mathcal{X}, \tilde{\mathcal{P}}, \varepsilon}, N_\gamma(P_0), f\Big).
\end{equation*}



Note that this proof is stated for a general space \(\mathcal{X}\) and a general universe \(\tilde{\mathcal{P}} = \Ptm\). It extends directly to the continuous case with \(\mathcal{X} = \mathbb{R}^n\) and \(\tilde{\mathcal{P}} = \Pth\), following the same reduction argument used in the proof of Theorem~\ref{thm: continuous global functional}.
\end{proof}


\subsection{Proof of Proposition \ref{prop: pointwise}}\label{appendix: proof pointwise}
\begin{proof}
    To prove this result, we need to define preliminary samplers and the universe set. Let
\[
\tilde{\mathcal{P}}_{c_1, c_2, \mu}
:= \left\{ P \in \mathcal{P}(\mathcal{X}) : P \ll \mu, \quad c_1 \leq \frac{dP}{d\mu} \leq c_2 \quad \mu\text{-a.e.} \right\},
\]
Define a sampler $\mathbf{Q}^\star_{c_1,c_2,\mu,\eps} \in \mathcal{Q}_{\mathcal{X}, \tilde{\mathcal{P}}, \eps}$ as follows:

For each  $P \in \tilde{\mathcal{P}}$, $\mathbf{Q}^\star_{c_1,c_2,\mu,\eps}(P) \coloneqq Q$ is defined as a probability measure such that $Q \ll \mu$  and
\[
\frac{dQ}{d\mu}(x) = \textnormal{clip} \left( \frac{1}{r_P} \frac{dP}{d\mu}(x) ; b, b e^{\eps} \right),
\]
where  $r_P > 0$ is a constant depending on  $P$ so that  $\int \frac{dQ}{d\mu} d\mu(x) = 1$.

\begin{proposition}[Proposition C.8. of \citet{park2024exactly}]\label{prop: pointwise optimal Park et al.}
For any $P \in \Ptm$ and any $f$-divergences $D_f$, we have
    \[
D_f\left(P \,\|\, \mathbf{Q}^\star_{c_1,c_2,\mu,\eps}(P)\right)
= \inf_{\bQ \in \Pt_{b,be^\eps,\mu}} D_f(P \,\|\, \bQ),
\]
for \[
b = \frac{c_2 - c_1}{(e^{\eps} - 1)(1 - c_1) + c_2 - c_1}.
\]
\end{proposition}
We want to prove that if we define the neighborhood as 
\begin{equation*}
N_\gamma(P_0) \coloneqq  \left\{ P \in \mathcal{P}(\mathcal{X}) :\quad E_\gamma(P \, \| \, P_0) =  E_\gamma(P_0 \, \| \, P) = 0  \right\},
\end{equation*}
and let \( \bQ^\star_\varepsilon \) denote the optimal sampler from Theorem~\ref{thm: local pure}, and let \( \bQ^\star_{g_\varepsilon} \) be the instantiation of Theorem~\ref{thm: local functional} with \( g = g_\varepsilon \). Then, for all \( P \in N_\gamma(P_0) \),
\[
D_f(P \,\|\, \bQ^\star_\varepsilon) \leq D_f(P \,\|\, \bQ^\star_{g_\varepsilon}).
\]
Note that this inequality is equivalent to the condition that for all \( P \in N_\gamma(P_0) \),
\[
D_f(P \,\|\, \bQ^\star_{\varepsilon,\ngb}) \leq D_f(P \,\|\, \bQ^\star_{g_\varepsilon,\ngb}),
\]
which is the ultimate result we want to prove in this proposition. 

Therefore, it suffices to show that for $\Ptm = \ngb$, we have $\bQ^\star_\eps = \mathbf{Q}^\star_{c_1,c_2,\mu,\eps}$ and that $\bQ^\star_{g_\varepsilon} \in \Pt_{b,be^\eps,\mu}$.

It follows that for $c_1 = \frac{1}{\gamma}$, $c_2 = \gamma$, and $\mu = P_0$, we have the equivalence $\Ptm = \ngb$. For this instantiation of $\Ptm$, the optimal sampler $\mathbf{Q}^\star_{c_1,c_2,\mu,\eps}$ for input $P$
is defined as a probability measure such that $Q \ll P_0$ and
\[
\frac{dQ}{dP_0}(x) = \textnormal{clip} \left( \frac{1}{r_P} \frac{dP}{dP_0}(x) ; b, b e^{\eps} \right), 
\]
where
\[
b = \frac{\gamma - \frac{1}{\gamma}}{(e^{\eps} - 1)(1 - \frac{1}{\gamma}) + \gamma - \frac{1}{\gamma}} = \frac{\gamma + 1}{e^\eps + \gamma} .
\]
This is exactly  \( \bQ^\star_\varepsilon \), the optimal sampler from Theorem~\ref{thm: local pure}. Thus, it remains to prove that $\bQ^\star_{g_\varepsilon} \in \Pt_{b,be^\eps,P_0}$ for $b =  \frac{\gamma + 1}{e^\eps + \gamma}$.

Combining Theorem~\ref{thm: local functional} and Corollary~\ref{cor: special case of f-LDP}, we obtain
\[
\bQ^\star_{g_\varepsilon}(P) \coloneqq \lambda^\star_{g_\varepsilon} P + (1 - \lambda^\star_{g_\varepsilon}) P_0, \quad 
\lambda^\star_{g_\varepsilon} = \frac{e^\varepsilon - 1}{(1 - \frac{1}{\gamma})e^\varepsilon + \gamma - 1} = \frac{\gamma(e^\varepsilon - 1)}{(\gamma - 1)(e^\eps + \gamma)}.
\]

Recall that
\[
\tilde{\mathcal{P}}_{b, be^\eps, P_0}
:= \left\{ Q \in \mathcal{P}(\mathcal{X}) : Q \ll P_0, \quad b \leq \frac{dQ}{dP_0} \leq be^\eps \quad P_0\text{-a.e.} \right\},
\]
and
\[
\ngb
:= \left\{ P \in \mathcal{P}(\mathcal{X}) : P \ll P_0, \quad \frac{1}{\gamma} \leq \frac{dP}{dP_0} \leq \gamma \quad P_0\text{-a.e.} \right\}.
\]

Therefore, if $P \in \ngb$, we have
\[
\frac{d\bQ^\star_{g_\varepsilon}}{dP_0} = \lambda^\star_{g_\varepsilon} \frac{dP}{dP_0} + (1 - \lambda^\star_{g_\varepsilon}).
\]
It follows that $\bQ^\star_{g_\varepsilon} \ll P_0$ and
\begin{align*}
    \frac{\lambda^\star_{g_\varepsilon}}{\gamma} + (1 - \lambda^\star_{g_\varepsilon})\leq \frac{d\bQ^\star_{g_\varepsilon}}{dP_0} \le \gamma \lambda^\star_{g_\varepsilon} + (1 - \lambda^\star_{g_\varepsilon}) \quad P_0\text{-a.e.}.
\end{align*}
Equivalently,
\begin{align*}
   \frac{e^\varepsilon - 1}{(\gamma - 1)(e^\eps + \gamma)} + \frac{\gamma^2 -e^\eps}{(\gamma - 1)(e^\eps + \gamma)}\leq \frac{d\bQ^\star_{g_\varepsilon}}{dP_0} \le \frac{\gamma^2(e^\varepsilon - 1)}{(\gamma - 1)(e^\eps + \gamma)} + \frac{\gamma^2 -e^\eps}{(\gamma - 1)(e^\eps + \gamma)} \quad P_0\text{-a.e.}.
\end{align*}
As a result,
\begin{align*}
    \frac{\gamma^2 -1}{(\gamma - 1)(e^\eps + \gamma)}\leq \frac{d\bQ^\star_{g_\varepsilon}}{dP_0} \le   \frac{(\gamma^2 -1)e^\eps}{(\gamma - 1)(e^\eps + \gamma)} \quad P_0\text{-a.e.}.
\end{align*}
Or equivalently,
\begin{align*}
    \frac{\gamma + 1}{e^\eps + \gamma}\leq \frac{d\bQ^\star_{g_\varepsilon}}{dP_0} \le   \frac{(\gamma + 1)e^\eps}{e^\eps + \gamma} \quad P_0\text{-a.e.}.
\end{align*}
This shows that $\bQ^\star_{g_\eps} \in \tilde{\mathcal{P}}_{b, be^\eps, P_0}$ and completes the proof.

\end{proof}


\section{Detailed experimental setup and additional experiments
}\label{sec: appendix experimental setup and edditional experiments}


\subsection{Mixture of Laplace distributions}\label{appendix: laplace mixture}



% \HG{$\frac{c_2 - c_1}{1 - c_1} \in \N$ does not naturally hold for mixture of Laplace distributions. I can make it hold by writing $\tilde{\mathcal{P}}_{\mathcal{L}} \subseteq \tilde{\mathcal{P}}_{0,\, \lceil e^{1/b}\rceil,\, h_{\mathcal{L}}}$ (see Example 3.1), but this would seem artificial. Anything you suggest I add in order to make the justification of $\frac{c_2 - c_1}{1 - c_1} \in \N$ assumption better?} \SA{I agree it may look artificial, so we can keep $c_1$ for Laplace in the above example. But did you example how did you address this issue for Fig 1?}\HG{For the experiments, I chose $c_1$ and $c_2$ in a way that $\frac{c_2 - c_1}{1 - c_1} \in \N$ by defining the local neighborhood using ceiling and floor functions.} \SA{Make sure that you mention it in the appendix.}\HG{Sure.}

\textbf{Proof of $\mathbf{\tilde{\mathcal{P}}_{\mathcal{L}} \subseteq \tilde{\mathcal{P}}_{e^{-1/b},\, e^{1/b},\, h_{\mathcal{L}}}}$ (Example \ref{example: Gaussian Laplace universe})}

Let the mixture of Laplace distributions with fixed scale parameter \( b \) be defined as
\[
\tilde{\mathcal{P}}_{\mathcal{L}} = \left\{ \sum_{i=1}^k \lambda_i \, \mathcal{L}(m_i, b) : k \in \mathbb{N},\, \lambda_i \geq 0,\, \sum_{i=1}^k \lambda_i = 1,\, \|m_i\|_1 \leq 1 \right\},
\]
where \( \mathcal{L}(m, b) \) denotes the \(n\)-dimensional Laplace distribution with mean \( m \in \mathbb{R}^n \) and scale parameter \( b > 0 \). We aim to show that
\[
\tilde{\mathcal{P}}_{\mathcal{L}} \subseteq \tilde{\mathcal{P}}_{e^{-1/b},\, e^{1/b},\, h_{\mathcal{L}}},
\]
where \( h_{\mathcal{L}} \) is the density of the zero-mean \(n\)-dimensional Laplace distribution with scale parameter \( b \).

The density of an \(n\)-dimensional Laplace distribution with location \( m \in \mathbb{R}^n \) and scale \( b \) is given by
\[
\mathcal{L}(x \mid m, b) = \frac{1}{(2b)^n} \exp\left( -\frac{\|x - m\|_1}{b} \right),
\quad x \in \mathbb{R}^n.
\]

Now, fix any \( m \in \mathbb{R}^n \) such that \( \|m\|_1 \leq 1 \). Then for any \( x \in \mathbb{R}^n \), we have the following bounds:
\begin{align*}
\frac{\|x\|_1 - \|m\|_1}{b}
&\leq \frac{\|x - m\|_1}{b}
\leq \frac{\|x\|_1 + \|m\|_1}{b} \\
\Rightarrow \quad
\frac{\|x\|_1 - 1}{b}
&\leq \frac{\|x - m\|_1}{b}
\leq \frac{\|x\|_1 + 1}{b}.
\end{align*}
Exponentiating both sides, we get:
\[
\exp\left( -\frac{\|x\|_1 + 1}{b} \right)
\leq \exp\left( -\frac{\|x - m\|_1}{b} \right)
\leq \exp\left( -\frac{\|x\|_1 - 1}{b} \right).
\]
Multiplying by the constant \( \frac{1}{(2b)^n} \), we obtain:
\[
\frac{e^{-1/b}}{(2b)^n} \exp\left( -\frac{\|x\|_1}{b} \right)
\leq \mathcal{L}(x \mid m, b)
\leq \frac{e^{1/b}}{(2b)^n} \exp\left( -\frac{\|x\|_1}{b} \right),
\]
which implies
\[
e^{-1/b} \, \mathcal{L}(x \mid \mathbf{0}, b) \leq \mathcal{L}(x \mid m, b) \leq e^{1/b} \, \mathcal{L}(x \mid \mathbf{0}, b).
\]
For any distribution \( P \in \tilde{\mathcal{P}}_{\mathcal{L}} \), we can write
\[
p(x) = \sum_{i=1}^k \lambda_i \, \mathcal{L}(x \mid m_i, b).
\]
Since each \( m_i \) satisfies \( \|m_i\|_1 \leq 1 \), the above inequality applies to each mixture component. Thus, for all \( x \in \mathbb{R}^n \),
\[
e^{-1/b} \, \mathcal{L}(x \mid \mathbf{0}, b)
\leq \sum_{i=1}^k \lambda_i \mathcal{L}(x \mid m_i, b)
\leq e^{1/b} \, \mathcal{L}(x \mid \mathbf{0}, b).
\]
Defining \( p_0(x) = \mathcal{L}(x \mid \mathbf{0}, b) \), we conclude that for every \( P \in \tilde{\mathcal{P}}_{\mathcal{L}} \),
\[
e^{-1/b} \leq \frac{p(x)}{p_0(x)} \leq e^{1/b}, \qquad \forall x \in \mathbb{R}^n,
\]
which confirms that \( \tilde{\mathcal{P}}_{\mathcal{L}} \subseteq \tilde{\mathcal{P}}_{e^{-1/b},\, e^{1/b},\, h_{\mathcal{L}}} \).


\textbf{Experimental details of Figure \ref{fig: Laplace pure and functional local vs global}}




The original distribution is a mixture of four two-dimensional Laplace distributions, each with scale parameter \( b = 2 \) and means at \( (1,0) \), \( (0,1) \), \( (-1,0) \), and \( (0,-1) \), respectively, with equal weights \( \frac{1}{4} \). That is, the input distribution is given by

\[
P_{\textsf{input}} = \sum_{i=1}^4 \frac{1}{4} \, \mathcal{L}(m_i, 2),
\]
where the \( m_i \) are defined as above.

From Example~\ref{example: Gaussian Laplace universe}, we know that \( P_{\textsf{input}} \in \tilde{\mathcal{P}} \), where
\[
\tilde{\mathcal{P}} = \left\{ \sum_{i=1}^k \lambda_i \, \mathcal{L}(m_i, b) : k \in \mathbb{N},\, \lambda_i \geq 0,\, \sum_{i=1}^k \lambda_i = 1,\, \|m_i\|_1 \leq 1 \right\},
\]
and furthermore, \( \tilde{\mathcal{P}} \subseteq \tilde{\mathcal{P}}_{e^{-1/b},\, e^{1/b},\, h_{\mathcal{L}}} \), where \( h_{\mathcal{L}} \) denotes the density of a two-dimensional Laplace distribution with mean zero and scale parameter \( b = 2 \).

For \( b = 2 \), we define the local and global minimax universes as
\[
\tilde{\mathcal{P}}_{\textsf{local}} = \tilde{\mathcal{P}}_{\frac{1}{2},\, 2,\, h_{\mathcal{L}}}
\quad \text{and} \quad
\tilde{\mathcal{P}}_{\textsf{global}} = \tilde{\mathcal{P}}_{\frac{1}{6},\, 6,\, h_{\mathcal{L}}}.
\]

These universes are chosen such that \( \tilde{\mathcal{P}}_{\textsf{local}} \subseteq \tilde{\mathcal{P}}_{\textsf{global}} \), \( P_{\textsf{input}} \in \tilde{\mathcal{P}}_{\textsf{local}} \cap \tilde{\mathcal{P}}_{\textsf{global}} \),  and the ratio \( \frac{c_2 - c_1}{1 - c_1} \in \mathbb{N} \), as required by Assumption~\ref{assumption: norm}.
Moreover, we note that \( \tilde{\mathcal{P}}_{e^{-1/b},\, e^{1/b},\, h_{\mathcal{L}}} \subseteq \tilde{\mathcal{P}}_{\frac{1}{2},\, 2,\, h_{\mathcal{L}}} \), i.e., the condition in Assumption \ref{assumption: norm} is achieved by slightly adjusting the bounds \( c_1 \) and \( c_2 \) using floor and ceiling functions.

We evaluate the performance of minimax-optimal samplers on the input distribution under two LDP settings: comparing local minimax-optimal samplers with global minimax-optimal samplers under both \( \nu \)-GLDP and \( \varepsilon \)-LDP. In the pure-LDP setting, we compare the global minimax-optimal sampler of \citet{park2024exactly} with the local minimax-optimal sampler from Theorem~\ref{thm: local pure} for \( \varepsilon = 1 \). In the \( \nu \)-GLDP setting, we compare our global minimax-optimal sampler from Corollary~\ref{cor: special case GDP} with the local minimax-optimal sampler from Theorem~\ref{thm: local functional}, using the special case \( g = G_\nu \) with \( \nu = 1.5 \).

Under both settings, Figure~\ref{fig: Laplace pure and functional local vs global} demonstrates that the local minimax-optimal sampler better preserves the input distribution compared to the global minimax-optimal sampler, given the same level of privacy.



% Let $\mathcal{L}(x \mid m, b)$ denote the probability density function of a two-dimensional Laplace distribution with location parameter $m = (\mu_1, \mu_2)$ and scale parameter $b > 0$, defined for $x = (x_1, x_2) \in \mathbb{R}^2$ as:
% \[
% \mathcal{L}(x \mid m, b) = \frac{1}{4b^2} \exp\left( -\frac{|x_1 - \mu_1| + |x_2 - \mu_2|}{b} \right).
% \]

% Define the set $\tilde{\mathcal{P}}$ of probability distributions as:
% \[
% \tilde{\mathcal{P}} = \left\{ \sum_{i=1}^k \lambda_i \mathcal{L}( m_i, b) : k \in \mathbb{N}, \, \lambda_i \geq 0, \, \sum_{i=1}^k \lambda_i = 1, \, \|m_i\|_1 \leq 1 \right\}.
% \]

% For any $m \in \mathbb{R}^2$ such that $\|m\|_1 \leq 1$, we have the following bounds:
% \begin{align*}
%   \frac{|x_1| + |x_2| - 1}{b} &\leq \frac{|x_1| - |\mu_1| + |x_2| - |\mu_2|}{b} \\
%   &\leq \frac{|x_1 - \mu_1| + |x_2 - \mu_2|}{b} \\
%   &\leq \frac{|x_1| + |\mu_1| + |x_2| + |\mu_2|}{b} \\
%   &\leq \frac{|x_1| + |x_2| + 1}{b}.
% \end{align*}

% Therefore, for all $x \in \mathbb{R}^2$, we have:
% \begin{align*}
%     \frac{e^{-1/b}}{4b^2} \exp\left( -\frac{|x_1| + |x_2|}{b} \right)
%     &= \frac{1}{4b^2} \exp\left( -\frac{|x_1| + |x_2| + 1}{b} \right) \\
%     &\leq \frac{1}{4b^2} \exp\left( -\frac{|x_1 - \mu_1| + |x_2 - \mu_2|}{b} \right) \\
%     &\leq \frac{1}{4b^2} \exp\left( -\frac{|x_1| + |x_2| - 1}{b} \right) \\
%     &= \frac{e^{1/b}}{4b^2} \exp\left( -\frac{|x_1| + |x_2|}{b} \right).
% \end{align*}

% Consequently, for all $x \in \mathbb{R}^2$, we obtain:
% \[
% e^{-1/b} \, \mathcal{L}(x \mid \mathbf{0}, b) \leq \mathcal{L}(x \mid m, b) \leq e^{1/b} \, \mathcal{L}(x \mid \mathbf{0}, b).
% \]

% Therefore, for any distribution in $\tilde{\mathcal{P}}$,
% \[
% e^{-1/b} \, \mathcal{L}(x \mid \mathbf{0}, b)
% \leq \sum_{i=1}^k \lambda_i \mathcal{L}(x \mid m_i, b)
% \leq e^{1/b} \, \mathcal{L}(x \mid \mathbf{0}, b).
% \]

% Choosing $P_0 = \mathcal{L}(x \mid \mathbf{0}, b)$, it follows that for every $P \in \tilde{\mathcal{P}}$,
% \[
% e^{-1/b} \leq \frac{P(x)}{P_0(x)} \leq e^{1/b}, \qquad \forall x \in \mathbb{R}^2.
% \]
% The same argument can be generalized to any $ n$-dimensional Laplace mixtures for $n \geq 2$.

\subsection{Gaussian mixtures}\label{appendix: gaussian ring}

\textbf{Proof of $\tilde{\mathcal{P}}_{\mathcal{N}} \subseteq \tilde{\mathcal{P}}_{0,\, 1,\, h_{\mathcal{N}}}$ (Example \ref{example: Gaussian Laplace universe})} 
% \HG{Double check Example 3.1!}

let $\tilde{\mathcal{P}}_{\mathcal{N}}  = \Big\{ \sum_{i=1}^k \lambda_i \mathcal{N}(m_i, \sigma^2 I_n)  :  \lambda_i \geq 0, \, \sum_{i=1}^k \lambda_i = 1, \, \|m_i\|_2 \leq 1 \Big\}$. We want to show that  $\tilde{\mathcal{P}}_{\mathcal{N}} \subseteq \tilde{\mathcal{P}}_{0,\, 1,\, h_{\mathcal{N}}}$ , where $h_{\mathcal{N}}(x) = \frac{1}{(2\pi \sigma^2)^\frac{n}{2}} \exp\big(-\frac{[\max(0, \|x\|_2 - 1)]^2}{2\sigma^2} \big)$.

Fix a component centre \(m\in\mathbb{R}^{n}\) with \(\|m\|_{2}\le1\).
Its Gaussian density at \(x\) is
\[
\mathcal{N}_{m,\sigma^{2}I_{n}}[x]
=\frac{1}{(2\pi\sigma^{2})^{n/2}}
  \exp\Bigl(
     -\tfrac{\|x-m\|_{2}^{2}}{2\sigma^{2}}
  \Bigr).
\]
By the triangle inequality,
\[
\|x-m\|_{2}
\;\ge\;
\bigl|\|x\|_{2}-\|m\|_{2}\bigr|
\;\ge\;
\max\bigl(0,\|x\|_{2}-1\bigr),
\]
because \(\|m\|_{2}\le1\).
Squaring and dividing by \(2\sigma^{2}\) yields
\[
\exp\Bigl(
  -\tfrac{\|x-m\|_{2}^{2}}{2\sigma^{2}}
\Bigr)
\;\le\;
\exp\Bigl(
  -\tfrac{[\max(0,\|x\|_{2}-1)]^{2}}{2\sigma^{2}}
\Bigr).
\]
Multiplying by the common normalizing constant
\((2\pi\sigma^{2})^{-n/2}\) gives
\[
\mathcal{N}_{m,\sigma^{2}I_{n}}[x]
\;\le\;
h_{\mathcal{N}}(x),\qquad\forall x\in\mathbb{R}^{n}.
\]

Now take an arbitrary mixture
\(P=\sum_{i=1}^{k}\lambda_{i}\,\mathcal{N}(m_{i},\sigma^{2}I_{n})
\in\tilde{\mathcal{P}}_{\mathcal{N}}\).
Its density is
\[
p(x)=\sum_{i=1}^{k}\lambda_{i}\,
      \mathcal{N}_{m_{i},\sigma^{2}I_{n}}[x].
\]
Because each component satisfies
\(\mathcal{N}_{m_{i},\sigma^{2}I_{n}}[x]\le h_{\mathcal{N}}(x)\) and the
weights obey \(\lambda_{i}\ge0,\ \sum_{i}\lambda_{i}=1\), we have
\[
0\le p(x)\le h_{\mathcal{N}}(x)\qquad\forall x\in\mathbb{R}^{n}.
\]
Therefore \(P\in\tilde{\mathcal{P}}_{0,\,1,\,h_{\mathcal{N}}}\), as
claimed.


 
% Here are all the figures we have:
% \begin{figure}[ht]
%   \centering
%   \includegraphics[width=0.8\linewidth]{figures/GDP_LocalvsGlobal_mu0.5_h0.png}
%   \caption{Comparison of global and local minimax-optimal samplers for functional LDP in the special case of \(\nu\)-GLDP with \(\nu = 0.5\)..}
%   \label{fig: GDP local vs global nu = 0.5}
% \end{figure}

% \begin{figure}[ht]
%   \centering
%   \includegraphics[width=0.8\linewidth]{figures/feps_LocalvsGlobal_eps0.1_h0.png}
%   \caption{Comparison of global and local minimax-optimal samplers for functional LDP in the special case of $g_\eps$ with \(\eps = 0.1\).}
%   \label{fig: f_eps local vs global eps = 0.1}
% \end{figure}

% \begin{figure}[ht]
%   \centering
%   \includegraphics[width=0.8\linewidth]{figures/feps_LocalvsGlobal_eps0.5_h0.png}
%   \caption{Comparison of global and local minimax-optimal samplers for functional LDP in the special case of $g_\eps$ with \(\eps = 0.5\).}
%   \label{fig: f_eps local vs global eps = 0.5}
% \end{figure}

% % \begin{figure}[ht]
% %   \centering
% %   \includegraphics[width=0.8\linewidth]{figures/pure_GlobalvsLocal_eps0.5_var0.5.png}
% %   \caption{Comparison of global and local minimax-optimal samplers for pure LDP for $\eps = 0.5$ and $\sigma^2 = 0.5$.}
% %   \label{fig: pure local vs global eps = 0.5 var = 0.5}
% % \end{figure}


% \begin{figure}[ht]
%   \centering
%   \includegraphics[width=0.8\linewidth]{figures/pure_GlobalvsLocal_eps0.5_var0.1.png}
%   \caption{Comparison of global and local minimax-optimal samplers for pure LDP for $\eps = 0.5$ and $\sigma^2 = 0.1$.}
%   \label{fig: pure local vs global eps = 0.5 var = 0.1}
% \end{figure}

% \begin{figure}[ht]
%   \centering
%   \includegraphics[width=0.8\linewidth]{figures/pure_GlobalvsLocal_eps0.5_var0.25.png}
%   \caption{Comparison of global and local minimax-optimal samplers for pure LDP for $\eps = 0.5$ and $\sigma^2 = 0.25$.}
%   \label{fig: pure local vs global eps = 0.5 var = 0.25}
% \end{figure}

% \begin{figure}[ht]
%   \centering
%   \includegraphics[width=0.8\linewidth]{figures/pure_GlobalvsLocal_eps0.5_var0.5.png}
%   \caption{Comparison of global and local minimax-optimal samplers for pure LDP for $\eps = 0.5$ and $\sigma^2 = 0.5$.}
%   \label{fig: pure local vs global eps = 0.5 var = 0.5}
% \end{figure}



% % \begin{figure}[ht]
% %   \centering
% %   \includegraphics[width=0.8\linewidth]{figures/pure_GlobalvsLocal_eps0.5_var0.1.png}
% %   \caption{Comparison of global and local minimax-optimal samplers for pure LDP for $\eps = 0.5$ and $\sigma^2 = 0.1$.}
% %   \label{fig: pure local vs global eps = 0.5 var = 0.1}
% % \end{figure}

% \begin{figure}[ht]
%   \centering
%   \includegraphics[width=0.8\linewidth]{figures/pure_LocalvsGlobal_eps0.1_var0.5.png}
%   \caption{Comparison of global and local minimax-optimal samplers for pure LDP for $\eps = 0.1$ and $\sigma^2 = 0.5$.}
%   \label{fig: pure local vs global eps = 0.1 var = 0.5}
% \end{figure}

% The \emph{Gaussian ring distribution} with \( k \in \mathbb{N} \) modes and component-wise variance \( \sigma^2 > 0 \) is defined as the uniform mixture of \( k \) Gaussian distributions in \( \mathbb{R}^2 \), where each component has equal weight, identical covariance matrix \( \sigma^2 I_2 \), and means that are equally spaced on the unit circle. Specifically, the mean of the \( i \)-th Gaussian component is given by
% \[
% m_i = \left( \cos \left( \frac{2\pi i}{k} \right), \, \sin \left( \frac{2\pi i}{k} \right) \right), \quad \text{for } i = 1, \dots, k.
% \]
% One of the Gaussian components is centered at \( (1, 0) \), corresponding to \( i = 0 \). The resulting density of the Gaussian ring distribution is:
% \[
% p(x) = \frac{1}{k} \sum_{i=1}^k \mathcal{N}_{m_i, \sigma^2 I_2}[x],
% \]
% where \( \mathcal{N}_{m, \Sigma}[x] \) denotes the density of a multivariate Gaussian distribution with mean \( m \) and covariance \( \Sigma \), evaluated at point \( x \). Explicitly, for any \( m \in \mathbb{R}^n \), we have:
% \[
% \mathcal{N}_{m, \Sigma}[x] = \frac{1}{(2\pi)^{n/2} \sqrt{\det(\Sigma)}} \exp \left( -\frac{1}{2} (x - m)^T \Sigma^{-1} (x - m) \right).
% \]
% We also use the notation \( \mathcal{N}(m, \Sigma) \) to refer to the corresponding Gaussian distribution itself.

% \paragraph{Specific configuration:} In our experiments (see Figures \HG{$\dots$}), we consider the Gaussian ring distribution with \( k = 3 \) modes and different component-wise variances \( \sigma^2 = 0.1 \), \( \sigma^2 = 0.25 \), and \( \sigma^2 = 0.5 \). This distribution is used as the original (pre-perturbation) data distribution for a client in the context of both our functional and pure LDP sampler evaluations. Our experiments compare global and local minimax-optimal samplers in both functional and pure LDP scenarios. For the global case, we follow the experimental setup of \citet{park2024exactly} to choose the reference distribution and the universe of probability distributions as follows:
% The universe \( \tilde{\mathcal{P}} \) comprises Gaussian mixture distributions, where each Gaussian component has a mean located within the unit ball centered at the origin and a shared covariance of \( 0.5 I_2 \). Formally, we define:
% \[
% \tilde{\mathcal{P}} = \left\{ \sum_{i=1}^k \lambda_i \mathcal{N}(m_i, \sigma^2 I_2) : k \in \mathbb{N}, \lambda_i \geq 0, \sum_{i=1}^k \lambda_i = 1, \|m_i\|_2 \leq 1 \right\},
% \]
% with \( \sigma^2 = 0.5 \). Additionally, it is shown in \citet{park2024exactly} that this class is a subset of the larger class \( \tilde{\mathcal{P}}_{0,1,\tilde{h}} \), where \( \tilde{h}(x) \) is defined by:
% \[
% \tilde{h}(x) = \frac{1}{2\pi\sigma^2} \exp \left( - \frac{[\max(0, \|x\| - 1)]^2}{2\sigma^2} \right).
% \]
It is worth noting that $h_\mathcal{N}$ is not a probability density. In order to make $h_{\mathcal{N}}$ a valid probability distribution, we normalize it by its integral and define
\[
c_2 = \int h_{\mathcal{N}}(x) \, dx, \quad \text{and} \quad h(x) = \frac{h_{\mathcal{N}}(x)}{c_2}.
\]
Accordingly, we have \( \tilde{\mathcal{P}}_{0,1,h_{\mathcal{N}}} = \tilde{\mathcal{P}}_{0,c_2,h} \). 

% For the local minimax case, we need to define a reference distribution in neighborhood of which we are definition the local minimax.  We define the reference measure $\tilde{P}_0$ as follows:
% \[\tilde{p}_0(x) = \frac{1}{2\pi\sigma^2} \exp \left( - \frac{\min\limits_{1\le i \le 3} \|x-m_i\|^2}{2\sigma^2} \right).\]
% To make $\tilde{P}_0$ a valid probability distribution, we get $P_0$ with the PDF:
% \[p_0(x) = \dfrac{\frac{1}{2\pi\sigma^2} \exp \left( - \frac{\min\limits_{1\le i \le 3} \|x-m_i\|^2}{2\sigma^2} \right)}{Z} \qquad Z = \int \tilde{p}_0(x) dx.\]
%  We can define the neighborhood $N_{\gamma}(P_0)$ for $\gamma = 3$ as:
% \[N_{3}(P_0) = \left\{ P \in \mathcal{C}(\mathbb{R}^2 ) :\quad \frac{1}{3} \leq \frac{p(x)}{p_0(x)} \leq 3 \quad \forall x \in \mathbb{R}^2 \right\}.\]
% For input distribution $p(x) = \frac{1}{3} \sum_{i=1}^3 \mathcal{N}_{m_i, \sigma^2 I_2}[x]$ we have:
% \[\frac{1}{3}\tilde{p}_0(x) \leq p(x) \leq \tilde{p}_0(x) \qquad \forall x \in \mathbb{R}^2.\]
% This, in turn, gives us the following inequality:
% \[\frac{1}{3}Zp_0(x) \leq p(x) \leq Z p_0(x) \qquad \forall x \in \mathbb{R}^2.\]
% Since $1 \leq Z \leq 3$, it follows that:
% \[\frac{1}{3}p_0(x) \leq p(x) \leq 3 p_0(x) \qquad \forall x \in \mathbb{R}^2.\]
% Therefore, $\sum_{i=1}^3 \mathcal{N}_{m_i, \sigma^2 I_2}[x] \in N_3(P_0)$. For the local minimax experiments, we consider the neighborhood \( N_3(P_0) \) and apply the samplers introduced in Theorems~\ref{thm: local functional} and~\ref{thm: local pure}, corresponding to the functional and pure LDP settings, respectively. \HG{We do not necessarily have $N_3(P_0) \subseteq \tilde{\mathcal{P}}_{0,1,\tilde{h}}$. We can generalize the universe to $\tilde{P}_{0,3,\tilde{h}}$.} For the functional LDP experiments, we consider two choices for the function \( g \): the GLDP function $G_{\nu}$ for $\nu = 0.5$ from  Corollary~\ref{cor: special case GDP}, and the function \( g_\varepsilon \) defined in Corollary~\ref{cor: special case of f-LDP} for $\eps = 0.1$ and $\eps = 0.5$. We compare local minimax sampler applied to Gaussian ring distribution with global minimax-optimal sampler in Theorem \ref{thm: global functional}.

\subsection{Finite sample space numerical results under pure LDP}\label{appendix: finite space}

\textbf{Experimental details of Figure \ref{fig: finite worst-case k = 20}}

We compare local and global minimax-optimal samplers in the finite setting \( \mathcal{X} = [k] \), where the global universe is defined as \( \tilde{\mathcal{P}}_{\textsf{global}} = \mathcal{P}([k]) = \tilde{\mathcal{P}}_{0, k, \mu_k} \), with \( \mu_k \) denoting the uniform distribution over \([k]\). The local neighborhood is specified as \( \tilde{\mathcal{P}}_{\textsf{local}} = \mathcal{N}_\gamma(\mu_k) = \tilde{\mathcal{P}}_{\frac{1}{\gamma},\, \gamma,\, \mu_k} \), where \( \gamma = \frac{k}{2} - 1 \), consistent with the finite-space version of Theorem~\ref{thm: local pure}.

For the selected values \( k \in \{10, 20, 100\} \), it is easy to verify that \( \gamma \in \mathbb{N} \), and the uniform distribution \( \mu_k \) is \( (\alpha, \frac{1}{\alpha}, 1) \)-decomposable, where
\[
\alpha = \frac{1 - c_1}{c_2 - c_1}
= \frac{1 - \frac{1}{\gamma}}{\gamma - \frac{1}{\gamma}}
= \frac{1}{\gamma + 1}
= \frac{2}{k}.
\]
Therefore, the decomposability condition in Theorem~\ref{thm: local pure} is satisfied, and we can apply its finite-space version by setting \( \mathcal{X} = [k] \).

The local minimax risk is then given by
\[
\mathcal{R}\big(\mathcal{Q}_{[k], \mathcal{P}([k]), \varepsilon}, N_\gamma(\mu_k), f\big) 
= \frac{1 - r_1}{r_2 - r_1} \, f(r_2) + \frac{r_2 - 1}{r_2 - r_1} \, f(r_1),
\]
where
\[
b \coloneqq \frac{\gamma - \tfrac{1}{\gamma}}{(e^\varepsilon - 1)(1 - \tfrac{1}{\gamma}) + \gamma - \tfrac{1}{\gamma}} 
= \frac{\gamma + 1}{\gamma + e^\varepsilon}, \quad
r_1 \coloneqq \frac{1}{\gamma b} = \frac{e^\varepsilon + \gamma}{\gamma(\gamma + 1)}, \quad
r_2 \coloneqq \frac{\gamma}{b e^\varepsilon} = \frac{\gamma(e^\varepsilon + \gamma)}{e^\varepsilon(\gamma + 1)}.
\]

On the other hand, the global minimax risk is given by~\citep[Theorem 3.1]{park2024exactly}:
\[
\mathcal{R}\big(\mathcal{Q}_{[k], \mathcal{P}([k]), \varepsilon}, \mathcal{P}([k]), f\big)
= \frac{e^{\varepsilon}}{e^{\varepsilon} + k - 1} 
  f\left( \frac{e^{\varepsilon} + k - 1}{e^{\varepsilon}} \right)
+ \frac{k - 1}{e^{\varepsilon} + k - 1} f(0).
\]

Therefore, in Figure~\ref{fig: finite worst-case k = 20}, we compare the theoretical worst-case \( f \)-divergence losses achieved by the local minimax-optimal sampler and the global minimax-optimal sampler:
\[
\mathcal{R}\big(\mathcal{Q}_{[k], \mathcal{P}([k]), \varepsilon}, N_\gamma(\mu_k), f\big)
\quad \text{and} \quad
\mathcal{R}\big(\mathcal{Q}_{[k], \mathcal{P}([k]), \varepsilon}, \mathcal{P}([k]), f\big).
\]



\textbf{Additional numerical results}

We replicated the procedure used in Section~\ref{experiment: finite} to generate Figure~\ref{fig: finite worst-case k = 20} , but now with sample sizes \(k = 10\) and \(k = 100\). As shown in Figures~\ref{fig: finite worst-case k = 10} and~\ref{fig: finite worst-case k = 100} and comparing with Figure \ref{fig: finite worst-case k = 20}, the local minimax sampler consistently attains a smaller worst-case loss than the global minimax sampler across different $\eps$ values. Also, the performance gap decreases for larger values of $k$.

\begin{figure}[ht]
  \centering
  \includegraphics[width=0.7\linewidth]{figures/result_finite_k10.pdf}
  \caption{Theoretical worst-case 
$f$-divergences of global and local minimax samplers under the pure LDP setting with uniform reference distribution $\mu_k$ over finite space ($k = 10$).
% (Left: KL divergence; Center: Total variation distance; Right: Squared Hellinger distance.)
}
  \label{fig: finite worst-case k = 10}
\end{figure}

\begin{figure}[ht]
  \centering
  \includegraphics[width=0.7\linewidth]{figures/result_finite_k100.pdf}
  \caption{Theoretical worst-case 
$f$-divergences of global and local minimax samplers under the pure LDP setting with uniform reference distribution $\mu_k$ over finite space ($k = 100$).
% (Left: KL divergence; Center: Total variation distance; Right: Squared Hellinger distance.)
}
  \label{fig: finite worst-case k = 100}
\end{figure}

\subsection{Experimental details of continuous sample space}\label{appendix: details of continuous}



In the continuous setting with \( \mathcal{X} = \mathbb{R} \), we fix the universe \( \tilde{\mathcal{P}}_{\textsf{local}} \) and evaluate the empirical worst-case \( f \)-divergence over 100 randomly generated client distributions \( P_1, \ldots, P_{100} \in \tilde{\mathcal{P}}_{\textsf{local}} \). Each \( P_j \) represents a client and is constructed as a mixture of a random number of one-dimensional Laplace components with scale parameter \( b = 1 \). The full generation procedure is described below.

We define the mixture class of Laplace distributions with fixed scale parameter \( b \) as
\[
\tilde{\mathcal{P}}_{\mathcal{L}} = \left\{ \sum_{i=1}^k \lambda_i \, \mathcal{L}(m_i, b) : k \in [K],\, \lambda_i \geq 0,\, \sum_{i=1}^k \lambda_i = 1,\, |m_i| \leq 1 \right\},
\]
where \( \mathcal{L}(m, b) \) denotes the one-dimensional Laplace distribution with mean \( m \in \mathbb{R} \) and scale parameter \( b > 0 \). To prevent an unbounded number of components in each mixture, we impose an upper bound \( K \) on the number of Laplace components per client.

Each \( P_j \in \tilde{\mathcal{P}}_{\mathcal{L}} \) is generated by randomly selecting \( k \), \( \lambda_i \), and \( m_i \) as follows:  
First, sample \( \tilde{k} \) from a Poisson distribution with mean \( k_0 \), and set \( k = \min(\tilde{k} + 1, K) \). Then, sample each \( m_1, \ldots, m_k \) independently from the uniform distribution on \( [-1, 1] \), and sample weights \( (\lambda_1, \ldots, \lambda_k) \) from the uniform distribution on \( \mathcal{P}([k]) \). In this experiment, to maintain consistency with~\citet{park2024exactly}, we use \( K = 10 \) and \( k_0 = 2 \).

We define the local and global universes as
\[
\tilde{\mathcal{P}}_{\textsf{local}} = \tilde{\mathcal{P}}_{1/3,\, 3,\, h_{\mathcal{L}}}
\quad \text{and} \quad
\tilde{\mathcal{P}}_{\textsf{global}} = \tilde{\mathcal{P}}_{1/9,\, 9,\, h_{\mathcal{L}}},
\]
where \( h_{\mathcal{L}} \) denotes the density of the Laplace distribution with mean zero and scale \( b = 1 \).

We evaluate the empirical worst-case divergence of each sampler using the maximum
\[
\max_{j \in [100]} D_f\bigl(P_j \,\|\, \mathbf{Q}(P_j)\bigr).
\]
The local minimax sampler is instantiated from Theorem~\ref{thm: local pure}, while the global minimax sampler corresponds to the optimal sampler from~\citep[Theorem 3.3]{park2024exactly}.



\subsection{Finite sample space numerical results under  GLDP}  \label{appendix: experiment GDP local global discrete}

In this section, we adopt the same experimental setup as in Section~\ref{experiment: finite} to evaluate the worst-case $f$-divergence of the local and global minimax samplers under $\nu$-GLDP constraints. To this end, we follow the same procedure using the same global and local universes. However, the computation of minimax risk differs: for the global minimax risk, we use the optimal value corresponding to the optimal sampler characterized in Corollary~\ref{cor: special case GDP}, while for the local minimax risk, we instantiate  Theorem~\ref{thm: local functional} with $g = G_\nu$ (both the finite sample space version of the results).


\begin{figure}[ht]
  \centering
  \includegraphics[width=0.7\linewidth]{figures/result_finite_GDP_k10.pdf}
  \caption{Theoretical worst-case 
$f$-divergences of global and local minimax samplers under the $\nu$-GLDP setting with uniform reference distribution $\mu_k$ over finite space ($k = 10$).
% (Left: KL divergence; Center: Total variation distance; Right: Squared Hellinger distance.)
}
  \label{fig: finite worst-case GDP k = 10}
\end{figure}


\begin{figure}[ht]
  \centering
  \includegraphics[width=0.7\linewidth]{figures/result_finite_GDP_k20.pdf}
  \caption{Theoretical worst-case 
$f$-divergences of global and local minimax samplers under the $\nu$-GLDP setting with uniform reference distribution $\mu_k$ over finite space ($k = 20$).
% (Left: KL divergence; Center: Total variation distance; Right: Squared Hellinger distance.)
}
  \label{fig: finite worst-case GDP k = 20}
\end{figure}

\begin{figure}[ht]
  \centering
  \includegraphics[width=0.7\linewidth]{figures/result_finite_GDP_k100.pdf}
  \caption{Theoretical worst-case 
$f$-divergences of global and local minimax samplers under the $\nu$-GLDP setting with uniform reference distribution $\mu_k$ over finite space ($k = 100$).
% (Left: KL divergence; Center: Total variation distance; Right: Squared Hellinger distance.)
}
  \label{fig: finite worst-case GDP k = 100}
\end{figure}

We compare the worst-case $f$-divergence under $\nu$-GDP for local and global minimax-optimal samplers across different values of $\nu \in \{0.1, 0.5, 1, 2\}$. The numerical results in Figures~\ref{fig: finite worst-case GDP k = 10}, \ref{fig: finite worst-case GDP k = 20}, and \ref{fig: finite worst-case GDP k = 100} demonstrate that for various sample sizes $k \in \{10, 20, 100\}$, the local minimax-optimal sampler consistently achieves lower minimax risk than the global sampler across all privacy parameter values.



\subsection{Continuous sample space numerical results under  GLDP} \label{appendix: experiment GDP local global continuous}

We follow the same experimental procedure described in Appendix~\ref{appendix: details of continuous}, with the only difference being that we now compare the local and global minimax-optimal samplers under $\nu$-GDP instead of $\varepsilon$-LDP. The local and global universes, as well as the distribution generation process, remain unchanged. For the global minimax sampler, we use the construction provided in Corollary~\ref{cor: special case GDP}, while the local minimax sampler is instantiated from Theorem~\ref{thm: local functional} with \( g = G_\nu \). We conduct the comparison across various values of \( \nu \in \{0.1,0.5,1,2\} \), evaluating the resulting samplers using three $f$-divergences: KL divergence, total variation distance, and squared Hellinger distance.


\begin{figure}[ht]
  \centering
  \includegraphics[width=0.7\linewidth]{figures/result_1DLaplaceMix_GLDP.pdf}
  \caption{Empirical worst-case 
$f$-divergences of global and local minimax samplers under $\nu$-GLDP setting, over 100 experiments on a 1-D Laplace mixture.}
  \label{fig: continuous worst-case GDP}
\end{figure}

As illustrated in Figure~\ref{fig: continuous worst-case GDP}, in the $\nu$-GDP setting—similar to the $\varepsilon$-LDP case—the local minimax sampler consistently achieves a smaller worst-case $f$-divergence than the global minimax sampler across all three divergences and privacy parameter values.


\subsection{Additional numerical results on continuous sample space }\label{appendix: high dimension laplace}
We extend the experiment from Section~\ref{subsec: experiment 1d continuous} to the two-dimensional setting \( \mathcal{X} = \mathbb{R}^2 \). Figure~\ref{fig: continuous worst-case 2-D} compares the worst-case empirical \( f \)-divergence between the global sampler of~\citet{park2024exactly} and the local sampler described in Theorem~\ref{thm: local functional} for $g = g_\eps$.

We follow the same experimental procedure as in Section~\ref{subsec: experiment 1d continuous}, detailed in Appendix~\ref{appendix: details of continuous}, to generate 100 client distributions. Each distribution is constructed as a mixture of 2-D Laplace components with fixed scale parameter \( b = 1 \). The only difference from the one-dimensional case is the extension to 2-D Laplace mixtures.

The mixture class of 2-D Laplace distributions with scale parameter \( b \) is defined as
\[
\tilde{\mathcal{P}}_{\mathcal{L}} = \left\{ \sum_{i=1}^k \lambda_i \, \mathcal{L}(m_i, b) : k \in [K],\, \lambda_i \geq 0,\, \sum_{i=1}^k \lambda_i = 1,\, \|m_i\|_1 \leq 1 \right\},
\]
where \( \mathcal{L}(m, b) \) denotes a two-dimensional Laplace distribution with mean \( m \in \mathbb{R}^2 \) and scale \( b > 0 \). To control the complexity of each mixture, we impose an upper bound \( K \) on the number of components per distribution.

Each distribution \( P_j \in \tilde{\mathcal{P}}_{\mathcal{L}} \) is generated by randomly sampling \( k \), \( \lambda_i \), and \( m_i \) as follows:  
First, sample \( \tilde{k} \sim \text{Poisson}(k_0) \), and set \( k = \min(\tilde{k} + 1, K) \). Then, sample each \( m_i \in \mathbb{R}^2 \) independently from the uniform distribution on the \(\ell_1\) ball \( \{ x \in \mathbb{R}^2 : \|x\|_1 \leq 1 \} \), and draw the weights \( (\lambda_1, \ldots, \lambda_k) \) uniformly from the probability simplex \( \mathcal{P}([k]) \). In this experiment, following the setup in~\citet{park2024exactly}, we use \( K = 10 \) and \( k_0 = 2 \).

The local and global universes are defined as
$
\tilde{\mathcal{P}}_{\textsf{local}} = \tilde{\mathcal{P}}_{1/3,\, 3,\, h_{\mathcal{L}}}
\quad $ and $
\tilde{\mathcal{P}}_{\textsf{global}} = \tilde{\mathcal{P}}_{1/9,\, 9,\, h_{\mathcal{L}}}$,
where \( h_{\mathcal{L}} \) denotes the density of the two-dimensional Laplace distribution with mean zero and scale parameter \( b = 1 \).










% \begin{figure}[ht]
%   \centering
%   \includegraphics[width=0.8\linewidth]{figures/result_1DLaplaceMix_b=1_c=3.png}
%   \caption{Empirical worst-case 
% $f$-divergences of global and local minimax samplers under the pure LDP setting, averaged over 100 experiments on a 1-D Laplace mixture.
% (Left: KL divergence; Center: Total variation distance; Right: Squared Hellinger distance.)}
%   \label{fig: continuous worst-case}
% \end{figure}



\begin{figure}[ht]
  \centering
  \includegraphics[width=0.7\linewidth]{figures/result_2DLaplaceMix.pdf}
  \caption{Empirical worst-case 
$f$-divergences of global and local minimax samplers under the pure LDP setting, over 100 experiments on a 2-D Laplace mixture.
% (Left: KL divergence; Center: Total variation distance; Right: Squared Hellinger distance.)
}
  \label{fig: continuous worst-case 2-D}
\end{figure}

In the two-dimensional case, similar to the one-dimensional setting, the local sampler outperforms the global sampler for nearly all fixed values of the privacy parameter.


% \textcolor{blue}{[I guess the problem that the experiment takes too long is the integration part. Also, currently integration for higher dimension is inaccurate, which seems to be the reason error increasing even when privacy budget increases. Maybe we could replace the integration to Monte-Carlo method. $\widehat{D_f(P||Q)}:=\frac{1}{N}\sum_i f(p(X_i)/q(X_i))$ where $X_i\sim Q$. For instance, if we set $N=1000$, I guess we could get a fairly smooth plot for $10\sim 100$-dimensional distribution when $\lambda_\eps^\star$ is not too close to 0. But changing the way of evaluation could incur too many subsequent works..]}


% We can show that $N_3(P_0) \subseteq \tilde{\mathcal{P}}_{0,1,\tilde{h}}$, which confirms that the local neighborhood is a subset of the global universe.
% \( \tilde{\mathcal{P}}_{0,1,\tilde{h}} \) is defined as \( \tilde{\mathcal{P}}_{0,c_2,h} \) where:
%     \[
%     c_2 = \int \tilde{h}(x) dx, \quad h(x) = \frac{\tilde{h}(x)}{c_2}, \quad \tilde{\mathcal{P}}_{0,c_2,h} = \left\{ P \in \mathcal{C}(\mathbb{R}^2 ) : \quad 0 \le \frac{p(x)}{h(x)} \le c_2 \quad \forall x \in  \mathbb{R}^2 \right\}.
%     \]

% To show the inclusion \( N_3(P_0) \subseteq \tilde{\mathcal{P}}_{0,1,\tilde{h}} \), we first prove that:
% \[
% \forall x \in \mathbb{R}^2, \quad p_0(x) \le \tilde{h}(x).
% \]
% Since \( \tilde{p}_0(x) \le \tilde{h}(x) \) pointwise by construction, we have:
% \[
% p_0(x) = \frac{\tilde{p}_0(x)}{Z} \le \frac{\tilde{h}(x)}{Z}.
% \]

% Using the upper bound on the integral \( Z = \int \tilde{p}_0(x) dx \le 3 \), we get:
% \[
% p_0(x) \le \frac{\tilde{h}(x)}{3}.
% \]

% Now take any \( P \in N_3(P_0) \). By definition:
% \[
% \frac{1}{3} \le \frac{p(x)}{p_0(x)} \le 3 \quad \forall x.
% \]

% Substituting the upper bound of \( p_0(x) \) by \( \tilde{h}(x)/3 \), we get:
% \[
% p(x) \le 3 p_0(x) \le 3 \cdot \frac{\tilde{h}(x)}{3} = \tilde{h}(x).
% \]

% Hence, for all \( x \), \( p(x) \le \tilde{h}(x) \), which implies:
% \[
% \frac{p(x)}{\tilde{h}(x)} \le 1.
% \]
% Moreover, since \( p(x) \ge \frac{1}{3} p_0(x) \ge 0 \), we also have \( \frac{p(x)}{\tilde{h}(x)} \ge 0 \).

% Thus, \( P \ll \tilde{h} \), and \( 0 \le \frac{dP}{d\tilde{h}} \le 1 \) \( \tilde{h} \)-almost everywhere, implying \( P \in \tilde{\mathcal{P}}_{0,1,\tilde{h}} \).









% \begin{proof}
%     The proof consists of three steps. In the first 
% \end{proof}
% \section{Technical Appendices and Supplementary Material}
% Technical appendices with additional results, figures, graphs and proofs may be submitted with the paper submission before the full submission deadline (see above), or as a separate PDF in the ZIP file below before the supplementary material deadline. There is no page limit for the technical appendices.

%%%%%%%%%%%%%%%%%%%%%%%%%%%%%%%%%%%%%%%%%%%%%%%%%%%%%%%%%%%%
\section{Instructions for reproducing results}\label{appendix: reproduce}

In this appendix, we provide instructions for reproducing the experiments and figures in the paper. For a detailed description of the code, please refer to the provided file \texttt{README.md}. For tasks that require substantial runtime, we specify the running times. Tasks that complete in less than 5 seconds are excluded from such reporting.

All experiments were conducted on a system running Ubuntu 22.04.4 LTS, equipped with an Intel(R) Xeon(R) CPU @ 2.20GHz and 16GB of RAM.

\subsection{Instructions for reproducing Figure \ref{fig: Laplace pure and functional local vs global}}

From the repository root, run 
\texttt{python -m experiments.exp\_LapMixture\_visual} 
to generate the output distributions, and then 
\texttt{python -m plotting.plot\_LapMixture\_visual} 
to produce Figure~\ref{fig: Laplace pure and functional local vs global}.



The measured running time in our environment for applying the minimax-optimal sampler to the original input distribution is approximately 300 seconds.

\subsection{Instructions for reproducing results for finite space}

The finite space results are organized into six figures, grouped into two categories. 
Figures \ref{fig: finite worst-case k = 20}, \ref{fig: finite worst-case k = 10}, and \ref{fig: finite worst-case k = 100} 
present the worst-case divergences under pure LDP for different values of \( k \in \{10, 20, 100\} \). 
In contrast, Figures~\ref{fig: finite worst-case GDP k = 10}, \ref{fig: finite worst-case GDP k = 20}, 
and \ref{fig: finite worst-case GDP k = 100} show the corresponding results under \( \nu \)-GLDP for the same values of \( k \).

To generate Figures~\ref{fig: finite worst-case k = 20}, \ref{fig: finite worst-case k = 10}, and \ref{fig: finite worst-case k = 100}, 
use the script \texttt{plotting/plot\_finite\_pure.py} with the \texttt{--k} argument to specify the desired value of \( k \). 
For example, the following commands can be used to generate the respective plots:
\begin{verbatim}
python -m plotting.plot_finite_pure --k 20
python -m plotting.plot_finite_pure --k 10
python -m plotting.plot_finite_pure --k 100
\end{verbatim}

To produce Figures~\ref{fig: finite worst-case GDP k = 10}, \ref{fig: finite worst-case GDP k = 20}, 
and \ref{fig: finite worst-case GDP k = 100}, run the script \texttt{plotting/plot\_finite\_GLDP.py} with the corresponding \( k \) values:
\begin{verbatim}
python -m plotting.plot_finite_GLDP --k 10
python -m plotting.plot_finite_GLDP --k 20
python -m plotting.plot_finite_GLDP --k 100
\end{verbatim}



\subsection{Instructions for reproducing results for continuous space}

To reproduce Figure~\ref{fig: continuous worst-case}, first run the script \texttt{experiments/exp\_1DLaplaceMix\_pure.py} with different values of the privacy parameter \( \eps \). The following commands correspond to \( \eps \in \{0.1, 0.5, 1.0, 2.0\} \), respectively:

\begin{verbatim}
python -m experiments.exp_1DLaplaceMix_pure --eps 0.1 --scale 1 --seed 1 
python -m experiments.exp_1DLaplaceMix_pure --eps 0.5 --scale 1 --seed 2 
python -m experiments.exp_1DLaplaceMix_pure --eps 1.0 --scale 1 --seed 3 
python -m experiments.exp_1DLaplaceMix_pure --eps 2.0 --scale 1 --seed 4 
\end{verbatim}

These scripts can be executed independently and in any order, or run in parallel. In our environment, running all four in parallel required approximately 600 seconds. After completion, run \texttt{python -m plotting.plot\_1DLaplaceMix\_pure} to generate the final plots.

To reproduce Figure~\ref{fig: continuous worst-case GDP}, execute the script \texttt{experiments/exp\_1DLaplaceMix\_GLDP.py} with various values of the privacy parameter \( \nu \). The following commands correspond to \( \nu = 0.1, 0.5, 1.0, \) and \( 2.0 \), respectively:

\begin{verbatim}
python -m experiments.exp_1DLaplaceMix_GLDP --nu 0.1 --scale 1 --seed 1 
python -m experiments.exp_1DLaplaceMix_GLDP --nu 0.5 --scale 1 --seed 2
python -m experiments.exp_1DLaplaceMix_GLDP --nu 1.0 --scale 1 --seed 3 
python -m experiments.exp_1DLaplaceMix_GLDP --nu 2.0 --scale 1 --seed 4 
\end{verbatim}

These scripts can be executed in any order or in parallel. In our environment, running all four in parallel required approximately 80 seconds. After completion, run \texttt{python -m plotting.plot\_1DLaplaceMix\_GLDP} to generate the final figure.

To reproduce Figure~\ref{fig: continuous worst-case 2-D}, run the script \texttt{experiments/exp\_nDLaplaceMix\_pure.py} with different values of the privacy parameter \( \eps \). The following commands correspond to \( \eps \in \{0.1, 0.5, 1.0, 2.0\} \):

\begin{verbatim}
python -m experiments.exp_nDLaplaceMix_pure --eps 0.1 --seed 1 --dim 2
python -m experiments.exp_nDLaplaceMix_pure --eps 0.5 --seed 2 --dim 2
python -m experiments.exp_nDLaplaceMix_pure --eps 1.0 --seed 3 --dim 2
python -m experiments.exp_nDLaplaceMix_pure --eps 2.0 --seed 4 --dim 2
\end{verbatim}

These can be run independently or in parallel. In our environment, running all four in parallel required approximately 40 seconds. Once the first step is completed, run \texttt{python -m plotting.plot\_nDLaplaceMix\_pure ----dim 2} to generate the final figure.

% \section{Broader impact}\label{appendix: broader impact}
% Our proposed optimal samplers are designed under LDP guarantees and can be applied to privacy protection in generative modeling—a rapidly growing area of interest. A key obstacle to the adoption of privacy-preserving algorithms in practice is the potential degradation in model performance. By characterizing minimax-optimal samplers under two variants of LDP (functional and pure) and across two formulations (global and local), we develop samplers that minimize utility loss under a given privacy budget. As a result, our samplers help overcome this challenge and support broader real-world applicability.

% Although an LDP sampler provides meaningful privacy guarantees, it cannot achieve perfect privacy without completely sacrificing utility—there is an inherent trade-off between the two. Moreover, Our samplers are designed for the single-sample setting, where each client releases only one privatized data point. Extending these guarantees to scenarios involving multiple, potentially correlated releases is an important direction for future work (see Section~\ref{sec: conclusion}). In practice, clients often contribute data through multiple channels, and aggregating these disclosures can lead to greater privacy leakage.


\fi

\end{document}
